\documentclass[11pt,a4paper]{article}
\usepackage[a4paper,top=22mm,bottom=23mm,left=23mm,right=23mm]{geometry}
\usepackage{fontspec}
\usepackage{amsmath,amssymb,amsthm,mathtools}
\usepackage{unicode-math}
\setmainfont{DejaVu Serif}
\setsansfont{DejaVu Sans}
\setmonofont{DejaVu Sans Mono}[Scale=0.85]
\setmathfont{Latin Modern Math}
\usepackage{booktabs,longtable,array,enumitem,graphicx,xcolor}
\usepackage{microtype}
\usepackage[unicode=true,colorlinks=true,linkcolor=black,citecolor=black,urlcolor=blue!45!black,bookmarksnumbered=true]{hyperref}
\hypersetup{pdfauthor={Kirtus Justus},pdfsubject={Mathematical models, observation and physical identification},pdfkeywords={relational structures,identifiability,quantum models,calibration,statistical inference}}
\newtheorem{theorem}{Theorem}[section]
\newtheorem{proposition}[theorem]{Proposition}
\newtheorem{lemma}[theorem]{Lemma}
\newtheorem{corollary}[theorem]{Corollary}
\theoremstyle{definition}
\newtheorem{definition}[theorem]{Definition}
\newtheorem{example}[theorem]{Example}
\theoremstyle{remark}
\newtheorem{remark}[theorem]{Remark}
\numberwithin{equation}{section}
\setlength{\parindent}{0pt}
\setlength{\parskip}{5pt plus 1pt}
\setlength{\emergencystretch}{2em}
\setcounter{tocdepth}{2}
\setlist{topsep=4pt,itemsep=2pt,parsep=0pt}
\newcommand{\ProgramTitle}{From Prephysical Hypotheses to Physical Models}
\newcommand{\AuthorLocation}{Port Coquitlam, British Columbia, Canada}

\hypersetup{pdftitle={From Prephysical Hypotheses to Physical Models: Mathematical Studies}}
\begin{document}
\begin{titlepage}
\vspace*{12mm}
{\huge\bfseries\raggedright\hyphenpenalty=10000 \ProgramTitle\par}
\vspace{9mm}
{\Large Mathematical Supplement\par}
\vspace{7mm}
{\large Structures, observation, statistical comparison and quantum models\par}
\vspace{14mm}
Kirtus Justus\par
\AuthorLocation\par
10 November 2026\par
\vspace{12mm}
\textbf{Abstract}\par
The studies collected here examine when mathematical descriptions can be
distinguished under specified observations. Finite relational quotients,
positive event kernels, equality-pattern laws, sampling procedures, quantum
state families and controlled stochastic models provide exact examples of
identification and nonidentification. Positive results include support
reductions, finite-order separation formulas, statistical error bounds and
conditional reconstruction results. Counterexamples show why normalization,
symmetry, marginal agreement, a variational objective or sampled dynamics alone
need not determine a physical model. Every conclusion is restricted to its
stated state space, composition rule, measurement access and calibration.
The mathematical constructions are not presented as experimental observations
or as a derivation of a unique physical ontology.\par
\vfill
Research direction and scholarly responsibility: Kirtus Justus.\par
Mathematical, computational and editorial assistance: ChatGPT (OpenAI).\par
This manuscript is prepared for scholarly assessment. Such preparation does
not establish independent peer review, experimental replication or institutional
endorsement.\par
\end{titlepage}
\setcounter{page}{2}
\section{Relational Structures and the Limits of Observation}
\label{str:chapter}

A mathematical description can be useful before its symbols have a physical interpretation. The distinction matters whenever a relation, quotient, positive kernel or symmetry is proposed as a starting point. This chapter states what these structures determine under explicit mathematical assumptions. None of the constructions below assigns a physical clock, spatial metric, detector, source law or quantum state space by definition.

\subsection{Finite structures and homomorphism equivalence}
\label{str:cores}
Fix a finite relational signature whose relation symbols have positive arity, with no constants or function symbols. Objects have finite nonempty carriers and arbitrary interpretations of those relations; arrows are ordinary homomorphisms. Write $A\preceq B$ when a homomorphism $A\to B$ exists and $A\simeq_h B$ when both directions exist. A finite core is a structure all of whose endomorphisms are automorphisms.

\begin{proposition}[Core and probe equivalence]
Every finite structure has a core, unique up to isomorphism in its homomorphism-equivalence class. Moreover, $A\simeq_h B$ if and only if, for every finite structure $X$, the Boolean statements $\operatorname{Hom}(X,A)\ne\varnothing$ and $\operatorname{Hom}(X,B)\ne\varnothing$ agree.
\end{proposition}
\begin{proof}
Choose an endomorphism with image of least cardinality and take an idempotent power on its eventual image. Its image is a retract. A nonbijective endomorphism of that image would contradict minimality. A bijective endomorphism of a finite relational structure is an automorphism: its induced permutation of each finite relation maps that relation injectively into itself and hence onto itself. If cores $C,D$ admit maps $f:C\to D$ and $g:D\to C$, both composites are automorphisms. Thus $f$ has a homomorphic inverse, for example $(gf)^{-1}g$, and is an isomorphism. Probe agreement implies the two required maps by taking $X=A,B$; the converse follows by composition. This is the standard finite-core result; see Bodirsky~\cite{bodirsky}.
\end{proof}

Products and disjoint unions give meet and join in the quotient poset. Indeed, a map $X\to A\times B$ is equivalent to a pair of maps to $A,B$, and a map $A\sqcup B\to X$ is equivalent to a pair of maps from $A,B$. Consequently
\[
 [A]\wedge[B]=[A\times B],\qquad [A]\vee[B]=[A\sqcup B].
\]
These operations are associative, commutative and idempotent at the level of classes. They do not assert that a raw product of two cores is itself a core. The supported identity is
\[
 \operatorname{core}(\operatorname{core}(A)\times\operatorname{core}(B))
 \cong\operatorname{core}(A\times B).
\]
For example, the graph product $K_2\times K_2$ is two disjoint edges and retracts onto one edge, although each factor is a core.

\begin{proposition}[Initial and terminal objects]
The empty-relation singleton $E_1$ is weakly initial in this full nonempty category, but is not initial. The category has no initial object. The full-relation singleton is terminal. The thin homomorphism quotient nevertheless has a least element $[E_1]$. If empty carriers are admitted under the stated signature convention, the empty structure is initial.
\end{proposition}
\begin{proof}
Every map $E_1\to A$ preserves its empty relations, so there are $|A|$ such arrows. Every function from any nonempty $A$ into the full two-element structure $F_2$ is a homomorphism, giving $2^{|A|}\ge2$ arrows. No $A$ can therefore satisfy initiality, which requires a unique arrow to every object~\cite{stacksinitial}. There is exactly one function into the full singleton, and it is a homomorphism. Passing to a thin category retains existence of arrows and discards their multiplicity, so $[E_1]$ becomes least. With an empty carrier allowed, its unique empty function into every object preserves every positive-arity relation.
\end{proof}

In particular, $E_1\sqcup A$ cannot be a disjoint-union unit in the full category: it has one more carrier element. Binary coproducts exist there; a nullary coproduct does not. A proof in the quotient poset must not silently substitute for this full-category universal property.

The quotient also discards information. Homomorphism counts distinguish one edge from two disjoint edges although their cores coincide. Deleting a vertex from one edge removes every edge; deleting one vertex from two disjoint edges leaves an edge. Thus arbitrary deletion does not descend to the same quotient. A chosen embedding of a core requires further data. These are precise losses of a chosen equivalence relation, not defects of every richer relational description.

If an assignment sends an idempotent quotient product to the ordinary tensor product of nonzero finite-dimensional Hilbert spaces and preserves the object isomorphism, its dimension obeys $d=d^2$, hence $d=1$. This elementary object constraint is narrower than categorical no-cloning theorems, which require additional categorical and naturality assumptions~\cite{abramsky}. Infinite-dimensional Hilbert spaces can satisfy $H\cong H\otimes H$ as objects; that fact does not construct a coherent functor preserving all arrows and compositions.

A positive finite-dimensional exception exists when zero Hilbert spaces are allowed. Regard the thin quotient as a meet-semilattice category, with arrows $a\to b$ for $a\le b$, tensor product $a\wedge b$ and unit its greatest element $\top$. Strong monoidal functors from this category to finite-dimensional complex Hilbert spaces with linear maps are classified, up to monoidal natural isomorphism, by upward-closed, meet-closed subsets containing $\top$. Given such a subset $\mathcal F$, assign $\mathbb C$ on $\mathcal F$ and the zero space elsewhere; arrows within $\mathcal F$ are the identity and every other existing arrow is the unique zero map. Upward closure prevents an arrow from a supported object to an unsupported one, and meet closure together with upward closure gives
$a\wedge b\in\mathcal F$ exactly when both $a,b\in\mathcal F$. The canonical tensor identifications therefore make a strong monoidal functor.

Conversely, object idempotence forces dimension zero or one. The image of the arrow $a\to\top$ is a nonzero counit when the object is one-dimensional, by the counit identity for the diagonal. Naturality then makes every arrow between supported objects nonzero and makes support upward closed. Tensor preservation makes it meet closed. Using these counits to identify the supported spaces with $\mathbb C$ identifies all supported arrows with the identity and respects the tensor maps. This proves the classification. Requiring every object to be nonzero leaves the constant one-dimensional functor. The surviving information is membership in a chosen order filter, not a physical qubit assignment.

\subsection{Operational rank under factorizing composition}
\label{rec:operational-rank}

For an operational object \(x\), let its preparations be \(p\), its
effects be \(e\), and its evaluation probabilities be \(e(p)\).
Define \(r(x)\) as the supremum of the ranks of finite evaluation
matrices \((e_i(p_j))_{i,j}\).  When this rank is finite, a nonsingular
square matrix of size \(r(x)\) witnesses it.  The preparations and
effects in such a witness need not allow perfect single-trial
discrimination.

\begin{theorem}[A rank obstruction to operational idempotence]
Suppose \(r(x)\) and \(r(y)\) are finite.  Assume that the composite
\(x\wedge y\) admits every independent preparation \(p\otimes q\) and
product effect \(e\otimes f\) needed for rank witnesses, and that their
evaluations factorize:
\[
 (e\otimes f)(p\otimes q)=e(p)f(q).
\]
Then
\[
 r(x\wedge y)\ge r(x)r(y).
\]
If composition is idempotent as an operational object, in particular
\(r(x\wedge x)=r(x)\), every finite positive rank \(r(x)\) is one.
\end{theorem}
\begin{proof}
If either rank is zero, the inequality is immediate. Otherwise choose nonsingular witness matrices \(A\) and \(B\) of sizes \(r(x)\)
and \(r(y)\).  The admitted product preparations and effects have
evaluation matrix \(A\otimes B\).  Its inverse is
\(A^{-1}\otimes B^{-1}\), so its rank is \(r(x)r(y)\).  This proves
the lower bound.  Applying it to \(y=x\) and using operational
idempotence gives \(r(x)\ge r(x)^2\).  A finite positive integer
satisfying this inequality is one.
\end{proof}

The statement generalizes the corresponding finite-Hilbert-dimension
obstruction to an explicitly specified preparation--effect pairing
\cite[Proposition~11]{rec:operational-study}.  It needs the stated
availability and factorization premises; a relational signature alone
does not supply them.  It makes no assertion that all relational
models, infinite-rank models or nonfactorizing compositions are trivial.

\subsection{Positive event kernels and their normalization}
\label{str:kernels}
Let $D\succeq0$ be a finite Hermitian matrix with $\mathbf1^*D\mathbf1=1$. For a set of alternatives $A$ define the real nonnegative score $q_D(A)=\mathbf1_A^*D\mathbf1_A$. If a partition $(A_j)$ satisfies $\mathbf1_{A_i}^*D\mathbf1_{A_j}=0$ for $i\ne j$, its scores sum to one. A Gram factorization $D=V^*V$ then gives orthogonal grouped vectors $V\mathbf1_{A_j}$ with squared norms equal to these scores. This conditional Born-form representation does not select one partition or a physical instrument. The surrounding mathematical setting is related to quantum measure theory and composition of decoherence functionals~\cite{sorkin,boes}; it is not an independent reconstruction of laboratory quantum mechanics.

The total-event normalization is not trace normalization and does not bound every event score by one. For an explicit example take $v=(1,i)/\sqrt2$, $D=vv^*$ and three tensor copies. The event consisting of the three binary words of Hamming weight one has amplitude $3i/(2\sqrt2)$ and score $9/8$, although $\mathbf1^*D^{\otimes3}\mathbf1=1$. A probability interpretation therefore needs an admissible-event rule or additional assumptions.

A separate four-alternative example is
\[
 D_4=\frac14\begin{pmatrix}1&0&0&1\\0&1&-1&0\\0&-1&1&0\\1&0&0&1\end{pmatrix}.
\]
It has eigenvalues $1/2,1/2,0,0$ and unit total-event score. Both partitions
$\{1,2\}\mid\{3,4\}$ and $\{1,3\}\mid\{2,4\}$ are strongly decoherent with probabilities $1/2,1/2$, while their common singleton refinement is not decoherent. The permutation exchanging alternatives two and three preserves the kernel and exchanges these two partitions. A deterministic rule selecting exactly one of them cannot respect every kernel automorphism. Allowing a collection of contexts or an additional symmetry-breaking reference changes the selection problem. The event $\{1,2,4\}$ has score $5/4$, providing another direct warning against reading every subset score as a probability.

For every real Boolean vector $b$,
\[
 b^*D^{\otimes n}b=b^*\overline D^{\otimes n}b.
\]
The left side is real and conjugating it gives the right side. Thus all identical-copy Boolean observations leave a conjugation ambiguity. Joint, nonrectangular events may still reveal information absent from single-copy events. For example,
\[
 D_t=\begin{pmatrix}1/2&it\\-it&1/2\end{pmatrix},\quad |t|\le1/2,
 \qquad b=(0,1,1,0)^T
\]
gives $b^*(D_t\otimes D_u)b=1/2+2tu$. The sign relative to a specified second preparation is observable in this enlarged menu; a common reversal of both signs is not. Access to an orientation reference is an additional premise.

More generally, complete one- and two-copy Boolean scores determine a finite Hermitian $D=R+iA$ up to conjugation. Singletons and two-element events recover the real symmetric matrix $R$. For distinct $i,j$ and $k,l$, the two-copy event $\{(i,k),(j,l)\}$ has off-diagonal contribution
$2\operatorname{Re}(D_{ij}D_{kl})=2(R_{ij}R_{kl}-A_{ij}A_{kl})$.
Subtracting its known diagonal and real terms recovers every product $A_{ij}A_{kl}$. The resulting rank-one outer product determines the vector of imaginary entries up to simultaneous sign, including the zero-vector case. Nonproduct joint events are essential here. This is tomography of the stipulated scores; it does not warrant a laboratory effect for an event with score above one.

\begin{proposition}[Bounded affine observations under a singular normalizer]
Let $W\succeq0$, $W\ne0$, and consider the entire cone section
$\mathcal K_W=\{D\succeq0:\operatorname{Tr}(WD)=1\}$. Every affine functional taking values in $[0,1]$ on all of $\mathcal K_W$ depends only on the compression of $D$ to $\operatorname{supp}W$. Under the transformation $\rho=W^{1/2}DW^{1/2}$, these functionals are ordinary effects on a rank-$\operatorname{rank}W$ density-matrix space. If $W$ has rank one, all such functionals are constant.
\end{proposition}
\begin{proof}
Absorb the affine constant into the coefficient of $W$, so the functional is $\operatorname{Tr}(FD)$ on the section. For $x$ in the support with $x^*Wx=1$ and $z$ in the kernel, test the rank-one matrices $(x+tz)(x+tz)^*$ for arbitrary complex $t$. Boundedness forces the quadratic kernel block and both linear cross terms of $F$ to vanish. The restriction of $W$ to its support is invertible. The allowed compressed matrices map onto all density matrices there, and boundedness on them says $0\preceq W^{-1/2}FW^{-1/2}\preceq I$. In rank one there is only one density matrix.
\end{proof}

The associated dynamical statement also has explicit premises. If a completely positive map $\Phi(D)=\sum_\ell A_\ell D A_\ell^*$ preserves the normalizer, $\Phi^*(W)=W$, then it descends to a CPTP map on the support quotient. For $z\in\ker W$,
$\sum_\ell\|W^{1/2}A_\ell z\|^2=z^*Wz=0$, so every $W^{1/2}A_\ell$ kills that kernel. On $R=\operatorname{supp}W$, put
$K_\ell=W_R^{1/2}P_RA_\ell P_RW_R^{-1/2}$.
Then $\sum_\ell K_\ell^*K_\ell=I_R$ and the output compression is
$\sum_\ell K_\ell\rho K_\ell^*$. This conclusion assumes the completely positive framework; it is not a physical derivation of quantum channels.

For the rank-one total-event normalizer, a positive balanced addition $P\succeq0$, $P\mathbf1=0$, remains balanced after any positive linear normalization-preserving map. Indeed, its positive image has zero total-event weight and hence annihilates $\mathbf1$. If a fixed cell $A$ belongs to an exactly strongly decoherent output partition for each of two compared states $B+gP$, its probability equals
$\mathbf1_A^*\Phi(B+gP)\mathbf1=\mathbf1_A^*\Phi(B)\mathbf1$,
independent of $g$. Which partitions are admissible may depend on $g$; this does not exclude every possible observation of the added sector.

This conclusion uses the unrestricted cone and global affine bounds. A compact restricted state set, a full-rank normalizer, or a different observation rule changes the question. For instance, entrywise nonnegative PSD matrices with total entry sum one form a compact alternative closed under ordinary tensor products and nonnegative column-stochastic grouping (including deterministic partition grouping). Such an alternative supplies a different state/effect model; it is not selected by singular normalization alone.

For classical probabilities, a column-stochastic matrix $T$ sends $p$ to $Tp$. Coherent congruence $D\mapsto TDT^T$ is a different operation, although it preserves total-event normalization. If both columns of $T$ equal $(1/2,1/2)^T$ and $D=\operatorname{diag}(1/2,1/2)$, congruence gives the all-ones matrix divided by four, whereas classical postprocessing gives $\operatorname{diag}(1/2,1/2)$. The former has off-diagonal terms and trace $1/2$. A classical quantum-channel realization instead has Kraus operators $\sqrt{T_{ij}}\,|i\rangle\langle j|$ and maps diagonal states to $\operatorname{diag}(Tp)$.

The normalization $D\mapsto D/\operatorname{Tr}D$ preserves tensor products but changes convex preparation weights: a mixture $pD+(1-p)E$ maps to a mixture of the normalized states with first weight
$p\operatorname{Tr}D/[p\operatorname{Tr}D+(1-p)\operatorname{Tr}E]$.
Replacing $J=\mathbf1\mathbf1^*$ by $J+\delta I$, $\delta>0$, gives a full-rank normalizer and informative effects on every direction. The limit $\delta\downarrow0$ need not be uniform over states with balanced strength growing like $1/\delta$. These are mathematically legitimate additional choices; neither the preparation rule nor the regularization is fixed by a signed sum.

\begin{proposition}[Deterministic and conditional aggregation]
Let
\[a=(1,1,0,0)^T/\sqrt2,\qquad b=(0,0,1,1)^T/\sqrt2,\]
and define
\[
 D_\pm=\frac{I_4}{4}\pm\frac{aa^*-bb^*}{8},\qquad
 M=\begin{pmatrix}1&1&0&0\\0&0&1&1\end{pmatrix}.
\]
No deterministic CPTP map takes both $D_\pm$ to $MD_\pm M^*$. The smallest possible maximum trace-distance error for this specified pair is $1/8$. A trace-nonincreasing operation with Kraus matrix $M/\sqrt2$ attains the desired conditional outputs with success probability $1/2$ on each input.
\end{proposition}
\begin{proof}
Both input and target pairs are density matrices. Their trace distances are $1/4$ and $1/2$, respectively. Trace-distance contraction and the triangle inequality give $1/2\le1/4+2\varepsilon$, so $\varepsilon\ge1/8$. The measure-and-group channel
$D\mapsto\operatorname{diag}(D_{11}+D_{22},D_{33}+D_{44})$ has error exactly $1/8$ on both targets. Finally, $(M/\sqrt2)^*(M/\sqrt2)\preceq I$ and its output is half the required target on either input. Failure probability belongs to the complete operation and cannot be discarded in an information-per-attempt calculation.
\end{proof}
The linear map on Hermitian matrices $D\mapsto MDM^*$ is onto the four-real-dimensional Hermitian $2\times2$ space, so its real kernel has dimension $16-4=12$. This linear statement alone is not an implementation claim.

\subsection{Marginal agreement and extension}
\label{str:marginals}
For finite alphabets of sizes $q_1,\ldots,q_n$, decompose each coordinate vector space into its constant direction and its zero-sum subspace. The intersection of the kernels of all one-coordinate summations is the tensor product of the zero-sum subspaces, of dimension $\prod_i(q_i-1)$. It is exactly the space invisible to every proper marginal. Small signed perturbations in this space preserve normalization and all proper marginals of a strictly positive law. Binary parity is the simplest example; agreement of smaller marginals does not imply joint independence.

For three ternary variables define
\[
 T(x_1,x_2,x_3)=\sum_{a=0}^2\prod_{i=1}^3(\mathbf1\{x_i=a\}-1/3),
 \qquad P_\varepsilon(x)=1/27+\varepsilon T(x).
\]
There are nine all-equal or all-different words with $T=2/9$ and eighteen mixed words with $T=-1/9$. Hence positivity holds exactly on $[-1/6,1/3]$, and every proper marginal is uniform.

\begin{proposition}[Exact four-variable extension interval]
The law $P_\varepsilon$ is the three-variable marginal of an exchangeable four-variable ternary law if and only if $-1/6\le\varepsilon\le2/15$.
\end{proposition}
\begin{proof}
Average any feasible extension over coordinate and symbol permutations. Let $a,b,c,d$ be its per-word masses for occupancy types $(4),(3,1),(2,2),(2,1,1)$. Put $A=(1+6\varepsilon)/27$ and $B=(1-3\varepsilon)/27$. The marginal constraints are
\[
 a+2b=A,\qquad3d=A,\qquad b+c+d=B.
\]
Nonnegativity implies $A\ge0$ and $B\ge A/3$. Conversely $b=0$, $a=A$, $d=A/3$, $c=B-A/3$ is a nonnegative solution whenever these inequalities hold. They are equivalent to the stated interval.
\end{proof}
At $\varepsilon=1/3$ the nearest four-extendible triple has TV distance $2/5$ if uniform pair marginals must be retained. Symmetrization reduces that minimization to the displayed one-parameter family, whose nearest admissible endpoint is $2/15$. If pair marginals may vary, the sharp distance is $1/3$: any symmetrized four-variable law obeys $Q(\text{all different})\le Q(\text{mixed})$, whereas the target gives these events probabilities $2/3,0$. The bounded signed event contrast yields the lower bound. It is attained by per-word weights $a=1/9$, $b=c=0$, $d=1/54$. These two distances answer different minimization problems.

A separate finite binary example starts with positive-count law
\[
 (p_0,p_1,p_2,p_3,p_4)=(1/162,16/81,16/27,16/81,1/162).
\]
Its pair correlation is $-5/27$. Exchangeable extension is the linear feasibility problem
\[
 p_k=\sum_{j=0}^N w_j\frac{\binom jk\binom{N-j}{4-k}}{\binom N4},
 \qquad w_j\ge0,\quad\sum_jw_j=1.
\]
It is feasible for $N=5,6$ and infeasible for $N=7$; this is a finite extendibility question, not a conclusion about a fundamental source. The complete extension proofs and sharp approximation distances appear in Section~\ref{fin:four-law} and Section~\ref{fin:distances}. No size-dependent quadratic Gibbs rule is inferred from a feasible extension.

Coherent restrictions on a cofiltered system of finite nonempty discrete sets have a nonempty inverse limit by compactness: each finite set of compatibility equations is satisfied by projecting a point in a common refinement. If there is a context $K$ refining all others with compatible unique restriction maps, that limit is canonically its allowed set $C(K)$. It need not be a singleton. A nondirected contextual cover need not have a global completion, so this compactness argument cannot be applied without its index-category assumptions.

The greatest-context conclusion does not require finite or nonempty fibers. More precisely, let contexts form a small poset with a greatest refinement $K$, sets $C(i)$, and maps $\rho_{ij}:C(j)\to C(i)$ for $i\le j$, with identity and composition laws. Projection to $K$ is a bijection
\[
 \varprojlim_i C(i)\;\cong\; C(K).
\]
Indeed, $c\in C(K)$ defines the compatible family $(\rho_{iK}(c))_i$, while compatibility forces every family to have this form. These constructions are inverses. Empty $C(K)$ gives an empty limit; a nonempty or singleton limit needs the corresponding property of $C(K)$. The greatest-context argument is separate from the compactness existence result for general cofiltered finite systems.

\subsection{Symmetry does not supply a probability or a metric}
\label{str:geometry}
There is no countably additive probability on a countably infinite transitive set invariant under every permutation: all singleton masses would be equal, giving either zero total mass or an infinite sum. Likewise, uniform probabilities on finite sets are not natural under every surjection, because unequal fibers push uniform mass to unequal weights. Restricted symmetry, equal-fiber maps and probabilities on quotient patterns are legitimate distinct constructions.

A positive-definite correlation matrix need not yield a metric under $d_{ij}=-\log R_{ij}$. For example,
\[
 R=\begin{pmatrix}1&3/5&21/100\\3/5&1&3/5\\21/100&3/5&1\end{pmatrix}
\]
has eigenvalues $1/4,79/100,49/25$, but $21/100<(3/5)^2$, so $d_{13}>d_{12}+d_{23}$. A Gram realization $R_{ij}=\langle u_i,u_j\rangle$ with unit vectors instead gives the Euclidean chord distance $\|u_i-u_j\|=\sqrt{2(1-R_{ij})}$ (a pseudometric if vectors coincide). Choosing this transform does not derive a Lorentzian metric or physical dimension.

For a prescribed path graph, its Laplacian eigenvalues involve $2-2\cos(\pi k/n)$ under a stated boundary convention. The occurrence of $\pi$ follows from that spectral problem; it supplies neither physical periodicity nor the graph's physical interpretation. Similarly, a chosen complex coordinate gives two real parameters, not a clock or a $2+1$-dimensional spacetime.

For the unoriented matching Gram matrices $G_n(d)_{\pi\sigma}=d^{\#\text{loops}(\pi,\sigma)}$, positivity in every degree forces $d$ to be a nonnegative integer. For integer $d$, contractions in $(\mathbb R^d)^{\otimes2n}$ give the Gram representation. The antisymmetrizer on $n$ strands has squared norm proportional to $d(d-1)\cdots(d-n+1)$; a positive noninteger $d$ makes this negative at $n=\lfloor d\rfloor+2$, and $G_1$ excludes negative $d$. Conversely any fixed finite collection of degrees can admit nonintegers: at sufficiently large integer $d$ the contraction vectors in those degrees are independent, and positive definiteness persists in a neighborhood by continuity. All-degree positivity restricts a parameter without selecting one integer, a physical alphabet or a probability rule.

\subsection{Sector preservation and restricted observations}
\label{str:sector}
Let $P$ be an orthogonal projection, $Q=I-P$, $\Phi$ a CPTP map and $E=\Phi^*(P)$. Define $\alpha=\|P(I-E)P\|$ and $\beta=\|QEQ\|$. Then
\[
 \|E-P\|\le F(\alpha,\beta)
 =\frac{|\alpha-\beta|+\sqrt{(\alpha-\beta)^2+4\min(\alpha,\beta)}}2.
\]
This is sharp. To see it, compress $E$ to the span of the normalized $P$ and $Q$ components of a unit vector. Its matrix is $\left(\begin{smallmatrix}1-a&z\\\bar z&c\end{smallmatrix}\right)$ with $a\le\alpha,c\le\beta$. Positivity of both $E$ and $I-E$ gives
$|z|^2\le\min\{(1-a)c,a(1-c)\}$. The norm of its difference from $\operatorname{diag}(1,0)$ is at most $F(a,c)\le F(\alpha,\beta)$; direct differentiation on $a\ge c$ and $a\le c$ verifies monotonicity. Equality is attained by the qubit effect with diagonal $1-\alpha,\beta$ and off-diagonal $\sqrt{\min\{(1-\alpha)\beta,\alpha(1-\beta)\}}$, realized by a measure-and-prepare channel. Symmetric leakage $\alpha=\beta=\epsilon$ permits discrepancy $\sqrt\epsilon$, while $\beta=0$ gives $\alpha$.

An orthonormal sequence $e_n$ converges weakly to zero, so every fixed compact effect $K$ satisfies $\langle e_n,Ke_n\rangle\to0$. A finite norm-net argument makes the convergence uniform over norm-precompact families of compact effects. This is a statement about a restricted fixed menu. Moving projections $|e_n\rangle\langle e_n|$, the full coordinate PVM, or the fixed noncompact projection onto the even coordinates can retain distinctions. Such escape is not annihilation and does not establish a physical persistence law.

Exact signed balance at each separately chosen size is also weaker than coherent restrictions. Uniformly balancing $2k$ signs does not give an exactly balanced law after deleting one factor, and arbitrary deletion of two factors can remove equal signs. A rule deleting an opposite pair can preserve balance, but it supplies a specific allowed operation. Likewise, covariance is not conservation of every observable: a depolarizing channel is covariant under every unitary while reducing every nontrivial Bloch expectation. A conservation or persistence claim requires its own dynamics and invariant quantity.

\subsection{What a physical completion must additionally specify}
\label{str:physical}
A proposal without an assigned foundational physical position, wavelength, frequency or clock may still have mathematical indices, spectral variables or nonconstant functionals. None is automatically a physical location, periodic motion or activity. Signed cancellation does not determine a norm, normalized source distribution or state/effect cone. A formal Ricker function
$\psi(x)=(1-x^2)e^{-x^2/2}$ has $\int\psi=0$ and $\int\psi^2=3\sqrt\pi/4$ on a supplied real coordinate; it does not generate that coordinate or a field equation.

Conventional physical models can be used to examine additional assumptions. With a supplied quantum Hamiltonian $p^2/(2m)+V(x)$ and the usual commutation and differentiability conditions, Ehrenfest's theorem gives $m\,d^2\langle x\rangle/dt^2=-\langle V'(x)\rangle$. If $V\in C^3$ and $|V'''|\le M$ on an interval containing the position support and its mean, Taylor's theorem gives
\[
 |\langle V'(X)\rangle-V'(\langle X\rangle)|\le\tfrac M2\operatorname{Var}(X).
\]
The linear term averages to zero. Quadratic potentials make the error zero; a cubic potential can attain the bound. The mean equation does not imply persistent localization: a free Gaussian wave packet spreads. Likewise, the canonically normalized Abelian-Higgs Lagrangian with potential $\lambda(|\phi|^2-v^2/2)^2$ gives $m_h^2=2\lambda v^2$, $m_V^2=q^2v^2$ after $\phi=(v+h)/\sqrt2$. These masses follow from the supplied field theory and normalization. Semiclassical horizon entropy or a chosen higher-dimensional metric likewise presupposes its gravitational structure and cannot fill a missing prephysical-to-observable correspondence.

Useful conditional predictions do not require proof of a unique metaphysical foundation. They do require a specified physical interpretation and a prospective observation that could test it. The results in this chapter identify several mathematical obstructions and positive alternatives; the physical selection among those alternatives remains a separate question.

\clearpage
\section{Finite Signed Graphs and Binary Extension Constraints}
\label{fin:chapter}

The models in this section are finite mathematical comparators.  The
graph, real or binary variables, variational functional, probability
measure, and any quantum representation are specified separately.
Absence of a time parameter from a functional is not a derivation of
physical activity without time.  In particular, a variational minimum
does not itself choose a sampling law or an observation instrument.

\subsection{A graph double well and its exact threshold}
\label{fin:graph-action}

For a finite undirected simple graph \(G=(V,E)\), let \(A_G,D_G\)
be its adjacency and degree matrices and \(Q_G=D_G+A_G\) its signless
Laplacian.  For real amplitudes \(a_i\), define
\[
 S(a)=\frac{\kappa}{2}\sum_{\{i,j\}\in E}(a_i+a_j)^2
       +\frac{\lambda}{4}\sum_i(a_i^2-m^2)^2,
 \qquad \kappa,\lambda,m>0 .
\]
The signless-Laplacian quadratic form is standard graph theory
\cite{fin:signless}.

\begin{theorem}[Threshold, bipartiteness, and balance]\label{fin:threshold}
Let \(q_{\min}\) be the least eigenvalue of \(Q_G\).
The origin is the unique global minimum of \(S\) if
\(\kappa q_{\min}\ge\lambda m^2\).
If \(\kappa q_{\min}<\lambda m^2\), global nonzero minimizers exist and
the origin is not a global minimum.
For connected \(G\), a zero-action configuration exists exactly when
\(G\) is bipartite.  In that case its only zero-action configurations
are the two alternating assignments of magnitude \(m\).
They have zero signed total exactly when the two bipartition classes
have equal size.
\end{theorem}
\begin{proof}
Expanding the squares gives
\[
 S(a)-S(0)=\frac12a^T(\kappa Q_G-\lambda m^2 I)a
                    +\frac{\lambda}{4}\sum_i a_i^4 .
\]
When the quadratic form is nonnegative, the quartic term is strictly
positive for every nonzero \(a\), even at equality of the threshold.
When the least quadratic coefficient is negative, substitute \(a=tv\)
for a corresponding unit eigenvector.  The negative order-\(t^2\)
term dominates the positive order-\(t^4\) term for sufficiently small
nonzero \(t\).  The functional is coercive, so it attains its global
minimum, which is then nonzero.
Every term of the original functional is nonnegative.  Zero action
requires \(|a_i|=m\) and \(a_i=-a_j\) on every edge.
On a connected graph these relations propagate a unique sign choice
up to global reversal and are consistent precisely when there is no
odd cycle, equivalently when the graph is bipartite.
For bipartition \(U\sqcup W\), the signed sums are
\(\pm m(|U|-|W|)\).
\end{proof}

On the four-cycle the minima are \(\pm(m,-m,m,-m)\).
Their Hessian is \(\kappa Q_G+2\lambda m^2I\), with eigenvalues
\[
 2\lambda m^2,\quad
 2\kappa+2\lambda m^2\ \text{(twice)},\quad
 4\kappa+2\lambda m^2 .
\]
The signed sum is zero, the mean square is \(m^2\), and
\(\sum_{\{i,j\}\in E}(a_i-a_j)^2=16m^2\).
These are positive relational variation and a strict variational
minimum with zero action.  They do not describe evolution under a
physical equation of motion.  On \(K_{1,3}\), the zero-action
configurations instead have total signed amplitudes \(\pm2m\).
An equally weighted ensemble of them has mean zero, which is a
different statement from configurationwise balance.

If the amplitudes are additionally restricted to \(a_i=m\sigma_i\),
\(\sigma_i\in\{-1,+1\}\), the onsite term vanishes and
\[
 S(\sigma)=2\kappa m^2 f_G(\sigma),\qquad
 \min_\sigma S=2\kappa m^2(|E|-\operatorname{MaxCut}(G)),
\]
where \(f_G\) counts equal-sign edges.  This discrete restriction and
its Ising/MaxCut problem are not the unrestricted continuous
minimization problem.

\subsection{Action-count statistics and graph identification}
\label{fin:partitions}

An additional Gibbs stipulation gives
\[
 Z_G(z)=\sum_{\sigma\in\{\pm1\}^V}z^{f_G(\sigma)},\qquad
 p_G(\sigma)=z^{f_G(\sigma)}/Z_G(z),\qquad z>0.
\]
No probability law has been inferred from the extrema alone.

\begin{proposition}[Trees, cycles, and the even-subgraph formula]
\label{fin:partition-theorem}
Every \(n\)-vertex tree has \(Z_G(z)=2(1+z)^{n-1}\).
For a cycle of length \(n\),
\[
 Z_{C_n}(z)=(1+z)^n+(z-1)^n .
\]
For any finite graph,
\[
 Z_G(z)=2^{|V|}\left(\frac{1+z}{2}\right)^{|E|}
 \sum_{\substack{F\subseteq E\\\deg_F(v)\text{ even for all }v}}
       \left(\frac{z-1}{z+1}\right)^{|F|}.
\]
\end{proposition}
\begin{proof}
On a rooted tree, choose the root sign and, independently as
combinatorial coordinates, whether each edge joins equal signs.
These choices bijectively determine the entire configuration.
Summing their weights gives \(2(1+z)^{n-1}\).
On a cycle the transfer matrix
\(\left(\begin{smallmatrix}z&1\\1&z\end{smallmatrix}\right)\)
has eigenvalues \(1+z,z-1\); its \(n\)th-power trace counts all closed
sign sequences and gives the stated formula.
Finally, for each edge,
\[
 z^{\mathbf1_{\{\sigma_i=\sigma_j\}}}
 =\frac{1+z}{2}\left(1+\frac{z-1}{z+1}\sigma_i\sigma_j\right).
\]
Expand the product over edges.  Summing over each vertex sign kills
every term in which its degree in the selected edge set is odd;
all remaining terms contribute \(2^{|V|}\).
\end{proof}

Summing out a pendant vertex contributes \(1+z\), regardless of where
it was attached.  In particular, a connected unicyclic graph with
cycle length \(\ell\) has
\((1+z)^{n-\ell}Z_{C_\ell}(z)\).
Thus action-count statistics can lose graph shape.  Such
partition-function equivalence is an established graph-identification
issue \cite{fin:ising}.

\subsubsection{A stipulated quantum observation of the action count}

To make one observation map explicit, take an orthonormal configuration
register and a two-dimensional sign register, and stipulate
\[
 |\Psi_G\rangle=\sum_\sigma\sqrt{p_G(\sigma)}\,|\sigma\rangle
          |\chi_{f_G(\sigma)}\rangle,\qquad
 |\chi_f\rangle=
 \frac{e^{i\theta f}|+\rangle+e^{-i\theta f}|-\rangle}{\sqrt2}.
\]
The Hilbert spaces, phase rule, tensor product and Born measurements are
additional inputs.

\begin{proposition}[Reduced sign law and conditional mixedness]
\label{fin:graph-readout}
The reduced sign state and two-copy SWAP-test probability are
\[
 \rho_G=\frac12\begin{pmatrix}1&C_G\\\overline C_G&1\end{pmatrix},
 \qquad C_G=\frac{Z_G(ze^{2i\theta})}{Z_G(z)},\qquad
 p_{\rm SWAP+}=\frac{3+|C_G|^2}{4}.
\]
Graphs with identical \(Z_G\) have identical sign-register laws for
all common parameters and all measurements on independently prepared
copies.  The pure conditional-sign ansatz is not in general closed
under eliminating graph vertices by ordinary partial trace.
\end{proposition}
\begin{proof}
Tracing the orthonormal configuration labels gives
\(\rho_G=\sum_\sigma p_G(\sigma)|\chi_f\rangle\langle\chi_f|\).
Its off-diagonal entry is half the characteristic function
\(E e^{2i\theta f}\), proving the ratio formula.
The SWAP identity is
\(p_{\rm SWAP+}=(1+\operatorname{Tr}\rho_G^2)/2\), and
\(\operatorname{Tr}\rho_G^2=(1+|C_G|^2)/2\).
Equality of the reduced state gives equality of all its independently
prepared tensor powers.

For the failure of closure, eliminate vertex three of a triangle and
fix either fully specified equal-sign configuration of the retained
vertices.  The eliminated sign has action counts \(f=1,3\) with
relative weights \(z,z^3\).  The conditional visibility is
\[
 C_{\rm eq}=\frac{e^{2i\theta}+z^2e^{6i\theta}}{1+z^2},\qquad
 |C_{\rm eq}|^2=
 \frac{1+z^4+2z^2\cos(4\theta)}{(1+z^2)^2}.
\]
For \(z>0\) and \(e^{4i\theta}\ne1\), this is strictly less than one.
The resulting conditional sign state is mixed, whereas the original
ansatz assigns a pure sign state to every complete graph configuration.
Retuning the smaller graph's parameters cannot remove that difference
while retaining this pure conditional ansatz.
\end{proof}

At \(z=1/2,\theta=\pi/8\), the equal retained-sign event has probability
\(5/13\); its conditional state has \(|C_{\rm eq}|^2=17/25\) and
purity \(21/25\).  Its distance from the nearest pure qubit state is
\((5-\sqrt{17})/10\), since the Bloch-vector length is \(\sqrt{17}/5\).
For the complete triangle and square sign states, respectively,
\[
 |C_{C_3}|^2=145/169,\quad |C_{C_4}|^2=801/1681,\qquad
 p_{C_3}=163/169,\quad p_{C_4}=1461/1681.
\]
An enlarged family of arbitrary joint density matrices is closed under
partial trace; this is a valid positive comparison, but it replaces
the restricted pure conditional ansatz with ordinary mixed-state
quantum theory.

The reduced state already has a conventional random-phase preparation:
\[
 \sum_fP(f)e^{i\theta fZ}|+\!x\rangle\langle+\!x|e^{-i\theta fZ}.
\]
It is not a classical simulation of all quantum experiments.
Writing \(F,F'\) for independent action counts also gives
\[
 1-|C_G|^2=E[1-\cos(2\theta(F-F'))]
 \le4\theta^2\operatorname{Var}(F).
\]
Thus \(1-p_{\rm SWAP+}\le\theta^2\operatorname{Var}(F)\);
the leading small-phase term is
\(\theta^2\operatorname{Var}(F)\).
At \(\theta=0\), or as \(z\to0^+\) for a fixed finite graph, the
SWAP probability tends to one.  The nonzero finite-parameter contrast
does not persist uniformly in those limits.

\begin{example}[Graph aliases with different squared magnetization]
\label{fin:graph-aliases}
Let \(M=\sum_i\sigma_i\).
On vertices \(0,\ldots,4\), take edge sets
\[
 A_5=\{01,12,20,03,04\},\qquad
 B_5=\{01,12,20,03,34\}.
\]
Their degree sequences differ, so they are nonisomorphic, but
both have \(Z_5=2z(1+z)^2(3+z^2)\).
On vertices \(0,\ldots,5\), take
\[
 A_6=\{01,05,12,15,23,24,34,45\},\qquad
 B_6=\{01,04,05,12,15,23,25,34\}.
\]
Both have degree sequence \((3,3,3,3,2,2)\).
Their two degree-two vertices are nonadjacent in \(A_6\) and adjacent
in \(B_6\), proving nonisomorphism.  Both have
\[
 Z_6=2z^2(1+z)^2(1+z^2)(z^2-2z+5).
\]
Here an edge \(ij\) denotes the unordered pair \(\{i,j\}\).
\end{example}

For a complete finite check of the last claims, let
\(W_G(u)=\sum_{\sigma:M^2=u}z^{f_G(\sigma)}\).
Fixing the sign at vertex zero and doubling enumerates all
configurations.  The resulting polynomials for the two \(A\) graphs are
\[
\begin{array}{c|l}
u&W_{A_5}(u)\\ \hline
1&4z^3+12z^2+4z\\
9&4z^4+4z^3+2z\\
25&2z^5
\end{array}
\qquad
\begin{array}{c|l}
u&W_{A_6}(u)\\ \hline
0&2z^6+4z^4+8z^3+6z^2\\
4&8z^5+10z^4+8z^3+4z^2\\
16&4z^6+8z^5\\
36&2z^8 .
\end{array}
\]
For the five-vertex pair the only differences \(W_{A_5}-W_{B_5}\)
are \(+2z(1+z)(1-z)^2\) at \(u=9\) and its negative at \(u=1\).
For the six-vertex pair they are \(+2z^2(1+z)(1-z)^3\) at \(u=4\)
and its negative at \(u=0\).
These identities specify all the \(B\) polynomials as well and prove
the equal partition functions by summation.  Dividing by the common
normalization gives, for \(0<z\le1\),
\[
 \operatorname{TV}_5=
 \frac{(1-z)^2}{(1+z)(3+z^2)},\qquad
 \operatorname{TV}_6=
 \frac{(1-z)^3}{(1+z)(1+z^2)(z^2-2z+5)}.
\]
At \(z=1/2\), these are \(2/39\) and \(4/255\).
The differences of \(E M^2\) are respectively eight and four times
these distances.  The six-vertex distance decays cubically as \(z\to1\),
with leading term \((1-z)^3/16\).
Observing \(M^2\) uses configuration-sensitive information absent from
the reduced sign register.  These are exact distributional contrasts,
not measured effects or instrumental precision claims.

\subsubsection{Complete labelled sign observations recover adjacency}
\label{fin:labelled-sign-recovery}

The graph aliases above concern a restricted observation.  There is a
positive identification result when the entire labelled sign law
\(p_G(\sigma)=z^{f_G(\sigma)}/Z_G(z)\) is available, with known
\(0<z<1\) and a finite simple graph on a fixed labelled vertex set.
Fix distinct vertices \(i,j\) and any assignment \(\tau\) of all other
signs.  Write \(p_{ab}=p_G(\sigma_i=a,\sigma_j=b,\tau)\) for
\(a,b\in\{+,-\}\).  Strict positivity makes the ratio well defined, and
\[
 \frac{p_{++}p_{--}}{p_{+-}p_{-+}}
       =z^{\,2A_{ij}},\qquad
 A_{ij}=\frac{\log(p_{++}p_{--}/p_{+-}p_{-+})}{2\log z}.
\]
Indeed, an edge not joining \(i\) and \(j\) contributes equally to the
two numerator and two denominator exponents: edges to a fixed outside
sign have one equal and one unequal endpoint assignment in either
product.  The edge \(ij\), if present, contributes two in the numerator
and zero in the denominator.  Normalization cancels.  Thus the ratio
is \(z^2\) on an edge and one otherwise, recovering every labelled
adjacency entry.

At \(z=1\) the law is uniform for every graph, so the contrast vanishes.
The zero-temperature endpoint can have support zeros and does not obey
this positivity argument.  Unlabelled observations or the reduced
action-count register do not supply these four probabilities.  Exact
population recovery therefore preserves the graph-alias counterexamples
and makes no finite-sample or physical-access guarantee.

\subsection{Opposite-sign adjacency and an operator-dependent spectrum}
\label{fin:spectra}

\begin{proposition}[Complete-bipartite spectra]\label{fin:bipartite-spectrum}
The rule \(i\sim j\) exactly when \(\sigma_i\ne\sigma_j\) produces
an empty graph or a complete bipartite graph \(K_{p,q}\).
For \(p,q>0\), \(N=p+q\), the symmetric normalized Laplacian has spectrum
\[
 \{0,\ 1^{(N-2)},\ 2\},
\]
whereas the unnormalized Laplacian has eigenvalues
\[
 0,\quad q\ \text{with multiplicity }p-1,\quad
 p\ \text{with multiplicity }q-1,\quad N.
\]
\end{proposition}
\begin{proof}
The two sign classes determine all edges, giving \(K_{p,q}\).
Its normalized adjacency matrix has off-diagonal blocks
\(J_{p,q}/\sqrt{pq}\) and its transpose.  On vectors constant within
each part its eigenvalues are \(+1,-1\); on the within-part
zero-sum spaces it vanishes.  Subtracting from identity proves the
normalized spectrum.
For the unnormalized Laplacian, a vector supported on the first part
and summing to zero has eigenvalue \(q\), and analogously the other
part gives \(p\).  The remaining two-dimensional constant-on-parts
space has eigenvalues \(0,N\).
\end{proof}

Thus \(K_{1,3}\) and \(K_{2,2}\) are normalized-Laplacian isospectral,
while their unnormalized spectra are \((0,1,1,4)\) and \((0,2,2,4)\).
The normalized heat trace and its formal dimension diagnostic are
\[
 K_N(s)=\frac{1+(N-2)e^{-s}+e^{-2s}}N,\qquad
 d_{\rm spec}(s)=
 \frac{2s[(N-2)e^{-s}+2e^{-2s}]}
      {1+(N-2)e^{-s}+e^{-2s}}.
\]
At each fixed \(s>0\), \(d_{\rm spec}(s)\to2s\) as \(N\to\infty\).
Choosing \(s=1\) gives the number two while the ordinary graph diameter
stays at most two.  There is no scale-independent plateau in this
fixed-\(s\) limit.  A chosen graph operator and one diffusion-scale
number therefore do not establish a physical spatial dimension or
a local continuum.

\subsection{One four-sign law and its exact extensions}
\label{fin:four-law}

For four signs, put \(M_4=\sum_i\sigma_i\) and
\[
 P_4(\sigma)=\frac{2^{-M_4^2/4}}{81/8}.
\]
Let \(K_N\) count the plus signs among \(N\) coordinates.  Grouping
the sixteen configurations by \(K_4\) gives
\begin{equation}\label{fin:target}
 P_4(K_4=k)=
 \left(\frac1{162},\frac{16}{81},\frac{16}{27},
                    \frac{16}{81},\frac1{162}\right)_k .
\end{equation}
The law is exchangeable and invariant under global sign reversal.
Its signed mean is zero and \(E M_4^2=16/9\).
Since \(M_4^2=4+2\sum_{i<j}\sigma_i\sigma_j\),
every distinct pair has correlation \(-5/27\).
Finite extendibility is a conventional probability problem
\cite{fin:extendibility}; its requirements should be stated for
complete marginals, not inferred from a few moments.

\begin{theorem}[Unique five- and six-coordinate extensions]
\label{fin:rigidity}
If every four-coordinate marginal of a law on \(N=5\) or \(6\)
labelled signs equals the complete law \(P_4\), then the larger law
is unique, even without assuming exchangeability.
For \(N=5\), each configuration with \(K_5=2\) or \(3\) has probability
\(4/81\); the all-minus and all-plus configurations each have probability
\(1/162\); all others have probability zero.
For \(N=6\), each configuration with \(K_6=3\) has probability \(4/81\),
the two all-equal configurations have probability \(1/162\), and all
others have probability zero.
\end{theorem}
\begin{proof}
Let \(p_j\) be the count law of a uniformly selected four-coordinate
subset and \(q_j=P(K_N=j)\).
Conditional on \(K_N=k\), this sampled count has probability
\(\binom kj\binom{N-k}{4-j}/\binom N4\), regardless of exchangeability.
For \(N=5\), the nonnegative integer-count polynomial
\[
 Q_5(k)=k(5-k)(k-2)(k-3)
\]
equals \(8\) at \(k=1,4\) and zero otherwise.  Direct substitution in
the hypergeometric probabilities gives
\[
 E Q_5(K_5)=8(q_1+q_4)
          =10(p_1+p_3)-\frac{20}{3}p_2 .
\]
The right side is zero for \eqref{fin:target}, so the support is
\(\{0,2,3,5\}\).
A specified four-coordinate pattern with one plus and three minuses
has probability \(4/81\).  Its missing sign must be plus because
total count one is forbidden.  This fixes each individual
two-plus configuration at \(4/81\).  Global sign reversal of the
argument fixes the three-plus configurations; no symmetry assumption
on the candidate law is needed because the prescribed marginal itself
has the reversed pattern probabilities.  An all-plus four-pattern
can extend only to the all-plus five-pattern, fixing its mass at
\(1/162\); likewise for all minus.

For \(N=6\), use
\[
 Q_6(k)=k(6-k)(k-3)^2,\qquad
 E Q_6(K_6)=20(q_1+q_5)+8(q_2+q_4)
          =30(p_1+p_3)-20p_2 .
\]
It is nonnegative and its expectation is zero, forcing support
\(\{0,3,6\}\).
A four-coordinate pattern with three pluses and one minus then forces
both omitted signs to be minus.  This fixes every three-plus
configuration at \(4/81\), and the extreme patterns fix the remaining
two masses at \(1/162\).
For existence in either size, the listed probabilities sum to one.
For a specified four-pattern with \(k=0,1,2,3,4\) pluses, summing over
the omitted coordinates gives, respectively,
\((1/162,4/81,8/81,4/81,1/162)\) per pattern, exactly \(P_4\).
\end{proof}

The corresponding count masses are
\[
 \begin{array}{c|l}
 N=5&(q_0,q_2,q_3,q_5)=(1/162,40/81,40/81,1/162)\\
 N=6&(q_0,q_3,q_6)=(1/162,80/81,1/162).
 \end{array}
\]
This conclusion requires all four-coordinate marginals.
A single designated marginal always has an extension by
adjoining an independent sign, and no uniqueness assertion applies to
that different problem.
The displayed exact extensions require zero probabilities.
Everywhere finite Gibbs actions with counting measure give strictly
positive probabilities, so cannot realize these exact extensions.

\begin{corollary}[No seven-coordinate exact extension]\label{fin:no-seven}
No law on seven labelled signs has every four-coordinate marginal
equal to \(P_4\).
\end{corollary}
\begin{proof}
Every pair would have correlation \(-5/27\), hence
\[
 E M_7^2=7+42(-5/27)=-7/9,
\]
contradicting nonnegativity.
\end{proof}

\subsection{Sharp finite and infinite approximation distances}
\label{fin:distances}

\begin{theorem}[Odd-even finite-extension staircase]\label{fin:staircase}
Let \(N=2m-1\) or \(2m\), \(m\ge4\).
The minimum total variation from \(P_4\) to the four-coordinate
marginal of an exchangeable \(N\)-sign law is
\[
 d_N=\frac{16}{27}
 -\frac{3m(m-1)}{2(2m-1)(2m-3)}.
\]
In particular \(d_7=d_8=74/945\).
The corresponding minimum over infinitely exchangeable laws is
\(47/216\).
\end{theorem}
\begin{proof}
Conditional on a population total \(k\), the balance event \(B=\{K_4=2\}\)
has probability
\[
 h_N(k)=\frac{\binom k2\binom{N-k}2}{\binom N4}.
\]
For interior \(k\), the ratio \(h_N(k+1)/h_N(k)\) is
\((k+1)(N-k-2)/[(k-1)(N-k)]\); it is at least one precisely up to
the central count.  Thus its maximum is attained at the central
count or counts and equals
\[
 h_{\max}=\frac{3m(m-1)}{2(2m-1)(2m-3)}.
\]
Every mixture over counts obeys the same bound.
Since \(P_4(B)=16/27\), the balance-event deficit is a total variation
lower bound.

For even \(N=2m\), choose exactly \(m\) pluses uniformly over positions.
For odd \(N=2m-1\), mix the two central counts \(m-1,m\) equally and
uniformly within each count.  Both have the same four-count law:
\[
 q_0=q_4=\frac{(m-2)(m-3)}
                    {4(2m-1)(2m-3)},\qquad
 q_1=q_3=\frac{m(m-2)}{(2m-1)(2m-3)},\qquad q_2=h_{\max}.
\]
For \(m\ge4\), the outer probabilities are at least \(1/70>1/162\),
the next pair at least \(8/35>16/81\), and the center is below \(16/27\).
For example, the first two lower comparisons reduce respectively to
\((m-4)(9m-17)\ge0\) and \((m-4)(m+2)\ge0\).
Only the central category is deficient, so total variation equals
the central deficit.  Exchangeable laws are uniform within each count
class, so their full-pattern total variation equals their count-law
total variation.

For an infinite exchangeable law, restriction to every \(N\) yields
the same four-law, so the finite balance bounds imply
\(Q(B)\le\lim_Nh_{\max}=3/8\).
Consequently the distance is at least \(16/27-3/8=47/216\).
Iid fair signs attain it: their count law
\((1,4,6,4,1)/16\) has all four noncentral categories above the target.
\end{proof}

The next plateaus are \(d_9=d_{10}=22/189\) and
\(d_{11}=d_{12}=41/297\), increasing to \(47/216\).
For an arbitrary, possibly nonexchangeable \(N\)-sign law, uniformly
sampling four indices still obeys the same balance-event bound.
At least one four-coordinate marginal must therefore be at least
\(d_N\) from the target, by data processing and convexity.
These are distances between ideal distributions.  They are not
experimental detector tolerances or physical population limits.

\subsection{A sharp bounded-action approximation}
\label{fin:bounded-action}

Let \(R\ge1\).  Consider all strictly positive probability laws \(p\)
on five signs satisfying
\(\max_\sigma p(\sigma)/\min_\sigma p(\sigma)\le R\).
For a Gibbs law \(p\propto e^{-s(\sigma)}\) relative to counting measure,
this is exactly an action-range bound
\(\max s-\min s\le\log R\).
Let \(d_*(R)\) be the minimum, over such laws, of the largest total
variation distance between any of their five complete four-coordinate
marginals and \(P_4\).

\begin{theorem}[Exact action-range tradeoff]\label{fin:range-theorem}
The minimum is attained and equals
\[
 d_*(R)=
 \begin{cases}
 \displaystyle\frac{16}{27}-\frac{3R}{3+5R},
                       &1\le R\le78/5,\\[6pt]
 \displaystyle\frac{32}{27(5R+2)},&R\ge78/5 .
 \end{cases}
\]
\end{theorem}
\begin{proof}
Average any admissible law over coordinate permutations and global
sign reversal.  The constraints \(p(\sigma)\le R p(\tau)\) are linear
and invariant under these transformations, so the ratio bound
survives.  Convexity of total variation and symmetry of the target
show that the largest marginal error cannot increase.
Thus it suffices to consider per-configuration probabilities \(a,b,c\)
for count classes \(\{0,5\},\{1,4\},\{2,3\}\), respectively.
They obey
\[
 2a+10b+20c=1,\qquad
 q=(a+b,\ 4(b+c),\ 12c,\ 4(b+c),\ a+b).
\]

First, \(a,b\ge c/R\) implies
\(c\le R/[4(5R+3)]\).  The central event therefore gives
\[
 \operatorname{TV}(P_4,q)\ge16/27-3R/(5R+3).
\]
Set
\[
 a=b=\frac1{4(5R+3)},\qquad c=\frac{R}{4(5R+3)}.
\]
These weights are positive, normalized and have ratio \(R\).
For \(1\le R\le78/5\), \(q_0=q_4\ge1/162\) and
\(q_1=q_3>16/81\), while \(q_2<16/27\).
Hence the total variation is exactly the central deficit, attaining
the first formula.

For the second lower bound, put \(v=3R/(5R+2)\in[0,1]\) and use
the test function \(f(2)=1\), \(f(0)=f(4)=v\), \(f(1)=f(3)=0\).
After eliminating \(a\) with normalization,
\[
 E_qf=12c+2v(a+b)
     =v+\frac{24(c-Rb)}{5R+2}\le v.
\]
Since \(E_{P_4}f=16/27+v/81\), the bounded-function characterization
of total variation gives
\[
 \operatorname{TV}(P_4,q)\ge16/27+v/81-v
                          =\frac{32}{27(5R+2)}.
\]
For \(R\ge78/5\), choose
\[
 b=\frac{20}{81(5R+2)},\qquad
 c=Rb,\qquad
 a=\frac1{162}-b=\frac{5R-38}{162(5R+2)}.
\]
They are positive and normalized.  Here \(b\le a\le c\), so their
maximum-to-minimum ratio is \(R\).
The outer marginal masses exactly match \(1/162\);
each near-central mass exceeds its target by
\(16/[27(5R+2)]\), while the central deficit is
\(32/[27(5R+2)]\).  The lower bound is attained.
At \(R=78/5\) the constructions and formulas agree.
\end{proof}

Thus the error is positive for every finite action range and decays
as \((32/135)R^{-1}+O(R^{-2})\) for large \(R\).
This concerns all four-subset marginals simultaneously.  An exact
extension of only one designated marginal can instead be formed by
adjoining an independent sign, with finite positive weights.
The probability-ratio optimization does not postulate a physical
action law or calibrate an instrument.

\subsection{Quadratic weights are not closed under deletion}
\label{fin:quadratic}

\begin{proposition}[A five-to-four closure obstruction]
\label{fin:quadratic-closure}
For \(b>0\), the size-dependent quadratic family
\[
 P_N(\sigma)\ \propto\ \exp[-b M_N^2/(2N)]
\]
is not projectively consistent under ordinary coordinate deletion.
Even allowing an adjusted quadratic coefficient at size four does not
make the five-to-four marginal a pure quadratic weight.
\end{proposition}
\begin{proof}
For retained magnetization \(M\in\{0,\pm2,\pm4\}\), summing the fifth
sign gives
\[
 w_b(M)=2e^{-b(M^2+1)/10}\cosh(bM/5).
\]
Every weight \(w(M)=C e^{-\alpha M^2}\) necessarily satisfies
\(w(4)w(0)^3/w(2)^4=1\), because both its constant and quadratic
exponents cancel.  Here the ratio is
\[
 \frac{\cosh(4b/5)}{\cosh^4(2b/5)}
 =1-\tanh^4(2b/5)<1.
\]
To obtain the identity, set \(x=2b/5\), use
\(\cosh(2x)=2\cosh^2x-1\), and divide by \(\cosh^4x\).
\end{proof}

The multiplicative defect is \(16b^4/625+O(b^6)\), so a lower-order
expansion can miss it.  At \(b=0\) the iid fair family is exactly
projective.  There is also a small-size exception: after reducing a
three-sign law to two signs, only two retained values of \(M^2\) remain, so a
normalization and a retuned quadratic coefficient can fit both
positive weights.  Three distinct \(M^2\) values in the five-to-four
experiment impose the extra compatibility equation that fails.
The valid bare probability extensions of Theorem~\ref{fin:rigidity}
are therefore not generated by simply reusing the positive quadratic
action.  Their existence and this particular action's failure are
compatible facts.

\clearpage
\section{Resources, Equality Patterns and Composition}
\label{pat:chapter}

This chapter uses classical probability models to distinguish invariance, consistency, observable information and controllable communication. The sample positions are mathematical labels. Counts of positions or repetitions are not assigned physical time, spatial volume or energy without a separate resource model. Logarithms in relative entropy are natural; reported information in bits divides by $\ln2$.

\subsection{One common reference and a complete preparation}
\label{pat:reference}
Let a finite or countable message $\Theta$ select a preparation law $P_\theta$ and let $\mu$ be one common normalized reference. If $D(P_\theta\|\mu)\le K<\infty$ for every admitted message, then any fixed stochastic observation of the complete preparation obeys
\[
 I(\Theta;Y)\le I(\Theta;X)\le K,\qquad
 I(\Theta;Y)_{\rm bits}\le K/\ln2.
\]
Indeed, with $\overline P=\sum_\theta\pi_\theta P_\theta$, the relative-entropy identity
\[
 \sum_\theta\pi_\theta D(P_\theta\|\mu)
 =I(\Theta;X)+D(\overline P\|\mu)
\]
and data processing prove the bound. The reference and the complete-preparation budget are premises. Choosing a separate reference for each message would invalidate this argument. Entropy and communication are used in their ordinary Shannon sense~\cite{shannon}.

Marginal budgets need not bound joint information. Let $\Theta$ and $X$ be independent fair bits and set $Y=X\oplus\Theta$. Both marginal laws are independent of the message, while $(X,Y)$ determines it. Likewise, a message encoded in distinct Bell states has identical single-qubit marginals. The message-conditioned correlation resource belongs to the complete preparation. It cannot be replaced by the correlations of the message-averaged state. Preshared entanglement and selection on successful attempts also change the resource accounting.

For $n$ individually addressable binary positions, a hard Hamming-weight budget of one permits $n+1$ perfectly distinguishable strings and gives $\log_2(n+1)$ bits. A mean Hamming-weight budget of one instead gives $n h_2(1/n)$ for $n\ge2$, attained by independent Bernoulli$(1/n)$ positions. The upper bound follows from entropy subadditivity and concavity of binary entropy. Both quantities grow with $n$; a scalar cost bound alone does not bound the number of distinguishable modes.

\begin{theorem}[Finite entropy under a mean cost]
\label{pat:entropycriterion}
Let $J$ be a finite or countable alphabet with finite costs $r_j\ge0$ and at least one zero-cost symbol. Admit all probability laws on $J$ and ideal identification of its symbols. For $E>0$, the supremum of $H(p)$ over $\sum_jp_jr_j\le E$ is finite if and only if
\[
 Z(\beta)=\sum_j e^{-\beta r_j}<\infty
\]
for some $\beta>0$. Finiteness for one positive $E$ implies finiteness for every finite positive $E$.
\end{theorem}
\begin{proof}
A finite $Z(\beta)$ gives the Gibbs law $q_j=e^{-\beta r_j}/Z$ and $H(p)\le\beta E+\log Z$. Conversely, an infinite finite-cost sublevel allows a fixed positive mass spread over arbitrarily many symbols, with the remaining mass on a zero-cost symbol; entropy is unbounded. Suppose all sublevels are finite and write $N(R)=\#\{j:r_j\le R\}$. If $\log N(R)/R$ were bounded for large $R$, summing integer-width cost shells would make $Z(\beta)$ finite for a sufficiently large $\beta$. If no such $\beta$ exists, take $R\to\infty$ with $\log N(R)/R\to\infty$, place mass $E/R$ uniformly on the nonreference symbols of cost at most $R$, and place the remainder on the zero-cost symbol. Its entropy is at least $(E/R)\log(N(R)-1)$ and is unbounded.
\end{proof}
For a finite alphabet, the endpoint $\beta=0$ gives the uniform law and $H(p)\le\log|J|$. On an infinite counting alphabet $Z(0)$ diverges, so that endpoint supplies no finite entropy bound.

Finiteness of one Gibbs sum is not an entropy-continuity theorem. Nor does the theorem determine physical costs. For example, $2^{\ell^2}$ alternatives at cost $\ell$ permit entropy $h_2(1/\ell)+\ell$ at mean cost one, despite every finite-cost sublevel being finite.

\subsection{A finite entropy supremum without a maximizer}
\label{pat:nonattainment}

The word \emph{supremum} in Theorem~\ref{pat:entropycriterion} cannot in general be replaced by \emph{maximum}. Take one zero-cost symbol and $g_n=\lfloor2^n/n^3\rfloor$ symbols of cost $n$ for each positive integer $n$, omitting empty levels. At $\beta_0=\log2$, put
\[
 Z_0=1+\sum_{n\ge1}g_n2^{-n},\qquad
 q_j=\frac{e^{-\beta_0r_j}}{Z_0},\qquad
 \bar r=\sum_jq_jr_j.
\]
The sum $Z_0$ is finite because $g_n2^{-n}\le n^{-3}$. Also $g_1=2$ gives $Z_0\ge2$, so
\[
 0<\bar r\le\frac12\sum_{n\ge1}\frac1{n^2}<1;
\]
the last inequality follows from $\sum_{n\ge1}n^{-2}<1+\int_1^\infty x^{-2}\,dx=2$.

At mean-cost budget one, Gibbs comparison gives $H(p)\le\log Z_0+\beta_0$. Let $U_n$ be uniform on the level of cost $n$, and for sufficiently large $n$ set
\[
 p_n=(1-\varepsilon_n)q+\varepsilon_nU_n,\qquad
 \varepsilon_n=\frac{1-\bar r}{n-\bar r}.
\]
This law has mean cost exactly one. Entropy concavity and $n^{-1}\log g_n\to\log2$ give
\[
 \liminf_n H(p_n)\ge H(q)+(1-\bar r)\beta_0
 =\log Z_0+\beta_0.
\]
Thus this upper bound is the supremum. Equality for an admissible law would require both mean cost one and $D(p\|q)=0$. The latter forces $p=q$, whose mean is strictly less than one. The supremum is not attained.

Moreover, $p_n\to q$ in total variation while the entropies approach a value strictly larger than $H(q)$. Finiteness of a Gibbs sum at one temperature therefore does not guarantee entropy continuity on the cost-bounded set. The example's Gibbs sums diverge for $0<\beta<\log2$; it does not satisfy the stronger all-positive-temperature hypothesis used in entropy-continuity results~\cite{beckerEntropy}. The integer-partition example below has its own stronger partition-product and mean-cost argument and is not contradicted by this construction. The example was supplied in the related scientific assessment and is proved here in full~\cite{justusAssessment}.

\subsection{Orientation, abstract channel uses and physical rates}
\label{pat:orientation}
The two-copy event score in Section~\ref{str:kernels} gives the hypothetical binary law $p(t,u)=1/2+2tu$ when an oriented reference $u$ is fixed. Choosing $u=1/2$ and inputs $t=\pm1/8$ gives outcome probabilities $5/8$ and $3/8$. As an ordinary binary symmetric communication channel, with that effect assumed available, its capacity is
\[
 C=1-h_2(3/8)=0.045565997075\ldots\quad\text{bits per use}.
\]
An unlimited number of independently accessible uses has unbounded total information at this positive rate. The copy resource and the observation effect cannot be omitted from a proposed total information limit.

Using two identical preparations instead gives $p(t)=1/2+2t^2$. Its Bernoulli Fisher information about $t$ and relative entropy from $t=0$ satisfy
\[
 \mathcal J(t)=\frac{64t^2}{1-16t^4},\qquad
 D(\operatorname{Bern}(1/2+2t^2)\|
 \operatorname{Bern}(1/2))=8t^4+O(t^8).
\]
The first formula is $(p')^2/[p(1-p)]$; the second follows by expanding the two logarithms at zero. Sign information is absent in this identical-copy law. With a fixed nonzero reference $u$, the corresponding Fisher information is
$16u^2/(1-16t^2u^2)$ and is nonzero at $t=0$. These are properties of the specified observation laws, not a derivation of an oriented apparatus.

For comparison, the conventional Gaussian channel formula
$C(B)=B\log_2(1+P/(N_0B))$ has limit $P/(N_0\ln2)$ as $B\to\infty$ at fixed signal power $P$, with $N_0$ defined so the noise power is $N_0B$~\cite{shannon}. The limit is the elementary limit $\log(1+x)/x\to1$. Bandwidth in hertz, power and information per second already presume physical time and a specified channel. No observer-independent rate or whole-system capacity follows from an uncalibrated mathematical order parameter.

A finite-dimensional channel kernel and a mathematical cost can define a capacity in bits per block without physical time or energy. The Gaussian vector example and its countercontrols are developed in \emph{Abstract Channels and Gaussian Capacity}. Deriving that kernel from a proposed foundation and interpreting a physical rate remain separate questions.

\subsection{Equality patterns, raw probabilities and deletion}
\label{pat:orbits}
For $k$ positions carrying labels in a countably infinite set, two tuples are related by a common permutation of the label set exactly when they have the same equality pattern. Thus the orbits are the $B_k$ set partitions of the $k$ positions. If position labels are also forgotten, only the integer partition of block sizes remains; there are $p(k)$ such types. Neither quotient gives an invariant normalized raw-label distribution. An invariant marginal on a countably infinite transitive label set would contradict countable additivity.

Uniformity at each size does not give a consistent family. Under uniform set partitions, two positions agree with probability $1/2$, whereas restricting the five equally likely three-position partitions gives $2/5$. Under uniform integer-partition types at size three, random deletion gives equality probability $(1+1/3+0)/3=4/9$, differing from the size-two uniform law by TV $1/18$.

A positive alternative is a deterministic paintbox: choose probabilities $p_i$, draw labels independently, and observe their equality partition. A different alternative is an Ewens law with parameter $\theta>0$, whose probability for a particular partition with block sizes $n_1,\ldots,n_b$ is
\[
 \frac{\theta^b\prod_{j=1}^b(n_j-1)!}{\theta(\theta+1)\cdots(\theta+k-1)}.
\]
These are consistent exchangeable laws, not unique selections by symmetry. Paintboxes and exchangeable partitions are treated systematically by Pitman~\cite{pitman}.

Separate pair probabilities do not determine joint pair observations. Independent fair binary labels and the Ewens$(1)$ law both give pair-equality probability $1/2$. For three positions their all-equal probabilities are $1/4$ and $1/3$, and their complete partition laws have TV $1/4$. Jointly recording two pair indicators already distinguishes them. The observation menu must specify whether outcomes are aligned on the same sample.

\subsection{Arbitrarily late first separation}
\label{pat:prouhet}
\begin{theorem}[Paintboxes with a prescribed finite cutoff]
For every integer $K\ge2$ there are distinct, full-support, fixed paintboxes whose complete partition laws agree through order $K$. At $m=K+1$, if $\Delta_m$ is the difference of their $m$th collision moments, the exact complete-pattern TV is
\[
 \operatorname{TV}_m=\frac{|\Delta_m|}{2}
 \sum_{b=1}^m {m\brace b}(b-1)!,\qquad
 |\Delta_m|=\frac{m!}{2^{m(m-1)/2}(2^m+1)^m}.
\]
\end{theorem}
\begin{proof}
Split $1,\ldots,2^m$ according to the parity of the binary digit sum of the predecessor. The polynomial identity
$\prod_{j=0}^{m-1}(1-e^{2^jt})=\sum_{a=0}^{2^m-1}(-1)^{s_2(a)}e^{at}$
has a zero of order $m$ at the origin. Its derivatives show equality of power sums through $m-1$ and difference $(-1)^m m!2^{m(m-1)/2}$ at order $m$; shifting $a$ to $a+1$ leaves that first difference unchanged. Both groups sum to $S=2^{m-2}(2^m+1)$. Give a core entry $a$ the probability $a/(2S)$ and append common tail probabilities $2^{-j-1}$, $j\ge1$, of total mass $1/2$. The tails cancel in moment differences, giving the formula for $\Delta_m$.

For a specified partition with $b$ blocks, its probability is a sum over distinct block labels. Inclusion-exclusion on identifications of those labels expresses it as a polynomial in collision moments. At total order $m$, the only term involving the differing moment $s_m$ merges all $b$ blocks and has coefficient $(-1)^{b-1}(b-1)!$. Every other term agrees. Summing absolute differences over the ${m\brace b}$ partitions with $b$ blocks proves the formula. Strict positivity of every finite partition event follows from the infinite positive tail.
\end{proof}
The first values, for $K=2,3,4,5$, respectively, are
\[\frac1{324},\quad\frac{39}{668168},\quad\frac{125}{556592256},\quad\frac{4869}{30891577600000}.\]
They are exact mathematical distinctions, however small.

Any sequence of independently restarted experiments restricted to the agreeing orders has identical history laws under these two models, including a common adaptive choice of the next allowed order. This follows inductively from equality of each conditional next-record law. Shared references across repetitions or growing joint arity define a different experiment.

At $K=2$ the core weights are $(1,4,6,7)/36$ and $(2,3,5,8)/36$, with the common geometric tail. The five three-position patterns have probabilities
\[
\begin{array}{c|cc}
 & A&B\\\hline
\text{all equal}&425/13608&439/13608\\
\text{each specified pair only}&445/3402&883/6804\\
\text{all distinct}&7843/13608&7871/13608
\end{array}
\]
(the middle row represents three patterns). For $n$ independent complete triple records, a test of size at most $0.05$ and power at least $0.95$ must satisfy
$nD(P_{\rm alt}\|P_{\rm null})\ge0.9\log19$, by data processing to the binary decision. Evaluation gives necessary counts $81{,}288$ with $A$ as null and $81{,}670$ with $B$ as null. These are not sufficient counts and do not apply to one larger joint equality partition.

\subsubsection{A separate finite-order ambiguity for exchangeable signs}
\label{pat:latent-sign-aliases}

A related construction concerns signed records, with a different
state space and observation rule from equality-pattern paintboxes.
Let a latent \(T\in[-1,1]\) have law \(\nu\), and conditional on \(T\)
draw signs independently with mean \(T\).  The resulting complete law is
\[
 P_\nu^{(m)}(s_1,\ldots,s_m)
 =2^{-m}\int_{-1}^1\prod_{i=1}^m(1+s_it)\,d\nu(t).
\]
These laws are exchangeable and consistent under every coordinate
deletion, because they are restrictions of one conditional iid process.

For \(0<\eta<1\), compare the fair prior on
\(\{\sqrt\eta,-\sqrt\eta\}\) with the prior assigning masses
\(1-\eta,\eta/2,\eta/2\) to \(0,1,-1\).  Their first three moments
agree, while their fourth moments are \(\eta^2\) and \(\eta\).
Expanding the displayed product proves equality of every joint law
through three signs and
\[
 P_B^{(4)}(s)-P_A^{(4)}(s)
 =\frac{\eta(1-\eta)}{16}\prod_{i=1}^4s_i,\qquad
 \operatorname{TV}(P_A^{(4)},P_B^{(4)})
 =\frac{\eta(1-\eta)}2.
\]

For an arbitrary finite cutoff \(k\), take an even \(n>k\), let \(L_n\)
be the Legendre polynomial and choose \(0<\varepsilon\le1\).  Define
\[
 d\nu_\pm(t)=\tfrac12[1\pm\varepsilon L_n(t)]\,dt,
 \qquad -1\le t\le1.
\]
These are symmetric probability measures.  To verify the needed
polynomial facts, Rodrigues' formula
\[
 L_n(t)=\frac1{2^n n!}\frac{d^n}{dt^n}(t^2-1)^n
\]
and integration by parts give orthogonality to every polynomial of
degree below \(n\).  The identity
\[
 L_n(t)=\frac1\pi\int_0^\pi
       \bigl(t+i\sqrt{1-t^2}\cos u\bigr)^n\,du
\]
(obtained by binomial expansion) gives \(|L_n(t)|\le1\).
Thus the two densities are nonnegative and normalized.  Their
moments below order \(n\) agree, so all smaller complete sign laws agree.
At order \(n\), only the top monomial in the product can contribute
to their difference.  Another integration by parts yields
\[
 \int_{-1}^1t^nL_n(t)\,dt
 =2^{-n}\int_{-1}^1(1-t^2)^n\,dt
 =\frac{2^{n+1}(n!)^2}{(2n+1)!}.
\]
Consequently
\[
 \operatorname{TV}(P_{\nu_+}^{(n)},P_{\nu_-}^{(n)})
 =\varepsilon\,\frac{2^n(n!)^2}{(2n+1)!}>0 .
\]
For \(n=4,\varepsilon=1/2\), this distance is \(4/315\).

The obstruction is to a fixed finite observation order over the entire
declared class.  Complete knowledge of all orders identifies every
moment of \(\nu\); on the compact interval, polynomial approximation
of continuous functions then identifies \(\nu\) itself.
The latent source, conditional sampling and accessible joint order
remain supplied assumptions.  Neither the sign construction nor the
paintbox construction selects a physical source.

\subsection{Entropy of quotient outputs and controllable messages}
\label{pat:coherententropy}
If every set partition at order $k$ has common cost $c_k$ and can be selected and distinguished, Theorem~\ref{pat:entropycriterion} gives
\[
 \exists\beta>0:\sum_k B_ke^{-\beta c_k}<\infty
 \quad\Longleftrightarrow\quad
 \liminf_{k\to\infty}\frac{c_k}{\log B_k}>0.
\]
For a positive lower ratio, choose $\beta$ so that the summand is at most $B_k^{-1}\le2^{-(k-1)}$. If the lower ratio is zero, infinitely many terms exceed $\sqrt{B_k}$ for each fixed $\beta$. Linear costs fail; costs of order $k\log k$ can suffice. This identifies a mathematical growth condition without explaining why a physical apparatus should obey it.

For integer-partition types with cost equal to order and the empty type allowed,
\[
 Z(\beta)=\prod_{j\ge1}(1-e^{-\beta j})^{-1},\qquad
 \sum_{j\ge1}\frac{j}{e^{\beta_*j}-1}=1.
\]
The unique mean-one entropy maximizer is the Gibbs law, and
\[
 H_*=(\beta_*+\log Z(\beta_*))/\log2
 =2.420267156348295\ldots,\quad
 \beta_*=1.074984419833593\ldots.
\]
The equations, rather than rounded decimals, define the exact optimum.

Requiring the conditional pattern laws at every order to come from one exchangeable consistent family gives a strict gap. Let $G_k$ be the Gibbs probability of order $k$, and let $c$ be the common probability of a uniformly selected pair agreeing in orders two and three. The two uniform type laws have pair probabilities $u_2=1/2,u_3=4/9$. Relative-entropy decomposition, data processing to that pair indicator, and minimization over order weights give
\[
 H_*-H(P)\ge-\frac1{\log2}\log\max_{0\le c\le1}
 \left[1-G_2-G_3+G_2e^{-d(c\|1/2)}+G_3e^{-d(c\|4/9)}\right]>0.
\]
Here $d$ is Bernoulli relative entropy. Both exponential terms are at most one, with equality at different $c$, proving strictness by compactness. The resulting relaxed upper value is approximately $2.41988092049132$ bits. A specified finite-order example uses the Ewens$(4/3)$ integer-type law
\[
 Q_k(\lambda)=\frac{k!\theta^{\sum_jm_j}}{(\theta)_{k\uparrow}}\prod_{j\ge1}\frac1{j^{m_j}m_j!},\qquad\theta=4/3,
\]
where $m_j$ counts blocks of size $j$. Put $h_k=-\sum_\lambda Q_k(\lambda)\log Q_k(\lambda)$ and choose order weights
\[
 \pi_k=\frac{\exp(h_k-1.047262\,k)}{\sum_{j=0}^{22}\exp(h_j-1.047262\,j)},\quad0\le k\le22,
 \qquad\pi_k=0\quad(k>22).
\]
The conditional type laws are restrictions of one coherent exchangeable family. Evaluation gives
\begin{align*}
\mathbb E K&=0.999996239443005554\ldots,\\
H&=2.395335833829052067\ldots\text{ bits}.
\end{align*} These numerical evaluations, including the mean's proximity to one, are not formal interval enclosures. The formulas specify the example exactly; no infinite-order lottery or exact mean-one retuning is implicit in the quoted digits. The exact coherent optimum remains undetermined.

If a sender controls only the order $K$ and the family $Q_k$ is fixed, then $M\to K\to\Lambda$ is a Markov chain and
\[
 I(M;K,\Lambda)=I(M;K),\qquad
 H(K,\Lambda)=H(K)+H(\Lambda\mid K).
\]
At $\mathbb EK\le1$, $K\in\mathbb N_0$, the largest order entropy is two bits, attained by the geometric law $P(K=k)=2^{-k-1}$. Conditional pattern randomness is not an extra controllable message alphabet. For the stated Ewens witness the decomposition is approximately $1.984978648746772$ order bits plus $0.410357185082280$ conditional pattern bits. If the empty order is disallowed, positive integer $K$ with mean at most one forces $K=1$.

\subsection{Composition and aligned observations}
\label{pat:composition}
Two quotient patterns do not determine a shared-label concatenation. Types $(2)$ and $(1)$ can arise from $(a,a)$ with $(a)$, yielding $(3)$, or with $(b)$, yielding $(2,1)$. Requiring disjoint fresh labels supplies a valid multiset-union operation, but adds a reference convention.

For independent paintboxes, common refinement is the equality law of independently drawn label pairs and has product weights. It need not stay in the original parameter family. Refining two Ewens$(1)$ laws gives all-equal probability $1/k^2$. Matching its second and third orders to a Pitman--Yor law forces $(\alpha,\theta)=(2/7,13/7)$. At order four the all-equal probabilities $1/16$ and $19/306$ differ; summing all fifteen pattern differences gives TV $59/4896$. Agreement at early orders does not establish closure at all orders.

Refinement commutes with restriction; join need not. If one input joins $1$ to $2$ and the other joins $2$ to $3$, their join connects $1$ and $3$, even after restricting the joined relation to $\{1,3\}$. Joining the restricted inputs leaves them separate. A set-valued disjoint-carrier amalgamation is associative as a set of possible completions. Uniformly selecting one completion at each binary step is a further stochastic rule and can fail associativity (already by TV $1/12$ on three singleton carriers).

Independent self-refinement leaves a random partition law fixed only if each pair-coincidence probability obeys $p_{ij}=p_{ij}^2$. On a countable carrier this forces a deterministic partition. Full exchangeability then leaves only the discrete and universal partitions. Perfectly coupling the two inputs as the same random partition removes the independence premise and permits any law.

To see an all-order coupling ambiguity, let independently sampled pairs $(X,Y)$ have probability matrix
\[
 R_t=\begin{pmatrix}t&1/2-t\\1/2-t&t\end{pmatrix},\quad0<t<1/2.
\]
Both component processes are fair and have unchanged full equality-pattern laws. Composite pair equality has probability $1/4+4(t-1/4)^2$. Comparing $t=1/4$ and $3/8$ gives TV $1/16$ for that binary observation and $1/8$ for the aligned pair of component equality indicators. On a countable alphabet the full-support coupling
$(1-\epsilon)p_ip_j+\epsilon p_i\mathbf1_{i=j}$ with $p_i=2^{-i}$ likewise preserves both components. Comparing $\epsilon=0,1/2$ changes composite equality from $1/9$ to $23/126$, a gap $1/14$.

Even the separate all-order laws of both components and their composite can all agree while aligned observations differ. Divide each of the following matrices by $40$:
\[
 A=\begin{pmatrix}1&2&3&4\\2&3&4&1\\3&4&1&2\\4&1&2&3\end{pmatrix},\quad
 B=\begin{pmatrix}1&2&3&4\\2&1&4&3\\3&4&1&2\\4&3&2&1\end{pmatrix}.
\]
Their row and column sums coincide, and their multisets of cell masses coincide. Thus all three separate equality-pattern laws agree at every order. The aligned four-position laws have TV $3/1250$. In particular, the event
$X_1=X_2\ne X_3=X_4$, $Y_1=Y_3\ne Y_2=Y_4$
has unnormalized numerator
$\sum_{i\ne k,j\ne\ell}M_{ij}M_{i\ell}M_{kj}M_{k\ell}$,
equal to $4360$ and $4488$ respectively; the probability gap is $1/20000$. This explicit finite enumeration is included among the computational checks.

After $m$ independent fair binary refinements there are $D=2^m$ equally weighted categories. The chance of any collision among $k$ labels is
\[
 1-\frac{(D)_k}{D^k}.
\]
It tends to zero for fixed $k$, but to $1-e^{-c^2/2}$ when $k/\sqrt D\to c>0$. The logarithm of the product $\prod_{j<k}(1-j/D)$ proves the latter limit. A fixed-order disappearance statement is therefore not uniform over growing observation order.

\subsection{Operations natural under every restriction}
\label{pat:naturality}
\begin{theorem}[Deterministic selection]
For a finite integer $r\in\mathbb N_0$, total operations $F_S:\operatorname{Eq}(S)^r\to\operatorname{Eq}(S)$, on every finite carrier and commuting with every injection-induced restriction, are exactly the discrete constant and intersections of a fixed subset of the $r$ input relations. The empty intersection is the universal relation. There are $2^r+1$ operations.
\end{theorem}
\begin{proof}
Pair restriction makes off-diagonal output membership a Boolean function of the $r$ input memberships. Let $\mathcal W$ be the coalitions of input coordinates that make the pair equal. On three positions, coalitions $A,A,B$ on edges $12,23,13$ are realizable when $A\subseteq B$, so transitivity makes $\mathcal W$ upward closed. Coalitions $A,B,A\cap B$ are also realizable, so $\mathcal W$ is intersection closed. If nonempty it is the principal filter generated by $\cap\mathcal W$; if empty the output is discrete. Each listed operation plainly satisfies the assumptions.
\end{proof}
At arity zero, the input tuple is unique and restriction to a pair selects exactly one of two constant rules: the discrete or universal relation. Both extend naturally to every finite carrier, so the same count $2^0+1=2$ applies.

Finite arity is essential. For countably many input relations, declare a pair equivalent when it is equivalent in all but finitely many inputs. The cofinite filter is closed under finite intersections, so this defines an equivalence relation and commutes with every restriction. It is not intersection over a fixed subset of inputs: acceptance of every coalition $\mathbb N\setminus\{j\}$ would force that subset to be empty, which would also accept the empty coalition. The cofinite rule does not. Finite carriers do not remove this infinite-arity alternative.

For two inputs, commutativity and pointwise idempotence select intersection in this deterministic class. The random kernel $\tfrac12\delta_R+\tfrac12\delta_Q$ shows that simply dropping determinism invalidates that conclusion.

\begin{theorem}[Stochastic selection and a local bound]
Let $K_S(R,Q)$ be a total kernel on finite equivalence relations, restriction-natural at every size, commutative and pointwise idempotent. Let $a$ be its two-position equality probability on mixed input statuses. Then $0\le a\le1/2$. Ordinary conditional composition on the input triple $(1,0,0)$ has associator TV $\epsilon_2=a(1-a)$. If $\kappa(R,Q)$ counts pairs of distinct blocks of $R\cap Q$ sharing a parent block in at least one input, then
\[
 \operatorname{TV}(K_S(R,Q),\delta_{R\cap Q})
 \le\min\{1,\kappa(R,Q)a\},\qquad a\le2\epsilon_2.
\]
Exact associativity therefore forces deterministic intersection.
\end{theorem}
\begin{proof}
On three positions take inputs identifying only $12$ and only $23$. Both separate $13$, so the output cannot identify $13$. Its events identifying $12$ and $23$ are disjoint, each of probability $a$, hence $2a\le1$. On two positions the two parenthesizations of $(1,0,0)$ have equality probabilities $a^2$ and $a$. Naturality preserves each common block with probability one, forbids mergers separated by both inputs and assigns probability $a$ to each eligible block merger. The union bound over those mergers gives the estimate. Since $a\le1/2$, $a(1-a)\ge a/2$. If the associator is zero, all pair memberships equal those of $R\cap Q$ almost surely; finiteness gives equality of the whole relation.
\end{proof}
If commutativity is omitted while exact associativity and pointwise idempotence are retained, the only possibilities are intersection and the two input projections. To prove this, let $a$ and $b$ be the two mixed-input equality probabilities for ordered statuses $(1,0)$ and $(0,1)$. The same three-position transitivity construction gives $a+b\le1$. Associativity on $(1,0,0)$ and $(0,0,1)$ gives $a^2=a$ and $b^2=b$. Thus $(a,b)$ is $(0,0)$, $(1,0)$ or $(0,1)$. Restricting to every pair then determines the entire finite output relation almost surely as intersection, the first input, or the second input, respectively.

The second application must use its specified conditional kernel given the intermediate relation. An unrecorded persistent random variable can violate that assumption. A two-position population bound suffices for this size-dependent inequality; an all-carrier associativity bound is unnecessary. Exact all-size restriction naturality remains a substantive premise, not something established by finitely many successful tests.

For $G_a=(1-2a)\delta_{R\cap Q}+a\delta_R+a\delta_Q$, expansion gives the signed associator $a(1-a)(\delta_T-\delta_R)$ and uniform distance to intersection at most $2a$. A different kernel independently marks common blocks horizontal, vertical or neutral with probabilities $t,t,1-2t$ and merges identically marked blocks sharing the appropriate parent. It agrees with $G_{t^2}$ on every two-position input. On $(\top_3,\bot_3)$ its probabilities are $t^3$ for one block, $t^2(1-t)$ for each specified pair block, and $1-3t^2+2t^3$ for three singletons. The TV gap is $3t^2(1-t)$. At $t=1/10$ an exactly-two-block event has probability $27/1000$ and is impossible under $G_{t^2}$. Independent repetition gives size zero and power $1-(973/1000)^m$; $110$ is the first integer reaching $0.95$. This optimum concerns that specified pair and zero-size tests, not certification of all stochastic kernels.

\subsubsection{Small-carrier agreement and a growing-carrier obstruction}
\label{pat:marked-growing}

The marked-common-block kernel has a stronger contrast on growing
carriers.  Fix \(0<t\le1/2\).  Its common blocks receive independent
marks \(H,V,N\) with probabilities \(t,t,1-2t\); \(H\) blocks merge
within a parent block of the first input, and \(V\) blocks within a
parent block of the second.  These operations preserve the common
blocks, are commutative, are pointwise idempotent, and commute with
restriction: surviving common blocks inherit the same independent
marks.  Conditional composition samples the stated kernel anew given
the intermediate relation.

Let \(C_{n,u}\) select each of \(n\) positions independently with
probability \(u\), merge the selected positions into one block, and
leave other positions single.  Selecting zero or one position gives
the discrete relation \(\bot_n\).  For the marked kernel \(L_t\),
\[
 L_t(\top_n,\bot_n)=C_{n,t},
\]
and its two composed laws on \((\top_n,\bot_n,\bot_n)\) are
\[
 L_t(L_t(\top_n,\bot_n),\bot_n)=C_{n,t^2},\qquad
 L_t(\top_n,L_t(\bot_n,\bot_n))=C_{n,t}.
\]
In the first expression a position must survive two independent
selections to belong to the merged block.  In the second, the inner
operation leaves the discrete relation unchanged.

For fixed positive \(t\), the largest block divided by \(n\) converges
in probability to \(t^2\) under the first law and \(t\) under the second.
A threshold strictly between these values has probability difference
tending to one.  Therefore the supremum of the marked kernel's
associator TV over all finite carriers and triples is exactly one,
although its two-position defect is only \(t^2(1-t^2)\).
For every fixed cutoff \(N\), the defect on carriers of size at most
\(N\) is at most \(4\binom N2t^2\): each kernel application differs from
deterministic meet by at most \(\binom N2t^2\), and coupling the two
applications in each grouping to the common deterministic output
gives the bound.  Thus it can be arbitrarily small at a prescribed
finite cutoff.

There is also an exact moving-size limit.  If \(t\to0\) and \(nt\to c>0\),
\[
 \operatorname{TV}(C_{n,t^2},C_{n,t})
       \longrightarrow 1-(1+c)e^{-c}.
\]
To prove it, the chance that the \(t^2\) selection contains at least
two points is at most \(\binom n2t^4\to0\), so its partition law
converges in TV to the discrete relation.  The other law has distance
\[
 1-(1-t)^n-nt(1-t)^{n-1}
\]
from that relation, which tends to the displayed Poisson expression.
The triangle inequality transfers the limit to the two-law distance.

For an exact finite calculation, combine the probabilities of binomial
counts zero and one into one category, keep counts \(2,\ldots,n\)
separate, and take the TV distance of those two count laws.  Conditional
on a count at least two, each selected subset has the same probability,
so this is also the full partition-law distance.
These growing-carrier failures do not invalidate the earlier
size-dependent bound derived from the two-position parameter under
exact restriction naturality.  Finite tests still cannot certify that
all-size structural premise.

For uniform independent permutations, the join of their cycle partitions is the orbit partition of the generated subgroup. The chance of transitivity failure is $1/n+O(n^{-2})$, while a fixed pair is disconnected with probability $2/n^2+6/n^3+O(n^{-4})$. These are different events. One exact route to the first statement is the recurrence $p_1=1$ and
$p_n=1-\sum_{k=1}^{n-1}\frac{k p_k}{n\binom nk}$,
obtained by choosing the orbit of a specified point; the fixed-pair probability is
$\sum_{k=1}^{n-1}\frac{k(n-k)p_k}{n(n-1)\binom nk}$.
They describe this supplied random-permutation model, not a general law of relational connectivity.

\clearpage
\section{Statistical Comparisons and Reference Errors}
\label{sam:chapter}

An observation may preserve an algebraic distinction while leaving it difficult to estimate. Conversely, an exactly consistent report can erase the distinction altogether. This chapter separates source sampling, query sampling, report errors and errors in the comparison reference. All numerical guarantees concern the specified mathematical experiments. They are not results of physical trials.

\subsection{Collision statistics under a specified source}
\label{sam:collisions}
Let $X_1,\ldots,X_n$ be independent labels with a common discrete distribution $p=(p_j)$, and let $s_m=\sum_jp_j^m$. For $2\le m\le n$, define
\[
 U_{n,m}=\binom nm^{-1}\sum_{|S|=m}\mathbf1\{X_i\text{ are all equal for }i\in S\}.
\]
Then $\mathbb E U_{n,m}=s_m$. Counting pairs of subsets by their intersection gives the exact variance
\begin{equation}
 \operatorname{Var}(U_{n,m})=
 \sum_{r=1}^m\frac{\binom mr\binom{n-m}{m-r}}{\binom nm}
 (s_{2m-r}-s_m^2).
 \label{sam:uvariance}
\end{equation}
Indeed, disjoint subsets have independent indicators. If the intersection has $r\ge1$ positions, the two all-equal events together require all $2m-r$ positions to have one label. For $m=3$, this becomes
\begin{equation}
 V_p(n)=\frac{9(n-3)(n-4)(s_5-s_3^2)+18(n-3)(s_4-s_3^2)+6(s_3-s_3^2)}{n(n-1)(n-2)}.
 \label{sam:v3}
\end{equation}

Consider two fixed sources with finite cores $(1,4,6,7)/36$ and $(2,3,5,8)/36$, respectively, and in each case a common tail $p_j=2^{-j-1}$, $j\ge1$. Each core and tail has total mass $1/2$. Their moments are
\[
 s_m=\sum_{a\text{ in core}}(a/36)^m+\frac1{2^m(2^m-1)}.
\]
The means of the triple statistic are $\mu_0=425/13608$ and $\mu_1=439/13608$, with gap $\Delta=1/972$. Their pair-equality probabilities agree. Thus pair agreement alone cannot separate these two sources, whereas triple agreement can.

For a threshold $c$ between the means, the one-sided variance inequality gives
\begin{equation}
 \alpha\le\frac{V_0(n)}{V_0(n)+(c-\mu_0)^2},\qquad
 \beta\le\frac{V_1(n)}{V_1(n)+(\mu_1-c)^2}.
 \label{sam:cantelli}
\end{equation}
At the midpoint $c=2/63$, $n=365450$ makes both displayed bounds at most $1/20$. Allowing an unequal division of the gap improves this sufficient certificate: a threshold meeting both bounds exists exactly when
\[
 \sqrt{19V_0(n)}+\sqrt{19V_1(n)}\le\Delta.
\]
The first integer satisfying this inequality is $329180$. The admissible interval is
$[\mu_0+\sqrt{19V_0(n)},\,\mu_1-\sqrt{19V_1(n)}]$.
These are optima of the stated variance-based certificate, not minimum sample sizes among all tests. The inequality can be checked without floating-point decisions: put $a=19V_0$, $b=19V_1$; require $\Delta^2-a-b\ge0$ and $(\Delta^2-a-b)^2\ge4ab$. A complete $U$-statistic is the average of its leave-one-out versions, so Jensen's inequality proves that its variance cannot increase with sample size. Checking the preceding integer therefore establishes the first sufficient size for this certificate.

Changing $r$ raw labels can affect at most the fraction
$1-\binom{n-r}{3}/\binom n3$ of triples. This is a label-replacement bound. It must not be used as a bound for an arbitrary number of erroneous pair reports: those are different error models.

\subsection{Extending a statistic to inconsistent pair reports}
\label{sam:extension}
Let $z_{ij}\in\{0,1\}$ be arbitrary pair reports, where one means equality. They need not define an equivalence relation. For a $k$-set $S$, let $m_S(z)$ count its negative reports and set
\[
 \phi_k(z|_S)=\max\{0,1-m_S(z)/(k-1)\},\qquad
 W_{n,k}(z)=\binom nk^{-1}\sum_{|S|=k}\phi_k(z|_S).
\]
For $k=3$, a triple with zero, one, two or three positive reports receives score $0,0,1/2,1$.

\begin{proposition}[Sharp report sensitivity]
On every genuine equality graph, $W_{n,k}=U_{n,k}$. Moreover,
\[
 |W_{n,k}(z)-W_{n,k}(z')|
 \le\min\left\{1,\frac{k\,d_H(z,z')}{n(n-1)}\right\}.
\]
The coefficient $k/[n(n-1)]$ is the least possible global single-report Lipschitz constant among real-valued extensions agreeing with $U_{n,k}$ on every equality graph.
\end{proposition}
\begin{proof}
A nonconstant equality pattern on $k$ positions has at least $k-1$ unequal pairs, so its score is zero; a constant pattern has score one. Each pair lies in $\binom{n-2}{k-2}$ subsets and changes each local score by at most $1/(k-1)$, giving the upper bound by telescoping edits. The all-equal partition and the partition consisting of a singleton and a block of $n-1$ differ on $n-1$ reports and have statistic difference $k/n$. Along a shortest edit path any agreeing extension therefore has an edge difference at least $k/[n(n-1)]$. The displayed extension attains it, for example by removing a single edge from the all-equal graph.
\end{proof}
This is an extension and sensitivity result. It asserts neither uniqueness of the extension nor optimal testing risk.

Consistency of reports alone does not warrant their accuracy. The two finite cores above can both be deterministically grouped into weights $(5,13)/36$, leaving their common tails unchanged: group $(1,4)$ and $(6,7)$ in the first case and $(2,3)$ and $(5,8)$ in the second. The complete observed equality laws are then identical at every sample size. The reported comparisons remain perfectly transitive. Each implementation's pair-error probability against true equality is $23/324$, while their mutual disagreement can be zero. This construction distinguishes calibration against a warranted reference from agreement between implementations.

\subsection{Sparse random queries and the complete variance}
\label{sam:sparse}
The experiment is specified as follows. Draw $n$ independent source labels. Fix a complete potential pair-report graph $z$, possibly as a function of that entire sample and detector randomness. Its number of incorrect entries is at most $\rho\binom n2$. After this assignment is fixed, draw $B$ triples independently and uniformly with replacement, independently of the graph, and average their half-credit scores to obtain $\widehat W_{n,B}$. Persistent sample identity and stable access to the same report assignment are premises.

Conditional on the sample and fixed graph, the sampled scores are independent and have mean $W_{n,3}(z)$ and variance $v(z)/B$. They are generally dependent after averaging over the original source sample. In the error-free case,
\[
 \operatorname{Var}(\widehat W_{n,B})=\left(1-\frac1B\right)V_i(n)+\frac{\mu_i(1-\mu_i)}B.
\]
This follows by conditional variance; it is also the usual incomplete-$U$-statistic variance identity~\cite{clemencon}. Increasing $B$ alone does not remove $V_i(n)$.

Put $b=\min\{1,3\rho/2\}$. The preceding sensitivity bound gives
$|W_{n,3}(z)-U_{n,3}|\le b$. Introduce the analytical variable
$R=U_{n,3}+\widehat W_{n,B}-W_{n,3}(z)$.
Its mean is $\mu_i$ and its query noise is uncorrelated with $U_{n,3}$ by conditional centering. Since scores lie in $[0,1]$,
\[
 \operatorname{Var}R=V_i(n)+\frac{\mathbb Ev(z)}B
 \le V_i(n)+\frac{\mu_i+b}B=:A_i.
\]
For $d_0=c-\mu_0-b>0$ and $d_1=\mu_1-c-b>0$, the unconditional guarantees are
\begin{equation}
 \alpha\le\frac{A_0}{A_0+d_0^2},\qquad
 \beta\le\frac{A_1}{A_1+d_1^2}.
 \label{sam:sparseerrors}
\end{equation}
No independence among report errors is assumed. The variable $R$ is a proof device, not a correction available to an observer without the true labels.

For the two specified sources choose
\[
 n=2000000,\quad B=10000000,\quad\rho=1/10000,\quad
 b=3/20000,\quad c=99179/3125000.
\]
Exact substitution into \eqref{sam:sparseerrors} gives
\begin{align*}
 \alpha_s&=\frac{25032226015092744528840625}{634998927977903024107964257}
 =0.0394208949215168\ldots,\\
 \beta_s&=\frac{27900695463010851341340625}{699384194430517969365104257}
 =0.0398932313386483\ldots.
\end{align*}
At most thirty million pair-query calls are needed, compared with
$M=\binom n2=1999999000000$ potential pair entries. This trades more raw samples and additional query variance for fewer queried entries. It is not an optimum in total physical resources. The computation evaluates exact moments and inequalities; it does not simulate or experimentally collect those millions of observations.

The requirement that the graph precede the query randomization is essential. After learning the chosen triples, an adversary could mark every queried pair positive, changing at most $3B=30000000$ entries. This is below the admitted error budget $199999900$, yet every queried triple would score one. The complete graph average still obeys the deterministic sensitivity bound; the query average is no longer an unbiased sample of a previously fixed graph. Stable repeated responses do not repair a graph selected after disclosure of the query indices.

\subsection{Reference sampling with a quality gap}
\label{sam:quality}
Suppose a further, separately warranted reference reveals the true equality of a sampled pair. If a fixed report population has $E$ errors among $M$ entries, a uniform sample of $s$ distinct entries has hypergeometric discrepancy count $K$. In particular,
\[
 \Pr(K=0)=\frac{\binom{M-E}s}{\binom Ms}\le(1-E/M)^s.
\]
A zero-discrepancy acceptance rule with $s=46050$ has false-pass probability below $1/100$ for error fraction at least $10^{-4}$. At that same boundary its probability of accepting a good-enough graph is also at most about one percent. Consequently this particular rule cannot ensure high power for a subsequent detector throughout that boundary class.

A separated pair of quality requirements permits useful guarantees. Sample $s=200000$ distinct entries and accept when $K\le8$. For error fraction at least $10^{-4}$ and at most $10^{-5}$, respectively,
\begin{align}
 \Pr(\mathrm{pass})&\le(49997/50000)^{200000}(5/2)^8<47/5000,
 \label{sam:badpass}\\
 \Pr(\mathrm{fail})&\le(200007/200000)^{200000}(2/9)^9<29/20000.
 \label{sam:goodfail}
\end{align}
To derive these bounds, compare the convex exponential moment of a sample without replacement with that of independent draws from the same finite population, as in Hoeffding's Theorem 4~\cite{hoeffding}. For the lower tail take $e^{-t}=2/5$, so
$\Pr(K\le8)\le e^{8t}(1-p+pe^{-t})^s$.
For the upper tail take $e^t=9/2$, giving
$\Pr(K\ge9)\le e^{-9t}(1-p+pe^t)^s$.
Monotonicity in $p$ supplies the two boundary substitutions. The strict decimal-free comparisons in \eqref{sam:badpass}--\eqref{sam:goodfail} are integer inequalities.

If detection occurs only after acceptance, \eqref{sam:sparseerrors} and the quality bounds yield
\[
 \alpha_{\rm total}\le\alpha_s+47/5000<1/20,
 \qquad
 \mathrm{power}_{\rm total}\ge1-\beta_s-29/20000>19/20,
\]
where the power statement assumes a uniformly good alternative class with actual report fraction at most $10^{-5}$. The graph must be fixed before both independent randomizations. A data-dependent good/bad event does not justify conditioning the original sample as if it remained independent. For a rigorous unconditional argument let $G=\{E\le M/10000\}$ and replace $R$ by
\[
 R^*=U_{n,3}+\mathbf1_G(\widehat W_{n,B}-W_{n,3}(z)).
\]
Conditional centering still gives mean $\mu_i$ and variance at most $A_i$, since
$\mathbb E[\mathbf1_Gv(z)]\le\mathbb E[\mathbf1_GW_{n,3}(z)]\le\mu_i+b$.
On $G$ the sparse error event is controlled as before; on its complement the false-pass bound holds pointwise in the fixed sample and graph. The power bound follows by a union bound under the stronger actual-quality premise.

Acceptance is a repeated-sampling statement about a fixed finite population. It is not a posterior probability that a particular accepted graph is good, and it does not certify future observations under drift. If the reference $r$ itself can err relative to truth $t$, then
$d_H(z,t)\le d_H(z,r)+d_H(r,t)$.
The second term needs its own warrant; agreement with a shared coarsening does not provide one.

\subsection{The cost of separating adjacent error counts}
\label{sam:adjacent}
The next result concerns a different demand: uniform separation of two adjacent finite counts, with no quality gap. Fix an unknown binary error vector of length $M=1999999000000$, fixed before the randomization. A truthful query reveals one entry. There are no auxiliary features, structural restrictions on error locations, or informative side channels. Put $E_0=M/10000=199999900$. Require acceptance probability at least $19/20$ for every vector with at most $E_0$ errors and at most $1/100$ for every vector with at least $E_0+1$ errors. Randomized decisions are allowed.

\begin{proposition}[Minimum worst-case query cap]
In this experiment the exact minimum number of distinct queries is
\[
 s_{\min}=M-105=1999998999895.
\]
The same minimum applies to adaptive index selection with that worst-case cap, when choices use only previously revealed answers and private randomness. It is not a bound on expected stopping time.
\end{proposition}
\begin{proof}
Average any proposed procedure over a uniform permutation of the population labels. Uniform guarantees are preserved. Under the uniform arrangement of $E$ marked entries, after $t$ queries with $k$ observed errors any unqueried index has conditional error probability $(E-k)/(M-t)$. Therefore adaptivity cannot change the answer-sequence law. A rule using at most $s$ queries can be padded to $s$ distinct queries. The sufficient statistic is the count
\[
 P_E(K=k)=\frac{\binom Ek\binom{M-E}{s-k}}{\binom Ms}.
\]
This averaging proves a minimax necessity; it does not assume that nature selects a random arrangement.

Write $r=M-s$ and let $U_E$ count errors in the omitted set of $r$ positions. Then $K=E-U_E$ and
$\Pr(U_E=u)=\binom Eu\binom{M-E}{r-u}/\binom Mr$.
At the two boundary counts set
\[
 g_r=\Pr(U_{E_0}=0),\quad b_0=\Pr(U_{E_0+1}=0),\quad b_1=\Pr(U_{E_0+1}=1).
\]
On the common support, at $K=E_0-u$, the bad-to-good likelihood ratio is
\[
 \frac{(E_0+1)(r-u)}{(M-E_0)(u+1)},
\]
which increases with $K$. Exchanging accepted probability mass between two counts shows that accepting smaller counts first minimizes bad acceptance at a given good-acceptance requirement. At $r=105,106$, acceptance on $K<E_0$ alone gives less than $19/20$, and acceptance at $K=E_0$ with probability
$\gamma_r=1-1/(20g_r)$ attains that requirement exactly. Its minimum bad acceptance is
\[
 B_r=1-b_0-b_1+\gamma_rb_1.
\]
Exact rational evaluation gives $B_{105}=0.00992053452729221435\ldots<1/100$ but
$B_{106}=0.01001448946901366265\ldots>1/100$.
At $r=105$, the simpler decision accepting all $K<E_0$, accepting $K=E_0$ with probability $19/20$, and rejecting larger counts has boundary acceptance probabilities
$0.9505222793491606\ldots$ and $0.00992601900893398\ldots$.
A monotone coupling in $E$ extends those guarantees to every required count. Less information cannot improve the attainable testing region, since a larger sample can simulate a smaller one by discarding observations. Thus the failure at $r=106$ proves necessity for every smaller cap, and the rule at $r=105$ proves sufficiency.
\end{proof}

For perspective, adjacent hypergeometric laws satisfy the exact identity
\[
 \operatorname{TV}(P_E,P_{E+1})=\frac{s}{M}\max_k h(k),\qquad
 h=\operatorname{Hypergeom}(M-1,E,s-1).
\]
Coupling inclusion of the additional marked position gives
$P_{E+1}(k)-P_E(k)=(s/M)[h(k-1)-h(k)]$; unimodality telescopes the absolute differences. The weaker bound $\operatorname{TV}\le s/M$ alone does not yield the sharp cap above. Knowledge of a specified changed location, a constrained error pattern, additional reference data, an expected-cost objective, or a nonzero quality gap defines a different experiment.

\subsection{Interpretation and reproducibility}
\label{sam:interpretation}
The sufficient detection budgets and the nearly complete inspection requirement are compatible: the former permits separated quality classes, while the latter requires uniform classification across a one-entry boundary. Both depend on a truthful reference and a specified observation mechanism. Neither establishes a prephysical source law, an independently available copy mechanism, or a physical clock. Rejecting one of the two specified paintboxes is evidence within that source comparison, not a test of every proposed relational foundation.

The accompanying computations reconstruct the collision moments, variance bounds, sensitivity examples, rational sparse-error certificate and omitted-count probabilities. Strict decisions use rational or integer inequalities. Displayed decimal expansions aid reading. The unrestricted sensitivity, conditional-variance and minimax assertions rest on the proofs above; finite computation is a check on their instances, not a substitute for them.

\clearpage
\section{Quantum State Families, Symmetry and Relative References}
\label{qua:chapter}

This chapter uses ordinary complex qubit Hilbert spaces, their tensor products, density operators, partial traces and Born measurements as explicit assumptions. The examples compare the consequences of different state families and selection rules. None derives that quantum apparatus from signed labels or from a relational quotient. In particular, preserving one classical outcome law need not preserve transverse quantum observations, and separately consistent state families need not satisfy one common selection principle.

\subsection{A shared latent source and a supplied quantum representation}
\label{qua:shared}
Fix $0<r<1$ and $c=\sqrt{1-r^2}$. A fair latent variable $T\in\{r,-r\}$ and conditionally independent signs $S_i\in\{-1,1\}$ with $\mathbb E[S_i\mid T]=T$ give an exchangeable family at every size. All restrictions are consistent. The latent prior and conditional sampling rule are part of this construction; signed cancellation alone does not determine them.

On qubits define the pure states
\[
 \rho_\pm=\tfrac12(I+cX\pm rZ),\qquad
 \Omega_N=\tfrac12(\rho_+^{\otimes N}+\rho_-^{\otimes N}).
\]
Their $Z$ measurement law is the preceding classical law. The family is separable and consistent under arbitrary deletions. Choose phases of the branch vectors so their inner product is $c$. The two nonzero eigenvalues of $\Omega_N$ are
\[
 \lambda_\pm=(1\pm c^N)/2,
\]
because the nonzero spectrum of an equal mixture of two unit vectors is that of their $2\times2$ Gram matrix divided by two. Thus $\Omega_N$ is not faithful on the full $2^N$-dimensional algebra for $N\ge2$. A modular evolution requiring an invertible full-space state must be restricted to its support or supplied with an explicit faithful regularization; neither convention selects a physical clock.

Writing $\bar\rho=(\rho_++\rho_-)/2$, one obtains
\[
 \Omega_2-\bar\rho^{\otimes2}=\frac{r^2}{4}Z\otimes Z,
 \qquad D_{\rm tr}(\Omega_2,\bar\rho^{\otimes2})=r^2/2.
\]
Here and throughout $D_{\rm tr}(\rho,\sigma)=\tfrac12\|\rho-\sigma\|_1$.
The correlation matrix of $\Omega_2$ is $\operatorname{diag}(c^2,0,r^2)$. The optimized CHSH value for spin measurements is therefore
$2\sqrt{c^4+r^4}\le2$, by the two-largest-squared-singular-values criterion~\cite{horodecki}.
For $0<r<1$ the two branch states do not commute. The absence of CHSH violation is consequently not a statement that every quantum correlation notion vanishes; the state is separable but is not diagonal in a fixed product basis.

A static projective family is not a copying dynamics. If one branch-blind CPTP map sent each $\rho_\pm^{\otimes N}$ to $\rho_\pm^{\otimes(N+1)}$, monotonicity of squared fidelity would require $c^{2(N+1)}\ge c^{2N}$, which is false. A preparation that knows the branch, or access to additional correlated resources, changes this problem.

A commutative signed algebra remains commutative under any algebra representation. It cannot supply noncommuting $X$ and $Z$ merely by renaming its elements. The complex crossed product of the two-point function algebra with its sign-swap action is $M_2(\mathbb C)$: the diagonal minimal projections together with the swap generate all matrix units. The swap operation and complex algebra are additional inputs. Likewise, the supported double well $\lambda(t^2-r^2)^2$ vanishes on both allowed latent values, so its coupling does not affect the already chosen source law. The proposed cost from $Q=mZ$ and $R=X$ obeys $[Q,R]^*[Q,R]=4m^2I$ and cannot select a state or transverse direction.

\subsection{Phase-sensitive observations with an unchanged signed law}
\label{qua:phase}
The supplied diagonal unitary $U_\phi=\exp(i\phi Z\otimes Z)$ commutes with every $Z$-outcome projector. It preserves the complete signed law of a pair while changing other measurements. Direct Pauli conjugation yields
\begin{equation}
 U_\phi\Omega_2U_\phi^*=\tfrac14\{II+c\cos(2\phi)(XI+IX)
 -c\sin(2\phi)(YZ+ZY)+c^2XX+r^2ZZ\}.
 \label{qua:phaseform}
\end{equation}
Tensor signs have been suppressed only in this equation. At $\phi=\pi/4$, the two-party negativity is $c^2/2$ and the optimized CHSH value is $2\sqrt{1+c^4}$. These follow from the partial-transpose spectrum and from the correlation matrix
\[
 T_\phi=\begin{pmatrix}c^2&0&0\\0&0&-c\sin(2\phi)\\0&-c\sin(2\phi)&r^2\end{pmatrix}.
\]
At $\phi=\pi/4$ its eigenvalues are $c^2,1,-c^2$.
For the concrete value $r=1/2$, partial-transpose negativity appears whenever $\sin(2\phi)\ne0$, whereas CHSH violation requires
\[
 \sin^2(2\phi)>(7-\sqrt7)/12.
\]
Indeed, the larger eigenvalue of the lower block is
$[1+\sqrt{1+48\sin^2(2\phi)}]/8$ and the other relevant singular value is $3/4$. A negative partial transpose, a CHSH violation, and a usable physical Bell experiment are distinct requirements.

\begin{proposition}[A matched projective family]
For $t=\pm r$, let $\rho_t=(I+cX+tZ)/2$,
$\sigma_t=U_\phi\rho_t^{\otimes2}U_\phi^*$ and $\tau_t=\operatorname{Tr}_2\sigma_t$.
Then
\[
 \Xi_{2k}=\tfrac12\sum_{t=\pm r}\sigma_t^{\otimes k},\qquad
 \Xi_{2k+1}=\tfrac12\sum_{t=\pm r}\sigma_t^{\otimes k}\otimes\tau_t
\]
is consistent under deletion of its last factor and preserves the complete shared-source signed law. Its pair assignment is an additional ordering convention.
\end{proposition}
\begin{proof}
The pair state is symmetric in its two factors, and $\operatorname{Tr}\tau_t=1$. Tracing the last factor of an even member gives the preceding odd member; tracing the last factor of an odd member gives the preceding even member. Diagonal phases do not change the $Z$ law and $\tau_t$ has $Z$ mean $t$, so every displayed family member has the stated classical law. Importantly, the odd unpaired factor is $\tau_t$, not the original $\rho_t$. Substituting $\rho_t$ would generally destroy the asserted consistency.
\end{proof}
At $\phi=\pi/4$,
$\tau_t=\tfrac12(I-ctY+tZ)$.
Two factors from different pairs have state $\tfrac12\sum_t\tau_t\otimes\tau_t$, which is separable. Its difference from a designated pair is
\[
 \frac{c^2}{4}\{XX-c(YZ+ZY)-r^2YY\}.
\]
The eigenvalues of the expression in braces are $2+c^2$, $-c^2$, $-c^2$, and $c^2-2$. Hence the trace distance is
$c^2(2+c^2)/4=(1-r^2)(3-r^2)/4$, equal to $33/64$ at $r=1/2$.
Arbitrary deletion therefore need not preserve a named pairing, even though the signed probabilities agree.

Applying the same phase to every pair at each size is a different proposal. For
$\widetilde\Omega_N=(\prod_{i<j}e^{i\phi Z_iZ_j})\Omega_N(\prod_{i<j}e^{i\phi Z_iZ_j})^*$,
the two-site $XX$ expectation is $c^2$, while the same expectation in the two-site reduction of the three-site state is $c^2\cos^2(2\phi)$. If $\sin(2\phi)\ne0$, consistency fails. If $\phi$ is an odd multiple of $\pi/2$, a one-site $X$ expectation changes sign between those sizes. Only the trivial $\phi\in\pi\mathbb Z$ avoids this particular obstruction for $0<r<1$. This excludes the stated uniform all-pairs prescription, not every coherent family.

\subsection{A finite tail constrains a maximally transverse completion}
\label{qua:tail}
Suppose a state $\Gamma_{k+L}$ has exactly the full shared-source signed law above and satisfies $\langle X_i\rangle=c$ for every one of the first $k$ factors. Let $e_L$ be the minimum classification error for the fair sign $T$ from $L$ conditionally independent $Z$ observations of mean $T$. A majority rule, randomized fairly on ties, attains this error. Then
\begin{equation}
 D_{\rm tr}(\Gamma_k,\Omega_k)\le\sqrt{2kr^2e_L}.
 \label{qua:tailbound}
\end{equation}
To see the mechanism, measure the tail and denote its conditional first-$k$ state by $\gamma_y$. Guess the sign $s_y\in\{-1,1\}$. For each first factor the average failure probability of the pure projector $\rho_{s_yr}$ is
\[
 \tfrac12\{1-c\langle X_i\rangle-r\,\mathbb E[s_yZ_i]\}
 =\tfrac12\{1-c^2-r^2(1-2e_L)\}=r^2e_L.
\]
The product-projector failure is at most the sum of its single-site failures. The pure-target trace-distance bound and concavity of the square root then give \eqref{qua:tailbound} (in fact this argument permits the smaller right side $\sqrt{kr^2e_L}$). The guessed signs are fair, so the average pure target is exactly $\Omega_k$. All first-$k$ transverse constraints are used; a constraint only on the first factor does not establish the bound. Exact full signed laws, conditional tail sampling and arbitrary precision in those constraints are substantial assumptions. The estimate is a sufficient restriction on finite completions, not a claim that a particular finite size is necessary for every Bell test.

\subsection{Bell resources and pair structure}
\label{qua:bell}
For deterministic simultaneous signs $a_0,a_1,b_0,b_1$, the CHSH combination
$a_0(b_0+b_1)+a_1(b_0-b_1)$ has magnitude two. The same bound holds after mixing such local response assignments with setting independence. Quantum dichotomic observables obey $2\sqrt2$ by the squared-CHSH operator bound, while a general nonsignalling binary box can attain four, for example the law $a\oplus b=xy$ with uniform marginals. These bounds use different response assumptions; a classical set notation alone does not impose the first factorization.

For ordinary qubit factors, let $\beta_{ij}$ denote the CHSH value optimized over nontrivial local Pauli spin observables. Independently optimized pair values obey
$\beta_{AB}^2+\beta_{AC}^2\le8$. The independently optimized qualification uses the strong three-qubit monogamy result of Cheng and Hall~\cite{chenghall}; a common-settings monogamy statement alone would not justify it. Hence the graph of pairs satisfying $\beta_{ij}>2$ is a matching. If one defines the nonlinear objective
\[
 A_N=-\sum_{i<j}\max\{0,\beta_{ij}-2\},
\]
then
$A_N\ge-\lfloor N/2\rfloor(2\sqrt2-2)$, attained by products of Bell pairs. A deficit smaller than $2\sqrt2-2$ from that optimum forces a maximum-cardinality violation matching. This follows by counting at most $\lfloor N/2\rfloor$ nonzero terms and applying the single-pair bound. It does not select a labelled matching.

The objective is not the expectation of a fixed Hamiltonian: both measurement optimization and the positive-part operation are nonlinear in the state. In an equal mixture of two distinct four-qubit singlet matchings, each edge present in one matching has Werner state
$\tfrac12P_{\rm sing}+\tfrac12I_4/4$.
Its negativity is $1/8$ and its optimized spin-CHSH value is $\sqrt2$, so it contributes no positive Bell excess. The other pairs are maximally mixed. This example makes the loss under mixing explicit without claiming that every symmetric nonlinear objective uniquely selects Bell pairs.

\subsection{Collective spin and the distribution of singlet weight}
\label{qua:spin}
Let $J_N=\tfrac12\sum_{i=1}^N\boldsymbol\sigma_i$ and $H_N=J_N^2$. Its eigenvalues are $j(j+1)$, with multiplicities
\[
 m_{N,j}=\binom N{N/2-j}-\binom N{N/2-j-1}
\]
for the allowed spins $j$. This follows by subtracting consecutive fixed-$J_z$ dimensions: each spin-$j$ copy contributes one vector to every weight from $-j$ to $j$. In particular, the minimum is zero for even $N$ and $3/4$ for odd $N$. The even zero-spin multiplicity is the Catalan number $C_k=\binom{2k}k/(k+1)$ for $N=2k$. For $N=2k-1$, the minimum-sector dimension is $2C_k$; for $N=2k$, it is $C_k$. At $N=2$ the singlet is unique and the gap to the triplet is two.

The Pauli identity gives
\[
 H_N=\frac{3N}{4}I+\frac12\sum_{i<j}\boldsymbol\sigma_i\cdot\boldsymbol\sigma_j,
 \qquad
 \sum_{i<j}P_{\rm sing}^{(ij)}=\frac{N(N+2)}8I-\frac12H_N.
\]
Consequently the mean pair singlet fraction is at most
$\frac{N+2}{4(N-1)}$ for even $N$ and $\frac{N+3}{4N}$ for odd $N$. Every state supported in the minimum sector saturates the corresponding average bound. The identity does not fix how that weight is distributed among pairs.

For four qubits, the product of singlets on pairs $(12),(34)$ has singlet fraction one on pair $12$. Coupling the two pair-triplet sectors to total spin zero gives an orthogonal total-singlet vector whose pair-$12$ reduction is the uniform triplet state, with singlet fraction zero. Superpositions of these vectors realize every fraction in $[0,1]$ at the same global energy zero: after tracing pair $34$, cross terms vanish because its singlet and triplet sectors are orthogonal. An arbitrarily small additional operator can therefore select substantially different pair structures within the same degenerate ground space.

\begin{proposition}[No unique symmetric ground ray for three or more qubits]
A Hamiltonian commuting with diagonal $SU(2)$ and every factor permutation has no unique pure ground ray when $N\ge3$.
\end{proposition}
\begin{proof}
A unique ray carries a one-dimensional representation of each symmetry. The connected group $SU(2)$ has no nontrivial continuous one-dimensional character, so the vector must be a spin singlet. A one-dimensional representation of $S_N$ is either symmetric or alternating. The symmetric qubit subspace has spin $N/2$; the alternating subspace is $\bigwedge^N\mathbb C^2=0$ for $N>2$. Neither can contain the required ray. This does not forbid a unique invariant mixed state or a symmetry-breaking preparation.
\end{proof}

\subsection{Ordered dimers, symmetric mixtures and restrictions}
\label{qua:projective}
An explicitly chosen ordered matching gives
\[
 \Gamma_{2k}=P_{\rm sing}^{\otimes k},\qquad
 \Gamma_{2k+1}=P_{\rm sing}^{\otimes k}\otimes I_2/2.
\]
These states minimize $H_N$ and are consistent under last-factor deletion. Thus failure of one symmetric family does not prove failure of all minimizing projective families. The matching remains a supplied preparation choice.

Write $P_{\min,N}$ for the minimum-spin projector and $d_N$ for its rank. Averaging $\Gamma_N$ over permutations gives
$\rho_N^*=P_{\min,N}/d_N$.
Within the relevant spin sector, permutation irreducibility on its multiplicity space and the already maximally mixed spin factor give this identity. Even-to-odd reduction of this canonical family is consistent. Odd-to-even reduction is not, and
\begin{equation}
 D_{\rm tr}\bigl(\operatorname{Tr}_{2k+1}\rho_{2k+1}^*,\rho_{2k}^*\bigr)
 =\frac{3k}{2(2k+1)}.
 \label{qua:oddgap}
\end{equation}
Only spins zero and one of the first $2k$ factors can couple to total spin $1/2$. The spin-zero contribution to the odd minimum sector has dimension $2C_k$, while the full sector has dimension $2C_{k+1}$. Its weight is therefore $C_k/C_{k+1}=(k+2)/[2(2k+1)]$. The other contribution has orthogonal support, giving \eqref{qua:oddgap}.

For $N=2k-1$ or $2k$, the dimer/canonical distances are
\[
 D_{\rm tr}(\Gamma_N,\rho_N^*)=1-1/C_k,
 \qquad
 \operatorname{TV}(Z\text{ laws})=1-\frac{2^k}{\binom{2k}k}.
\]
For even $N$, the dimer state has rank one inside a $C_k$-dimensional uniformly mixed sector. Its $Z$ law is uniform on the $2^k$ words with one positive sign per specified pair, while the canonical law is uniform on the $\binom{2k}k$ balanced words. For odd $N$, the dimer rank is two inside dimension $2C_k$, and the same ratio follows from the two nearest-balanced count sectors. These are global quantum and labelled classical distances, respectively. If permutations are declared gauge redundancies, named-pair measurements require a relative reference; if they are physical symmetries, choosing a matching is additional state preparation. The formal Casimir does not decide between those interpretations.

Another family, the maximum-spin mixture
\[
 \Sigma_N=\frac{P_{{\rm sym},N}}{N+1}
 =\int(|\psi\rangle\langle\psi|)^{\otimes N}\,d\psi,
\]
uses normalized Haar measure on pure qubits. Invariance makes the integral scalar on the irreducible symmetric subspace, and its trace fixes the scalar. The formula proves separability and arbitrary-deletion consistency. In a $Z$ measurement the positive count is uniform on $0,\ldots,N$ because the pure-state population is uniform on $[0,1]$ and its beta-binomial integral is $1/(N+1)$. Distinct-site odd signed moments vanish and the even moment of order $2m$ is $1/(2m+1)$. Exchangeability thus permits correlation; it does not mean independent unconditioned draws.

\subsection{Infinite extension and finite-temperature limits}
\label{qua:exchangeability}
An indefinitely exchangeable complex quantum state has the de Finetti form
$\omega_N=\int\tau^{\otimes N}\,d\mu(\tau)$~\cite{caves}.
For $N\ge2$, the canonical minimum-spin state satisfies the exact distance formula
\begin{equation}
 \inf_{\omega_N\text{ indefinitely exchangeable}}
 D_{\rm tr}(\rho_N^*,\omega_N)=1-d_N/2^N.
 \label{qua:distanceexchange}
\end{equation}
To prove it, diagonalize $\tau$ by a common rotation, with eigenvalues $p,1-p$. For $N=2k$, the minimum-sector overlap is $d_N[p(1-p)]^k$; for $N=2k+1$, it is $(d_N/2)[p(1-p)]^k$. In either case it is at most $d_N/2^N$. Measuring the minimum-sector projector yields the lower bound, and the product state $I/2^N$ attains it because $\rho_N^*$ is uniform on that projector. The same lower bound holds for any state wholly supported there, but attainment is asserted only for the canonical uniform state. Arbitrary finite permutation invariance is a weaker requirement than indefinite extension.

The thermal prescription $\rho_N(\beta)\propto e^{-\beta H_N}$ also need not be projective. At one common finite $\beta>0$, the two-site singlet probabilities of $\rho_2(\beta)$ and of the reduced $\rho_3(\beta)$ are
\[
 p_2=\frac1{1+3e^{-2\beta}},\qquad
 p_{3\to2}=\frac1{2(1+e^{-3\beta})}.
\]
Both states are Werner states, so their trace distance is their singlet-probability difference. With $x=e^{-\beta}$,
\[
 p_2-p_{3\to2}=\frac{(1-x)^2(1+2x)}{2(1+3x^2)(1+x^3)}>0.
\]
At $\beta=\log2$ this is $8/63$. At $\beta=0$ the family is the consistent product $I/2^N$. No statement follows here about size-dependent temperatures, other Hamiltonians or other thermal models.

Small normalized energy does not control a global state distance. The product $I/2^N$ has $\operatorname{Tr}(H_NI/2^N)=3N/4$, whose value per unordered pair is $3/[2(N-1)]\to0$, while its distance from $\rho_N^*$ is $1-d_N/2^N\to1$. Conversely, the shared pure-product family $\Omega_N$ lies entirely in maximum spin and has energy $N(N+2)/4$. At $\phi=\pi/4$ the phase-pair state is triplet-supported with pair energy two. For the matched family of Section~\ref{qua:phase}, specifically at $\phi=\pi/4$,
\[
 \operatorname{Tr}(\Xi_NH_N)=\frac{3N}{4}+\frac{k}{2}
 +\frac12\left[\binom N2-k\right]r^2(2-r^2),\qquad k=\lfloor N/2\rfloor.
\]
The paired correlation sum is one and every other pair has sum $r^2(2-r^2)$, proving the formula. These families cannot be generated by one unchanged minimum-Casimir prescription. Choosing a different objective is possible, but must be stated as an additional model.

\subsection{Matching orbits and shared references}
\label{qua:references}
On four qubits the three two-singlet matchings span the two-dimensional total-singlet space. Let their density operators be $\eta_1,\eta_2,\eta_3$. Their pair overlaps are one on the diagonal and $1/4$ off it; their equal mixture is the maximally mixed state on that space. These identities follow by expanding the singlets in the computational basis. The permutation group $S_4$ acts transitively on the three matchings, and a matching has stabilizer of order eight. That stabilizer contains odd permutations, so permutation parity does not define a function of the matching alone.

For two independent systems in specified density states $\rho,\sigma$, the symmetric SWAP effect has probability
$[1+\operatorname{Tr}(\rho\sigma)]/2$.
It gives one for aligned pure matchings and $5/8$ for distinct matchings. A common random matching reference produces
$\frac13\sum_i\eta_i\otimes\eta_i$, whereas independently randomized references produce
$(\frac13\sum_i\eta_i)^{\otimes2}$.
Their trace distance is $1/4$ and their SWAP probabilities are one and $3/4$. In a Bloch representation the three rays form a planar trine; the difference matrix has eigenvalues $1/4,0,0,-1/4$, which proves the distance. Shared randomization is not independent randomization.

Relative references can be described without assigning an absolute frame. If $U$ is a unitary group representation, the overlap
$\operatorname{Tr}(\rho U_r\sigma U_r^*)$ is invariant under
\[
 \rho\mapsto U_g\rho U_g^*,\quad
 \sigma\mapsto U_h\sigma U_h^*,\quad r\mapsto grh^{-1}.
\]
Compatible three-reference assignments satisfy $r_{13}=r_{12}r_{23}$. Independent vertex changes remove all edge labels on a connected tree by fixing a root and propagating frames. On a graph with cycles, the residual loop products transform by simultaneous conjugation. Their individual conjugacy classes, or their parities, need not determine the joint loop tuple. For example, in $S_3$ two noncommuting transpositions and two equal transpositions have the same individual conjugacy classes but different product classes. This is formal group-valued transport, not a derivation of physical gauge fields or curvature.

The matching orbit has two stabilizer double-coset classes, $E$ and $D$, of sizes eight and sixteen. Uniform independent lifts induce
$D*D=(E+D)/2$.
On the three-element orbit, a $D$ move chooses uniformly one of the two other matchings; two such moves return with probability $1/2$. This stochastic product is associative because it comes from convolution with the uniform stabilizer measure. It is not deterministic quotient multiplication. A groupoid representation preserves inverses; noninvertible erasure needs extra mixing or an instrument and cannot be called a groupoid arrow merely by changing its name.

\subsection{An orbit alphabet is more restrictive than a positive matrix}
\label{qua:packing}
For $m\ge2$ preparations chosen from the three exact matching rays, let $n_1+n_2+n_3=m$ be their counts. The sum of pair overlaps is
\[
 \sum_{a<b}\operatorname{Tr}(\eta_{i_a}\eta_{i_b})
 =\frac{3\sum_i n_i^2+m^2-4m}{8}.
\]
It is minimized by balanced counts, giving
\[
 \frac{m(m-2)}4+\begin{cases}0,&3\mid m,\\1/4,&3\nmid m.\end{cases}
\]
For four matching preparations the minimum is therefore $9/4$. Allowing arbitrary density states on the same two-dimensional support lowers the sharp minimum to $m(m-2)/4$, equal to two at $m=4$. If $S=\sum_i\rho_i$, then
\[
 \sum_{i<j}\operatorname{Tr}(\rho_i\rho_j)
 =\tfrac12\{\operatorname{Tr}S^2-\sum_i\operatorname{Tr}\rho_i^2\}
 \ge\tfrac12(m^2/2-m),
\]
and equally spaced pure states on a Bloch great circle attain it. Positive-semidefinite candidate overlap data need not be realizable by a specified three-state alphabet. If the matching labels are indefinitely exchangeable, their conditional independent representation instead gives mean pair overlap
$1/4+(3/4)\sum_i p_i^2\ge1/2$, also after averaging over $p$.

Pair overlaps do not fix an oriented triple invariant. For three pure Bloch vectors along $+X,+Y,+Z$, all pair overlaps are $1/2$, but
\[
 \operatorname{Im}\operatorname{Tr}(\rho_1\rho_2\rho_3)=1/4.
\]
Reflecting $Y$ reverses this sign without changing pair overlaps. On three copies define the cyclic permutation $V$ with orientation chosen so $\operatorname{Tr}(V\rho_1\otimes\rho_2\otimes\rho_3)=\operatorname{Tr}(\rho_1\rho_2\rho_3)$, and let
$K=(V-V^*)/(2i)$, $E=(I+K)/2$.
Since the spectrum of $K$ is contained in $\{0,\sqrt3/2,-\sqrt3/2\}$, $E$ is an effect. It gives probabilities $5/8$ and $3/8$ for the two orientations. Averaging over every permutation removes this oriented effect's contrast, because odd permutations reverse $K$. It does not make the entire anonymized preparations identical: a calibrated $+Y$ single-factor effect averaged over the three positions has probabilities $2/3$ and $1/3$. Removing one reference-dependent observable is weaker than equality of all admitted observations.

All these comparisons require a declared group action, copy convention and effect menu. A formal matching, relative phase or loop variable does not by itself identify physical adjacency, charge, chirality, time or spatial geometry.

\clearpage
\section{Observation, Memory, Control, and Time in Finite Models}
\label{op:chapter}

This section studies what a specified collection of observations determines about
a finite mathematical model.  Its objects are conventional probability kernels,
finite hidden-state models, quantum states, and continuous-time Markov generators.
An observation law, an intervention, and a calibration premise are distinct
inputs.  Equality of observable laws will always refer to an explicitly stated
experimental menu.  Identifiability within such a menu does not establish that
the chosen states, controls, environmental factors, or time coordinate describe
a particular physical system.

We use row probability vectors.  For probability measures on a finite set,
\(\operatorname{TV}(p,q)=\frac12\sum_x|p_x-q_x|\).
Matrix \(1\)-norms below are entrywise unless specified otherwise; the Frobenius
norm is denoted by \(\|\cdot\|_F\).  The vector \(\mathbf1\) has all entries one.
All logarithms in likelihood calculations are natural; binary entropy is
explicitly measured in bits.  Population probabilities are exact mathematical
inputs unless a finite-sample statement is given.

\subsection{Controlled reduction and the limits of short records}
\label{op:reduction}

Let \(X,Z\) be finite sets, \(q:X\to Z\) a surjection, and \(P_a\) a
row-stochastic transition matrix on \(X\) for each allowed action \(a\).
The full process is Markov conditional on the selected action.  A policy may
randomize and may depend on the complete visible history, but it has no extra
model-specific hidden-state information.  Write
\[
 M_a(x,z')=\sum_{y:q(y)=z'}P_a(x,y).
\]
The familiar strong-lumpability condition is the action-free version of the
following statement; see Wolfer and Watanabe~\cite{op:wolfer}.

\begin{theorem}[Exact controlled reduction]\label{op:lumpability}
There are kernels \(Q_a\) on \(Z\) that reproduce the complete visible-history
law for every initial distribution on \(X\) and every common visible-history
policy if and only if, for every action \(a\), \(M_a(x,\cdot)\) is constant
on each fiber of \(q\).  In that event the common row is
\(Q_a(q(x),\cdot)\).
\end{theorem}
\begin{proof}
For necessity, choose point initial distributions at two states \(x,x'\) in
one fiber and apply the same fixed action.  Both reduced initial states are
\(q(x)\); their next-observation laws must therefore be equal.
For sufficiency, condition on any visible history of positive probability
and the next selected action.  The conditional law of the hidden state is
supported in the fiber of the current observation.  Mixing identical pushed
rows over that posterior gives the stated \(Q_a\) row.  The policy uses the
same conditional distribution of actions in both descriptions.  Induction
over the finite history proves equality, including the recorded actions.
\end{proof}

If only specified starting states are admitted, the fiber condition need
only be imposed on the appropriately reachable domain.  Reduction for one
special initial law is a weaker property and does not imply the all-start
statement.

\begin{example}[Fair pairs with predictive memory]\label{op:four-state}
On states \(00,01,10,11\), observe the final bit and use
\[
 P=\begin{pmatrix}
 3/4&1/4&0&0\\
 0&0&1/4&3/4\\
 1/4&3/4&0&0\\
 0&0&3/4&1/4
 \end{pmatrix}.
\]
The uniform distribution is stationary because each column sums to one.
Both append choices have positive probability; the shift graph is strongly
connected and has self-loops, so the chain is irreducible and aperiodic.
In stationarity every adjacent observed pair has probability \(1/4\).
Nevertheless,
\[
 P(Y_1=a,Y_2=b,Y_3=c)
 =\frac18\left[1+\frac12(-1)^{a+b+c}\right],
 \qquad
 \operatorname{TV}(P_{Y_{1:3}},\operatorname{Unif}\{0,1\}^3)=\frac14 .
\]
For example, the next-zero probabilities after \(00\) and \(10\) are \(3/4\)
and \(1/4\), respectively.  The visible process is not first-order Markov.
The best uniform row approximation by a kernel depending on the last bit
has error \(1/4\): the two relevant Bernoulli rows have distance \(1/2\),
and their fair midpoint attains half that distance.
\end{example}

Hidden state is not by itself an obstruction.  For example,
\[
 P((z,h),(z',h'))=Q(z,z')\,\mathbf1_{\{h'=1-h\}}
\]
has a perpetually alternating hidden bit and an exactly autonomous visible
kernel \(Q\).

\begin{theorem}[Uniform row error controls complete histories]
\label{op:history-bound}
Suppose, uniformly over all admitted \(x,a\),
\[
 \operatorname{TV}\bigl(M_a(x,\cdot),Q_a(q(x),\cdot)\bigr)\le\delta,
 \qquad 0\le\delta\le1.
\]
Start with the same visible law and use the same visible-history policy.
After \(N\) transitions the total variation distance between complete
visible histories, including actions, is at most
\[
 1-(1-\delta)^N\le N\delta .
\]
The first bound is sharp.
\end{theorem}
\begin{proof}
Couple the initial observations identically.  As long as the histories
agree, use the same policy randomness.  The next visible law of the full
model is a posterior mixture of pushed rows, hence is within \(\delta\)
of the reduced row by convexity of total variation.  A maximal coupling
keeps the next observations equal with probability at least \(1-\delta\).
Consequently the coupled histories agree with probability at least
\((1-\delta)^N\).  The coupling inequality proves the claim.
For sharpness, let the reference chain remain at state \(0\), with state
\(1\) absorbing if ever entered.  Let the comparison chain escape from
\(0\) to absorbing \(1\) with probability \(\delta\) on each step.
The row premise holds at both states.  The reference history is all zero,
and its probability under the comparison model is exactly
\((1-\delta)^N\).
\end{proof}

Thus a test of size \(\alpha\) for the reduced history law has size at most
\(\min\{1,\alpha+1-(1-\delta)^N\}\) for the full model.  An average
stationary row error supplies no such uniform conditional premise.

An intervention on fine states need not define an intervention on coarse
states.  As a simple exact example, let \(q(1)=q(2)=0\), \(q(3)=q(4)=1\).
The point masses at \(1,2\) have identical coarse images.  A permutation
interchanging \(2,3\) while fixing \(1,4\) sends them to distinct coarse
point masses.  No map of the original coarse distribution can represent
that operation on both inputs.  The quantum counterpart is equally direct:
\(Z\)-basis dephasing sends both \(|+\rangle\langle+|\) and
\(|-\rangle\langle-|\) to \(I/2\), while a Hadamard operation followed by
that dephasing distinguishes them perfectly.  Specifying the admitted
interventions is therefore part of specifying a reduced model.

\begin{example}[A full-support non-descent witness]\label{op:non-descent}
For the same partition \(\{1,2\},\{3,4\}\), take the strictly positive
fine laws
\[
 p=(3/5,3/20,1/5,1/20),\qquad
 p'=(3/20,3/5,1/5,1/20).
\]
Both have coarse law \((3/4,1/4)\).
Interchanging the second and third fine labels sends them to coarse
laws \((4/5,1/5)\) and \((7/20,13/20)\), of total variation \(9/20\).
Thus the failure is not confined to boundary point masses.
Exactly eight of the twenty-four fine permutations descend for all
fine distributions: the permutation must send each whole input block
to one output block, giving two block permutations and two internal
permutations in each block, hence \(2(2!)^2=8\).
\end{example}

\subsection{A parity-shift process and its exact innovation coordinates}
\label{op:parity}

Fix \(L\ge1\) and \(-1<\eta<1\).  The hidden state is the last \(L\) bits,
initially uniform.  From \(x=(x_1,\ldots,x_L)\), append \(b\) with
probability
\begin{equation}\label{op:parity-kernel}
 P_\eta\bigl(x,(x_2,\ldots,x_L,b)\bigr)
 =\frac{1+\eta(-1)^{x_1+\cdots+x_L+b}}2 .
\end{equation}
This family is a parity form of the two-faced chains studied by
Ryabko~\cite{op:ryabko}, with parameter \(p=(1+\eta)/2\).

\begin{proposition}[Finite-depth agreement]\label{op:finite-depth}
The uniform law on \(L\)-bit states is stationary.  Every contiguous
observed block of length at most \(L\) is uniform, whereas a block of
length \(L+1\) has law
\[
 2^{-(L+1)}\left[1+\eta(-1)^{\sum_{j=1}^{L+1}y_j}\right].
\]
Its distance from the uniform law is \(|\eta|/2\).
For \(\eta\ne0\), the observed process has Markov order exactly \(L\).
\end{proposition}
\begin{proof}
Each destination state has two predecessors, differing in their dropped
bit.  Their transition probabilities sum to one.  Thus a uniform state
distribution remains uniform.  Marginalization gives every shorter
contiguous block law.  Multiplying the initial \(2^{-L}\) probability by
the append probability gives the displayed \(L+1\) law, whose \(2^{L+1}\)
absolute discrepancies are all \(2^{-(L+1)}|\eta|\).
The transition rule depends on no more than the last \(L\) bits.
For \(\eta\ne0\), changing the oldest of those bits while keeping the last
\(L-1\) fixed changes the next-bit probability.  All such histories have
positive probability, proving that smaller order is impossible.
\end{proof}

The statement concerns contiguous blocks.  It does not assert independence
of every collection of \(L\) widely separated observations.  Since \(L\)
can be arbitrarily large, no common finite validation depth distinguishes
the whole family from fair independent bits.

For a record of length \(L+m\), define parity innovations
\[
 E_t=Y_t\mathbin{\oplus}\cdots\mathbin{\oplus}Y_{t+L},
 \quad 1\le t\le m,
\]
where \(\oplus\) is addition modulo two.

\begin{theorem}[Innovation factorization]\label{op:innovations}
The map from \(Y_{1:L+m}\) to \((Y_{1:L},E_{1:m})\) is bijective and
\[
 P(Y_{1:L}=x,E_{1:m}=e)
 =2^{-L}\prod_{t=1}^m\frac{1+\eta(-1)^{e_t}}2.
\]
Consequently the innovations are independent and identically distributed,
with \(P(E_t=0)=(1+\eta)/2\), and are independent of the initial block.
\end{theorem}
\begin{proof}
The recursion
\(Y_{t+L}=E_t\oplus Y_t\oplus\cdots\oplus Y_{t+L-1}\)
recovers the record uniquely.  Its probability is the initial block
probability times the \(m\) append probabilities in
\eqref{op:parity-kernel}.  Each append factor depends only on \(e_t\),
giving the asserted product law.  Bijection alone would not establish
independence; the factorization of probabilities is essential.
\end{proof}

Writing \(h_2\) for binary entropy in bits and \(d_{\rm Ber}\) for
Bernoulli relative entropy in nats, this gives
\[
 H(Y_{1:L+m})=L+m h_2((1+\eta)/2),\qquad
 D(P_{Y_{1:L+m}}\Vert\operatorname{Unif})=
 m d_{\rm Ber}((1+\eta)/2\Vert1/2).
\]
These are entropy and divergence per indexed observation, not a physical
communication capacity or a rate per unit of laboratory time.
Cancellation of overlapping parities also gives
\[
 (-1)^{Y_t+Y_{t+L+1}}=(-1)^{E_t+E_{t+1}},\qquad
 P(Y_t=a,Y_{t+L+1}=b)=\tfrac14[1+\eta^2(-1)^{a+b}].
\]
The two retained outcomes therefore have distance \(\eta^2/2\) from a
fair independent pair, although their acquisition spans \(L+2\) process
positions.  Repeated overlapping sparse pairs are not automatically
independent trials.
For \(L\ge2\), pairs at lags \(1,\ldots,L\) are independent and fair:
for lag \(L\), marginalizing any other coordinate of the \(L+1\) parity
law removes its perturbation, and smaller lags follow from the uniform
block laws.  The qualifier \(L\ge2\) is necessary.

\subsubsection{Ideal fixed-horizon testing}

\begin{lemma}[Finite likelihood-ratio ordering]\label{op:np}
For two strictly positive probability laws \(p_0,p_1\) on a finite
sample space, a randomized test of prescribed size \(\alpha\) has maximum
power by rejecting in decreasing order of \(p_1/p_0\), randomizing only
on a boundary level.
\end{lemma}
\begin{proof}
If the threshold is \(c\) and the proposed test is \(\phi_*\), then
\((\phi_*-\phi)(p_1-cp_0)\ge0\) pointwise for every other test
\(0\le\phi\le1\).  Summation and equal size imply
\(\sum\phi_*p_1\ge\sum\phi p_1\).  A boundary randomization fills the
required size.  This is the finite Neyman--Pearson argument
\cite{op:neyman}.
\end{proof}

With known \(L\), error-free readout, and alternatives \(\eta=0\) versus
\(\eta=1/2\), let \(C_m=\sum_{t=1}^m\mathbf1_{\{E_t=0\}}\).
The likelihood ratio of the complete record is \(3^{C_m}/2^m\);
the initial block is ancillary.  Thus
\[
 C_m\sim\operatorname{Bin}(m,1/2)\quad\hbox{or}\quad
 \operatorname{Bin}(m,3/4).
\]
At size \(1/20\), the optimum powers for \(m=39,40\) are respectively
\(0.9498706279714081\ldots\) and \(0.9530027582141789\ldots\).
At \(m=40\), reject for \(C_m\ge26\) and, when \(C_m=25\), reject with
probability \(4423136507/16760560440\).
The exact power is
\[
 \frac{720068525355674046204729}{755578637259143234191360}.
\]
Optimal full-record power cannot decrease with the horizon, because an
extra observation can be discarded.  The failure at \(39\) and success
at \(40\) therefore establish the minimum within this randomized
fixed-horizon experiment.  It uses \(L+40\) raw outputs.  The nonrandomized
rule \(m=42,C_m\ge27\) is a sufficient alternative, with size
\(0.04421477349\ldots\) and power \(0.95837130194\ldots\); no global
nonrandomized minimum is asserted.

\subsubsection{Independent readout errors create dependent innovations}

Now observe \(\widetilde Y_t=Y_t\oplus N_t\), where the \(N_t\) are iid
Bernoulli\((\epsilon)\), independent of the entire source, and
\(0\le\epsilon<1/2\).  Put \(a=1-2\epsilon\),
\(W_t=(-1)^{E_t}\), \(V_j=(-1)^{N_j}\), and
\[
 Z_t=(-1)^{\widetilde E_t}
     =W_t\prod_{j=t}^{t+L}V_j.
\]

\begin{proposition}[Noisy innovation law]\label{op:noisy-cov}
The mean and covariance function are
\[
 \mu=\eta a^{L+1},\quad \Gamma_0=1-\mu^2,\quad
 \Gamma_h=
 \begin{cases}
 \eta^2(a^{2h}-a^{2L+2}),&1\le h\le L,\\
 0,&h\ge L+1.
 \end{cases}
\]
Collections of \(Z\)'s separated by more than \(L\) indices are
independent.  For \(C_m=\sum_{t=1}^m(1+Z_t)/2\),
\[
 \operatorname{Var}(C_m)=\frac14
 \left[m\Gamma_0+
 2\sum_{h=1}^{\min\{L,m-1\}}(m-h)\Gamma_h\right].
\]
\end{proposition}
\begin{proof}
The \(W_t\) are iid of mean \(\eta\), while the independent \(V_j\) have
mean \(a\).  In \(Z_tZ_{t+h}\), the \(L+1-h\) common detector factors
square to one when \(h\le L\), leaving \(2h\) factors and expectation
\(\eta^2a^{2h}\).  For larger \(h\), the defining independent variables
are disjoint, giving expectation \(\mu^2\) and the stronger setwise
independence statement.  Summing the covariances gives the variance.
\end{proof}

The observed initial \(L\)-bit block remains uniform and independent of
the complete observed innovation sequence: condition on all innovations
and all detector bits, and use the uniform independent source initial
block.  Thus the ordered noisy innovations are sufficient within this
specified family.  Their count generally is not.  For
\(L=2,\eta=1/2,\epsilon=1/10\),
\[
 P(\widetilde E_{1:3}=001)=681/5000,\qquad
 P(\widetilde E_{1:3}=010)=609/5000,
\]
although both words have the same count.  At \(m=40\), the count mean is
\(628/25\) and its variance \(185736/15625=11.887104\), rather than the
attenuated-binomial value \(9.34464\).

\begin{proposition}[Exact noisy count recursion]\label{op:count-recursion}
Let \(A_t(h,k)\) be the joint probability of the most recent \(L\) detector
bits \(h\) and \(k\) even observed innovations after \(t\) steps.
Initialize
\[
 A_0(h,0)=\epsilon^{|h|}(1-\epsilon)^{L-|h|}.
\]
For \(b\in\{0,1\}\), put \(r_b=\epsilon^b(1-\epsilon)^{1-b}\),
\(h'=(h_2,\ldots,h_L,b)\), and
\(q(h,b)=[1+\eta(-1)^{|h|+b}]/2\).
For each current mass, add
\[
 A_t(h,k)r_bq(h,b)\ \hbox{to } A_{t+1}(h',k+1),\qquad
 A_t(h,k)r_b[1-q(h,b)]\ \hbox{to } A_{t+1}(h',k).
\]
Then \(P(C_m=k)=\sum_h A_m(h,k)\).
\end{proposition}
\begin{proof}
The next detector bit is independent of the current detector history.
The next source innovation is independent of all preceding innovations
and all detector bits.  Conditional on the detector-window parity, its
observed even probability is exactly \(q(h,b)\).
The displayed update is the law of total probability on these two
new independent inputs.  Induction proves the result.
\end{proof}

For rational parameters this is an integer recurrence after keeping one
common denominator at each step.  Applying Lemma~\ref{op:np} to the count
laws, sorting their actual likelihood ratios rather than presuming their
order, gives the following finite certificate for
\(L=2,\eta=1/2,\epsilon=1/10\) against the fair iid null.
Every \(m=1,\ldots,182\) has optimal count-only power below \(0.95\) at
size \(0.05\); \(m=183\) has power \(0.95047965645\ldots\).
At \(m=182\) it is \(0.94920331752\ldots\).
An easily specified sufficient rule at \(183\) rejects for \(C_m\ge104\)
and with probability \(43/50\) at \(C_m=103\); its size and power are
\(0.04986347138\ldots\) and \(0.95033569239\ldots\).
The exact recurrence and rational comparisons are supplied with the
computational materials.  This minimum is for tests of \(C_m\) only.
Unlike the complete record, \(C_{m+1}\) need not recover \(C_m\);
monotonicity of the count experiments is not used.
Here \(183\) innovations require \(185\) raw outputs.
A nonrandomized sufficient rule uses \(188\) innovations and rejects
for \(C_m\ge106\); its size is \(0.04659560696\ldots\) and power
\(0.95203135112\ldots\).  No nonrandomized minimum is claimed.

The ideal \(40\)-innovation rule has power only
\(0.49258836198\ldots\) under this noisy alternative, while retaining its
null size.  Replacing its count law by a binomial with the attenuated
mean is incorrect.  If instead parity windows start \(L+1\) indices
apart, their defining source innovations and detector windows are
disjoint, so the selected parities are iid.  That thinning consumes
additional raw outputs.

At the population level, known \(L\), \(\eta\ne0\), and \(a>0\) allow
recovery from \(\mu=E Z_t\) and \(\nu=E Z_tZ_{t+1}\):
\[
 a=(\mu^2/\nu)^{1/(2L)},\qquad
 \epsilon=(1-a)/2,\qquad \eta=\mu/a^{L+1}.
\]
Indeed \(\mu^2/\nu=a^{2L}\) and \(\nu=\eta^2a^2\).
The cases \(\eta=0\) and \(a=0\) are degenerate, and the inversion is
poorly conditioned near them.  These identities do not supply a
finite-data confidence interval or validate the iid detector model.

\subsection{Classical realization rank and finite-sample certificates}
\label{op:rank}

Let \(H_L(x,y)=P(Y_{1:L}=x,Y_{L+1:2L}=y)\).
Any \(D\)-state classical hidden realization with a sufficient state at
this boundary factors the matrix as
\[
 H_L(x,y)=\sum_{s=1}^D P(Y_{1:L}=x,S_L=s)
                    P(Y_{L+1:2L}=y\mid S_L=s).
\]
Hence \(\operatorname{rank}H_L\le D\).
Every source, detector, or controller variable that carries predictive
information across the cut belongs in that sufficient state.  General
hidden-model realization theory distinguishes this necessary rank
condition from the existence of a positive realization
\cite{op:ohta,op:vanluyten}.

\begin{theorem}[Parity-block spectrum and exact state count]
\label{op:parity-rank}
In the ideal parity-shift model the singular values of \(H_L\) are
\[
 2^{-L}|\eta|^j,\qquad j=0,\ldots,L,
\]
with multiplicity \(\binom Lj\).
For \(\eta\ne0\), the minimum number of states of an exact classical
transition-emitting realization is \(2^L\).  It is one when \(\eta=0\).
\end{theorem}
\begin{proof}
Define separate bijections of the row and column bit strings by
\(u_t=x_t\oplus\cdots\oplus x_L\) and
\(v_t=y_1\oplus\cdots\oplus y_t\).
The \(t\)th innovation following the cut is \(u_t\oplus v_t\).
The innovation factorization therefore gives, up to separate row and
column permutations,
\[
 H_L=2^{-L}B_\eta^{\otimes L},\qquad
 B_\eta=\frac12\begin{pmatrix}1+\eta&1-\eta\\1-\eta&1+\eta\end{pmatrix}.
\]
The singular values of \(B_\eta\) are \(1,|\eta|\); tensor products give
the asserted list.  Separate row and column permutations preserve
singular values, though they need not be a similarity transformation.
For nonzero \(\eta\) the rank is \(2^L\), providing the lower bound.
The shift-state construction \eqref{op:parity-kernel} has exactly that
many states and provides the matching positive upper realization.
For \(\eta=0\) the outputs are iid fair bits.
\end{proof}

This \(2^L\) result belongs to the parity-shift model.  A stationary
random-phase concatenation of perturbed \(k\)-bit blocks is a different
process, with a different \(2k-1\) realization result.  The parameters
\(L,k\), their record laws, and their state counts are not interchangeable.
Neither count is a bound on physical constituents or on quantum memory.

\begin{theorem}[Readout and the noisy Walsh spectrum]\label{op:noisy-rank}
Independent symmetric readout with \(a=1-2\epsilon>0\) preserves the
rank \(2^L\) for \(\eta\ne0\).
For \(s\in\{0,1\}^L\), define
\[
 r_i(s)=\bigoplus_{j=1}^i s_j,\qquad
 \ell_i(s)=\bigoplus_{j=i}^L s_j .
 \]
The noisy matrix has singular values, indexed by all \(s\),
\[
 \sigma_s=2^{-L}|\eta|^{|s|}
                  a^{|r(s)|+|\ell(s)|}.
\]
In particular \(\sigma_{\min}\ge (|\eta|a^2/2)^L>0\).
\end{theorem}
\begin{proof}
Let \(D_\epsilon=\left(\begin{smallmatrix}1-\epsilon&\epsilon\\
\epsilon&1-\epsilon\end{smallmatrix}\right)\).
Independent noise on both blocks gives
\(\widetilde H_L=D_\epsilon^{\otimes L}H_L
(D_\epsilon^{\otimes L})^T\).
Its two eigenvalues are \(1,a\), so it is invertible for \(a>0\).
For the full spectrum, expand each transition factor to obtain
\[
 H_L(x,y)=2^{-2L}\sum_s\eta^{|s|}
             (-1)^{r(s)\cdot x+\ell(s)\cdot y}.
\]
The maps \(s\mapsto r(s)\), \(s\mapsto\ell(s)\) are bijections.
The normalized Walsh vectors
\(2^{-L/2}(-1)^{v\cdot x}\) form orthonormal bases.
In these bases, \(H_L\) is a weighted permutation matrix with weights
\(2^{-L}\eta^{|s|}\).  Readout multiplies its row and column Walsh
coefficients by \(a^{|r(s)|}\) and \(a^{|\ell(s)|}\).
The absolute weights are the singular values.  Since
\(|s|\le L\) and \(|r(s)|+|\ell(s)|\le2L\), the lower bound follows.
\end{proof}

The same \(2^L\)-state realization with memoryless noisy emission gives
the upper bound after readout.  At \(\epsilon=1/2\) the observed process
is iid fair and the rank collapses.  Exact invertibility does not imply
a well-conditioned statistical inversion.

\begin{lemma}[Singular-tail distance]\label{op:tail}
For a real matrix \(H\) with singular values in decreasing order, define
\(T_D(H)=\sum_{j>D}\sigma_j(H)^2\).
Then
\[
 d_D(H):=\inf_{\operatorname{rank}Q\le D}\|H-Q\|_F
       =\sqrt{T_D(H)},\qquad
 |d_D(H)-d_D(G)|\le\|H-G\|_F.
\]
\end{lemma}
\begin{proof}
If \(Q\) has rank at most \(D\), its nullspace contains a subspace of
dimension equal to the number of columns minus \(D\).  Let \(P\) be
the orthogonal projection onto such a subspace.
Then \(\|H-Q\|_F^2\ge\|(H-Q)P\|_F^2=\|HP\|_F^2\).
In an eigenbasis of \(H^TH\), the latter is a weighted sum of its
eigenvalues with weights in \([0,1]\) summing to the rank of \(P\).
It is at least the sum of the smallest required eigenvalues,
which is \(T_D(H)\).  Keeping the largest \(D\) terms of a singular-value
decomposition attains this lower bound.  This is the
Eckart--Young approximation theorem~\cite{op:eckart}.
Finally, distance to any fixed nonempty subset is \(1\)-Lipschitz by
the triangle inequality; taking infima proves the last statement.
\end{proof}

The rank-\(D\) approximation just used is allowed to be any real matrix.
It need not be nonnegative, normalized, or realizable by a hidden process.

\begin{theorem}[Total variation bounds with marginal information]
\label{op:marginal-rank}
Let \(H,Q\) be normalized nonnegative joint matrices, with
\(\operatorname{rank}Q\le D\).  Write
\[
 \delta=\operatorname{TV}(H,Q),\quad
 r=\operatorname{TV}(H_X,Q_X),\quad
 c=\operatorname{TV}(H_Y,Q_Y),\quad T=T_D(H).
\]
Then
\[
 T\le\|H-Q\|_F^2\le
 \min\{2\delta^2,\delta(\delta+r+c)\},
\]
and therefore
\[
 \delta\ge
 \max\left\{\sqrt{T/2},
 \frac{\sqrt{(r+c)^2+4T}-(r+c)}2\right\}.
\]
When both marginals match, the second bound becomes \(\delta\ge\sqrt T\).
\end{theorem}
\begin{proof}
Put \(A=H-Q\).  Its positive and negative entry masses both equal
\(\delta\).  The sum of squares within each sign is at most the square
of its total mass, proving \(\|A\|_F^2\le2\delta^2\).
For a positive entry \(A_{ij}=x\), the negative mass in row \(i\) is at
least \(x-r\), since the positive part of any row sum is at most \(r\).
The negative mass in column \(j\) is at least \(x-c\).
These sets of negative entries are disjoint, because their only possible
intersection is the positive entry \((i,j)\).  Consequently
\(\delta\ge2x-r-c\).  Applying the same argument to \(-A\) gives
\(\|A\|_{\max}\le(\delta+r+c)/2\).
Thus \(\|A\|_F^2\le\|A\|_{\max}\|A\|_1
\le\delta(\delta+r+c)\).
Lemma~\ref{op:tail} supplies \(T\le\|A\|_F^2\);
solving the two quadratic inequalities proves the result.
\end{proof}

Known upper bounds on \(r+c\) may be substituted, because the second
lower bound decreases with that sum.  For zero row and column sums,
\(\|A\|_F\le\delta\) is sharp, as the signed \(2\)-by-\(2\) cycle
\(\left(\begin{smallmatrix}x&-x\\-x&x\end{smallmatrix}\right)\) shows.
That norm sharpness does not show that a rank-constrained probability
model attains the combined bound.

For \(L=2,\eta=1/2,\epsilon=1/10\), the exact noisy matrix is
\begin{equation}\label{op:noisy-four}
 \widetilde H_2=\frac1{2000}
 \begin{pmatrix}
 209&105&81&105\\
 81&105&209&105\\
 105&81&105&209\\
 105&209&105&81
 \end{pmatrix}.
\end{equation}
Its singular values are \(1/4,8/125,8/125,1/25\), and
\[
 T_1=153/15625,\qquad T_2=89/15625,\qquad T_3=1/625.
\]
Both marginals are uniform on four outcomes.  A normalized rank-one
nonnegative matrix with those same marginals is uniquely the matrix
whose entries are \(1/16\), and its total variation distance is exactly
\(21/125\).  The iid fair-bit process attains this restricted rank-one
competitor.  The number does not apply when competitor marginals may
vary; the bounds for two or three states have not been asserted attainable.

\subsubsection{From population rank to finite samples}

\begin{lemma}[Bounded-differences inequality]\label{op:bounded}
Let \(X_1,\ldots,X_n\) be independent, and suppose changing only \(X_i\)
changes \(f\) by at most \(c_i\).  Then
\[
 P(f-Ef\ge t)\le
 \exp\left(-\frac{2t^2}{\sum_i c_i^2}\right).
\]
\end{lemma}
\begin{proof}
The conditional Doob martingale of \(f\) has each martingale difference
in an interval of length at most \(c_i\): conditional expectations for
two possible values of \(X_i\) differ by at most \(c_i\), after coupling
the still independent remaining coordinates identically.
For a centered variable \(U\) in an interval of length \(c\),
\(\log E e^{\lambda U}\le\lambda^2c^2/8\).
Indeed the second derivative of its log moment-generating function is
the variance under exponential tilting, which is at most \(c^2/4\);
the function and its first derivative vanish at zero, so integration
twice gives the bound.  The variance bound follows from
\(E[(U-a)(b-U)]\ge0\), which gives
\(\operatorname{Var}(U)\le(EU-a)(b-EU)\le(b-a)^2/4\).
Apply the exponential bound conditionally to the martingale differences,
iterate, use Markov's inequality, and minimize over
\(\lambda=4t/\sum c_i^2\).  This is the standard bounded-differences
method~\cite{op:mcdiarmid}.
\end{proof}

\begin{theorem}[Empirical matrix radius and rank test]\label{op:finite-rank}
Fix a matrix shape with \(K\) cells, an observation map, and a sample size
\(n\).  Observe \(n\) independent records having the same categorical
joint law \(H\), and let \(\widehat H\) be their frequency matrix.
Dependence within a record is allowed.  For \(0<\alpha<1\), set
\[
 \varepsilon_\alpha(n,K)=
 \frac{\sqrt{1-1/K}+\sqrt{\log(1/\alpha)}}{\sqrt n}.
\]
Then \(P(\|\widehat H-H\|_F\le\varepsilon_\alpha)\ge1-\alpha\).
The test rejecting rank at most \(D\) when
\(d_D(\widehat H)>\varepsilon_\alpha\) has size at most \(\alpha\).
If \(d_D(H)>\varepsilon_\alpha+\varepsilon_\beta\), its power is at
least \(1-\beta\).
On the same confidence event, simultaneously for every \(D\),
\[
 \inf_{\substack{Q\ge0,\ \sum Q=1\\\operatorname{rank}Q\le D}}
 \operatorname{TV}(H,Q)\ge
 \frac{[d_D(\widehat H)-\varepsilon_\alpha]_+}{\sqrt2}.
\]
\end{theorem}
\begin{proof}
The multinomial second-moment identity gives
\[
 E\|\widehat H-H\|_F^2
 =\frac{1-\|H\|_F^2}{n}\le\frac{1-1/K}{n}.
\]
Jensen's inequality bounds the mean norm by the square root.
Replacing one record changes \(\widehat H\) by Frobenius norm at most
\(\sqrt2/n\).  The reverse triangle inequality gives the same
sensitivity for its error norm.  Lemma~\ref{op:bounded} therefore gives
an upper tail \(\exp(-nt^2)\); use
\(t=\sqrt{\log(1/\alpha)/n}\).
Under the rank null, \(d_D(H)=0\); Lemma~\ref{op:tail} proves the size
bound.  On the radius event with probability at least \(1-\beta\),
\(d_D(\widehat H)\ge d_D(H)-\varepsilon_\beta>\varepsilon_\alpha\),
which proves power.  No additional \(\alpha+\beta\) penalty is needed
for that power statement.  Finally,
\(d_D(H)\ge[d_D(\widehat H)-\varepsilon_\alpha]_+\), and
Theorem~\ref{op:marginal-rank}'s unrestricted bound applies to every
normalized rank-\(D\) competitor.
\end{proof}

At the single predeclared matrix and horizon, the confidence event
simultaneously controls all \(D\); repeated looks at changing horizons
are a different experiment.  Complete categorical outcomes, including
any failure outcome required by the experiment, must be represented.
A numerically computed singular tail requires an error allowance if
rounding can affect the threshold.

For \eqref{op:noisy-four}, \(K=16\) and \(\alpha=\beta=1/20\).
Let \(C=\sqrt{15/16}+\sqrt{\log20}\).
The condition is \(n>4C^2/T_D\).  Its first integer solutions are
\[
 \begin{array}{c|ccc}
 D&1&2&3\\ \hline
 n&2976&5116&18213\\
 \text{raw outputs }4n&11904&20464&72852
 \end{array}
\]
These are sufficient numbers of iid four-output records for this
concentration bound, not optimal testing budgets.  At \(n=18213\) the
single common power event suffices for all three rank rejections.
Experiments separately stopped at the three listed horizons do not
automatically constitute one test of size \(0.05\).

\subsection{Sequential evidence for complete record laws}
\label{op:sequential}

The null in this subsection is the complete iid fair-bit law,
\(P_0(Y_{1:n}=y)=2^{-n}\).
Fair one-bit marginals alone do not define this null.
Fix coherent candidate laws \(P_\theta\) before seeing data and
nonnegative mixing weights totaling one.

\begin{theorem}[Likelihood mixtures and an anytime threshold]
\label{op:mixture}
Each \(R_{\theta,n}=2^nP_\theta(Y_{1:n})\) is a nonnegative
\(P_0\)-martingale starting at one.  Their fixed mixture
\(\mathcal E_n=\int R_{\theta,n}\,dW(\theta)\) is also such a martingale,
and
\[
 P_0\{\exists n\ge0:\mathcal E_n\ge1/\alpha\}\le\alpha.
\]
\end{theorem}
\begin{proof}
Conditional on a prefix under \(P_0\), the next bit is fair.
Coherence gives
\[
 E_0[R_{\theta,n+1}\mid Y_{1:n}]
 =2^n\sum_{b=0}^1P_\theta(Y_{1:n},b)=R_{\theta,n}.
\]
Nonnegative integration permits conditional Tonelli, proving the mixture
claim.  Stop at the first threshold crossing or at a fixed \(N\).
The expected stopped value is one by finite optional stopping, so
\(P_0(\text{crossing by }N)\le\alpha\).
Let \(N\) increase.  The same argument holds for nonnegative
supermartingales and is Ville's inequality; see
Howard et al.~\cite{op:howard}.
\end{proof}

For the noisy parity family with fixed \(\eta=1/2,\epsilon=1/10\)
and unknown \(L\), take \(w_L=1/[L(L+1)]\).
When \(L\ge n\), every \(n\)-bit observed law is uniform.
Since \(\sum_{L=n}^\infty w_L=1/n\), the mixture has the finite expression
\begin{equation}\label{op:length-mixture}
 \mathcal E_n=\frac1n+
 \sum_{L=1}^{n-1}\frac{2^nP_L(Y_{1:n})}{L(L+1)},\qquad n\ge1,
 \quad \mathcal E_0=1.
\end{equation}
This identity does not guarantee inexpensive evaluation: a component
may require \(2^L\) hidden states.  Replacing the fixed mixture by a
post-data maximum does not preserve its martingale proof.

\begin{proposition}[A sufficient \(1950\)-output horizon]\label{op:1950}
At the fixed true alternative \(L=2,\eta=1/2,\epsilon=1/10\), the test
crossing \(20\) in \eqref{op:length-mixture} has power by \(1950\)
raw outputs at least \(0.9519745696463602\ldots\).
\end{proposition}
\begin{proof}
The \(650\) disjoint three-bit parity windows have independent
even-parity probability \(p=157/250\).  Let \(K\) be their even count and
let
\[
 S=(157/125)^K(93/125)^{650-K}
\]
be their likelihood ratio against fair parity.
For the full-record component \(R=2^{1950}P_2(Y_{1:1950})\),
strict positivity gives
\(E_2(S/R)=E_0S=1\).
The mixture contains \(R/6\).  Failure to cross by this horizon implies
in particular \(R<120\).
For any \(c>0\),
\[
 P_2(R<120)\le P_2(S<120c)+P_2(S/R>c)
              \le P_2(S<120c)+1/c .
\]
With \(c=40\), exact integer comparisons show
\(S\ge4800\) if and only if \(K\ge384\).
Therefore the miss probability is at most
\[
 \sum_{k=0}^{383}\binom{650}{k}
       (157/250)^k(93/250)^{650-k}+\frac1{40}
 =0.04802543035363984\ldots<\frac1{20}.
\]
All quantities in the last inequality are rational.  The calculation
does not require evaluating all nonconstant mixture components.
\end{proof}

This is a sufficient horizon for one specified alternative, not a
minimum, an expected stopping time, or a bound uniform in \(L\).
It is a different experiment from the fixed-horizon \(183\)-innovation
count test.

\begin{theorem}[Pointwise eventual detection]\label{op:eventual}
Fix a nondegenerate parity candidate with \(L<\infty\), \(\eta\ne0\),
\(0\le\epsilon<1/2\), and positive mixture weight \(w_L\).
If that candidate generates the data, its likelihood mixture eventually
crosses every fixed finite threshold almost surely.  More precisely,
\[
 \liminf_{n\to\infty}\frac{\log\mathcal E_n}{n}
 \ge \frac{d_{\rm Ber}(p_L\Vert1/2)}{L+1}>0,\qquad
 p_L=\frac{1+\eta(1-2\epsilon)^{L+1}}2 .
\]
\end{theorem}
\begin{proof}
Use \(m=\lfloor n/(L+1)\rfloor\) disjoint parity windows, and let \(S_m\)
be their likelihood ratio.  The iid strong law gives
\(\log S_m/n\to d_{\rm Ber}(p_L\Vert1/2)/(L+1)\).
For the full-record ratio \(R_{L,n}\), the same change-of-measure
identity as above gives \(E_L[S_m/R_{L,n}]=1\).
For every \(b>0\),
\[
 P_L(R_{L,n}<S_me^{-bn})\le e^{-bn}.
\]
The sum over \(n\) is finite, so the first Borel--Cantelli lemma
implies that this event occurs only finitely often.
Apply the assertion to a countable sequence of positive rational
\(b\) decreasing to zero and use
\(\mathcal E_n\ge w_LR_{L,n}\).
The fixed multiplicative weight has vanishing logarithm per \(n\).
\end{proof}

No uniform finite horizon follows: for \(L\ge n\) the full length-\(n\)
law equals the null exactly.  Alternatives approaching \(\eta=0\) or
\(\epsilon=1/2\) also approach the null on each fixed horizon.
For a concrete warning about the null, the process \(Y_t=B\) for one
fair bit \(B\) has fair marginals.  In the ideal-readout
\(\eta=1/2\) mixture, its \(L=1\) contribution is
\(\frac12(3/2)^{n-1}\), exceeding \(20\) at \(n=11\).
It is not an iid-null process; rejecting it does not violate the theorem.

\subsubsection{Conditional null envelopes and useful growth}

Suppose instead that the null supplies a predictable envelope
\[
 p_t=P(Y_t=1\mid\mathcal F_{t-1})
 \in[1/2-\delta_t,1/2+\delta_t],\quad 0\le\delta_t\le1/2 .
\]
For a predictable Bernoulli forecast \(q_t\), put \(q_t(1)=q_t\) and
\(q_t(0)=1-q_t\).  The factor
\[
 F_t=\frac{2q_t(Y_t)}
           {1+2\delta_t|2q_t-1|}
\]
has conditional null expectation at most one: its numerator has
expectation \(1+2(p_t-1/2)(2q_t-1)\), maximized exactly at the appropriate
endpoint of the interval.  Its products obey the crossing guarantee of
Theorem~\ref{op:mixture}.  This sharp normalization does not ensure
favorable evidence growth.

\begin{theorem}[Log-optimal binary factor for an interval null]
\label{op:clipped}
For \(0<\ell\le u<1\) and an alternative \(0<q<1\), let
\(p_*=\min\{u,\max\{\ell,q\}\}\).  Define
\[
 G_q(1)=q/p_*,\qquad G_q(0)=(1-q)/(1-p_*).
\]
It has expectation at most one for every Bernoulli parameter in
\([\ell,u]\), and maximizes \(E_q\log g(Y)\) over all nonnegative
factors \(g\) valid throughout that interval.
\end{theorem}
\begin{proof}
If \(q>u\), the expectation of \(G_q\) increases with the actual
Bernoulli parameter and equals one at \(u=p_*\).
If \(q<\ell\), it decreases and equals one at \(\ell=p_*\).
If \(q\in[\ell,u]\), it is identically one.
For every valid \(g\), Jensen's inequality gives
\[
 E_q\log\frac{g(Y)}{G_q(Y)}
 \le\log E_q\frac{g(Y)}{G_q(Y)}
 =\log E_{p_*}g(Y)\le0.
\]
Zero factors are handled by the extended value \(-\infty\) for the log.
\end{proof}

This is a binary specialization of the established growth-optimal
e-variable framework~\cite{op:larsson}.  Predictable choices of
\(\ell_t,u_t,q_t\) again yield a nonnegative null supermartingale.
The criterion is expected one-step log growth, not finite-horizon
Neyman--Pearson power.

For \(q=13/20\) and null interval \([2/5,3/5]\), the two factors are
\[
 (F(0),F(1))=(35/53,65/53),\qquad
 (G(0),G(1))=(7/8,13/12).
\]
Their expected log growth under iid Bernoulli\((13/20)\) is negative
and positive, respectively.  The exact sign certificates are
\[
 65^{13}35^7<53^{20},\qquad 13^{13}7^7>12^{13}8^7.
\]
The strong law makes the first product tend to zero and the second
diverge almost surely.  The first can still cross a threshold early.

A conditional-bias envelope is not a memorylessness null.  In the noisy
parity model, conditioning first on the full source history and then
averaging over the observed past yields
\[
 \left|P(\widetilde Y_t=1\mid\widetilde Y_{<t})-\frac12\right|
 \le\frac{|\eta|(1-2\epsilon)}2 .
\]
For the numerical alternative this is \(1/5\).
A null envelope at least that wide contains the dependent process, so
a test valid for every law in it must retain its null error guarantee
on that process.  The envelope itself needs justification for the
actual observation filtration; marginal calibration does not supply it.

\subsection{Visibility, environmental records, and access to fragments}
\label{op:records}

Consider an ordinary qubit and an environment with a fixed tensor
decomposition into \(N\) fragments.  A stipulated controlled interaction
maps the two qubit basis states to branch records
\[
 |0\rangle\longmapsto |0\rangle|e_0\rangle^{\otimes N},\qquad
 |1\rangle\longmapsto |1\rangle|e_1\rangle^{\otimes N},
 \quad \langle e_0|e_1\rangle=\gamma,\quad |\gamma|\le1.
\]
For a balanced input, tracing out the fragments leaves coherence of
magnitude \(|\gamma|^N/2\).  A phase setting with the convention
\(p_0(\phi)=\frac12[1+\operatorname{Re}(e^{i\phi}\gamma^N)]\)
has visibility \(|\gamma|^N\) and qubit purity
\((1+|\gamma|^{2N})/2\).
These formulas assume the stated state preparation, tensor partition,
controlled interaction, and Born-rule measurement.  They derive no
macroscopic amplification mechanism.

\begin{proposition}[A classical realization of the binary visibility law]
\label{op:visibility-alias}
For real \(0\le\gamma\le1\), all final binary laws
\[
 p_0(N,\phi)=\tfrac12(1+\gamma^N\cos\phi),
 \qquad N\in\{0,1,\ldots\},\ \phi\in\mathbb R,
\]
are reproduced by a two-state classical Markov chain and a
setting-dependent binary response.
\end{proposition}
\begin{proof}
Let \(Z_0=+1\).  Flip the sign independently at each step with
probability \((1-\gamma)/2\), so that the transition matrix is
\[
 K_\gamma=\frac12
 \begin{pmatrix}1+\gamma&1-\gamma\\1-\gamma&1+\gamma\end{pmatrix}.
\]
Then \(E Z_N=\gamma^N\), by the one-step conditional expectation
\(E(Z_{n+1}\mid Z_n)=\gamma Z_n\).
Given \(Z_N=z\) and the calibrated setting \(\phi\), report outcome zero
with probability \((1+z\cos\phi)/2\).  Averaging proves the formula.
At \(N=0\) the chain has made no transition, including when \(\gamma=0\).
\end{proof}

For complex \(\gamma=re^{i\beta}\), one may instead choose independent
classical phases \(\Theta_j=\beta\) or \(\beta+\pi\) with probabilities
\((1+r)/2,(1-r)/2\).  Since \(E e^{i\Theta_j}=\gamma\), the response
\([1+\cos(\phi+\sum_j\Theta_j)]/2\) has the required averaged law.
This construction uses accumulated phase; it is not the same assertion
of a homogeneous two-state realization with one fixed response for all
complex \(\gamma\).
Related random-unitary descriptions of quantum noise are established
\cite{op:crow}; their scope differs from this explicitly classical
simulation of a specified final binary family.  Neither construction
simulates arbitrary quantum preparations, joint measurements, Bell
experiments, or every environmental observable.

\begin{theorem}[A final channel does not identify fragment factorization]
\label{op:environment}
Fix the physical fragment partition and let \(0<\gamma<1\),
\(s=\sqrt{1-\gamma^2}\).  For \(N\ge2\), the product branches
\[
 E_0=|0\rangle^{\otimes N},\qquad
 E_1=(\gamma|0\rangle+s|1\rangle)^{\otimes N}
\]
and the branches \(F_i=U_NE_i\), where
\(U_N=\operatorname{CZ}_{12}\otimes I_{3:\!N}\), induce exactly the
same complete final qubit channel, including its action with a retained
ancilla.  Nevertheless \(F_1\) is not a product across the first two
fragments.
If the transformed old fragments are never accessed or recoupled, the
two descriptions also give the same complete histories for common
operations on the retained system, ancillas, and fresh fragments.
\end{theorem}
\begin{proof}
For any joint state \(\rho\) of a retained system and environment,
\[
 \operatorname{Tr}_E[(I\otimes U_N)\rho(I\otimes U_N^\dagger)]
 =\operatorname{Tr}_E\rho .
\]
This identity is valid with an arbitrary retained ancilla, proving the
complete-channel statement.  It also preserves the branch overlap
\(\langle E_0|E_1\rangle=\gamma^N\).
For the first two fragments, the coefficient matrix of \(F_1\) is
\[
 \begin{pmatrix}\gamma^2&\gamma s\\
                 \gamma s&-s^2\end{pmatrix},
 \qquad \det=-2\gamma^2s^2\ne0 .
\]
A pure two-fragment state is a product precisely when its coefficient
matrix has rank one; thus this branch is not a product.
Appending the remaining product factors does not remove that fact.
For a later history satisfying the stated access restriction, the
common environmental unitary commutes with every later retained-system
or fresh-fragment operation.  Move it through the product of the
corresponding instrument maps and remove it under the final trace.
This proves equality for each fixed outcome/action history; common
history-dependent randomized policies preserve it by conditioning.
\end{proof}

The history assertion requires absence of later coupling to the
transformed fragments.  Equality of final reduced channels by itself
does not prove equality under arbitrary later access.  Also, the
distinguished property here is tensor-product factorization of
conditional quantum fragment states.  It is not a claim that every
chosen classical fragment readout becomes statistically dependent.
For example, the common controlled phase leaves the branch's
\(Z\)-basis probability distribution unchanged.  If arbitrary entangling
changes of tensor factorization are declared to be a gauge freedom,
the target factorization property itself changes.
At \(\gamma=0\) or \(1\), or with fewer than two fragments, the stated
nonproduct witness is absent; common-unitary invariance of the reduced
channel still holds.

\begin{example}[A calibrated fragment measurement separates the pair]
\label{op:fragment-readout}
For \(\gamma=3/5,s=4/5,N=2\), the branch-one vectors are
\[
 E_1=(9,12,12,16)^T/25,\qquad
 F_1=(9,12,12,-16)^T/25.
\]
Both are normalized and have overlap \(9/25\) with \(|00\rangle\).
Their coefficient determinants are zero and \(-288/625\).
The first-fragment reduced off-diagonal entries are \(12/25\) and
\(-84/625\), so a local \(X\) measurement has
\[
 P_{E_1}(X=+)=49/50,\qquad P_{F_1}(X=+)=457/1250,
 \qquad \text{gap}=384/625.
\]
This requires anchored preparation of branch one, access to the specified
fragment, and a calibrated \(X\) readout.  If each reported binary law
is within total variation \(\varepsilon\) of its stated ideal law,
the gap is at least \(384/625-2\varepsilon\); strict positivity is a
sufficient separation condition.  This is pair separation, not
identification of an arbitrary environmental model.
\end{example}

\subsection{Positive hidden realizations and exact intervention contracts}
\label{op:controls}

A finite transition-emitting, or Mealy, model has nonnegative matrices
\(T_y\), indexed by a finite output alphabet, whose sum \(P=\sum_y T_y\)
is row-stochastic.  From an initial row distribution \(\pi\),
\[
 P(Y_{1:m}=y_1\cdots y_m)=\pi T_{y_1}\cdots T_{y_m}\mathbf1.
\]
A silent control is a row-stochastic matrix inserted between output
steps, producing no extra recorded outcome.  Linear reachability means
that the rows \(\{\pi T_w\}\) span the whole row space; linear
observability means that \(\{T_w\mathbf1\}\) spans the whole column
space, where \(w\) includes the empty word.  We use \emph{linearly minimal}
when both properties hold.  It is stronger than merely displaying no
obvious redundant state.

If \(S\) is invertible with \(S\mathbf1=\mathbf1\), put
\[
 T'_y=S^{-1}T_yS,\qquad \pi'=\pi S .
\]
Whenever these transformed quantities are valid probabilities, every
passive word law is unchanged by telescoping.  If each common fixed
control commutes with \(S\), all controlled word laws are unchanged as
well.  If instead controls are themselves allowed to vary, simultaneous
transformation \(K'_j=S^{-1}K_jS\) gives the same telescoping identity.
The distinction between exact fixed controls and uncertain controls
is substantive.  Similarity equivalence for Mealy models and its
distinction from other hidden-model classes are established
\cite{op:vanluyten}.

\begin{example}[A strictly positive minimal passive pair]
\label{op:mealy-pair}
Take
\[
 T_0=\frac1{20}\begin{pmatrix}6&2\\1&3\end{pmatrix},\quad
 T_1=\frac1{20}\begin{pmatrix}4&8\\7&9\end{pmatrix},\quad
 \pi=(4/9,5/9),\quad
 S=\frac1{10}\begin{pmatrix}9&1\\2&8\end{pmatrix}.
\]
The transformed matrices are
\[
 T'_0=\frac1{1400}\begin{pmatrix}449&151\\19&181\end{pmatrix},
 \qquad
 T'_1=\frac1{280}\begin{pmatrix}67&93\\125&115\end{pmatrix},
 \qquad \pi'=\pi S .
\]
Both representations are strictly positive and stationary:
\(\pi P=\pi\) and \(\pi'P'=\pi'\).
Their common two-output matrix is
\[
 J=\frac1{900}\begin{pmatrix}81&179\\179&461\end{pmatrix},
 \qquad \det J=53/8100.
\]
The factorization of \(J\) through a hidden boundary state implies
minimal dimension at least two; the displayed models attain it.
Their total-transition diagonal multisets are
\(\{1/2,3/5\}\) and \(\{14/25,27/50\}\), respectively.
They therefore are not related by a permutation of hidden-state labels.
\end{example}

\begin{proposition}[Exactly which fixed controls separate that pair]
\label{op:pair-control}
Let the same anchored control
\[
 K(a,b)=\begin{pmatrix}1-a&a\\b&1-b\end{pmatrix},
 \qquad 0\le a,b\le1,
\]
act in both models of Example~\ref{op:mealy-pair}.
All controlled output-history laws agree if and only if \(b=2a\).
The next-zero contrast after one such control is
\[
 \pi KT_0\mathbf1-\pi'KT'_0\mathbf1=(2a-b)/35.
\]
\end{proposition}
\begin{proof}
Direct multiplication gives the contrast and
\[
 SK-KS=\frac1{10}
 \begin{pmatrix}b-2a&2a-b\\b-2a&2a-b\end{pmatrix}.
\]
If \(b=2a\), the similarity telescopes through every interleaving
of outputs and controls.  If \(b\ne2a\), the displayed single-output
experiment already distinguishes the pair.  Common randomized or
adaptive control policies preserve the equality in the first case by
conditioning on their observed histories.
\end{proof}

The reset \(R_2=\left(\begin{smallmatrix}0&1\\0&1\end{smallmatrix}\right)\)
maximizes the absolute single-output contrast over this control square:
it gives \(2/35\), with zero probabilities \(1/5\) and \(1/7\).
This is an optimum for that contrast and that fixed pair.
It does not identify the entire model class.  For example,
\[
 U=\begin{pmatrix}9/10&1/10\\0&1\end{pmatrix}
 \quad\hbox{satisfies}\quad UR_2=R_2U=R_2 .
\]
The further transformed model has
\[
 \pi U=(2/5,3/5),\quad
 U^{-1}T_0U=\frac1{1800}\begin{pmatrix}531&229\\81&279\end{pmatrix},
 \quad
 U^{-1}T_1U=\frac1{1800}\begin{pmatrix}297&743\\567&873\end{pmatrix}.
\]
It is strictly positive, remains minimal, and has total-transition
diagonal multiset \(\{23/50,16/25\}\).
It is a further nonpermutation partner agreeing with the original on
all passive and exact-\(R_2\) controlled laws.

\begin{theorem}[One silent control does not identify the unrestricted
positive minimal class]\label{op:single-control}
Let \(n\ge2\), let \((\pi,\{T_y\})\) be a strictly positive linearly
minimal finite-alphabet realization, and assume its unanchored
preparation and output matrices may vary in a neighborhood of their
current values within the model class.  A single exact common
row-stochastic silent control \(K\) leaves a continuum of distinct
positive realizations with the same complete controlled laws.
Stationarity, if imposed on \(\pi\), is preserved.
\end{theorem}
\begin{proof}
For \(K\ne I\), choose \(S_t=I+t(K-I)\).
It satisfies \(S_t\mathbf1=\mathbf1\), commutes with \(K\), and is
invertible for sufficiently small \(t\).
At \(t=0\), all entries of \(T_y\) and \(\pi\) are positive.
Continuity and finiteness preserve positivity for a sufficiently small
interval of \(t\).  Normalization follows from \(S_t\mathbf1=\mathbf1\).
The transformed total transition matrix and preparation satisfy
\(\pi'P'=\pi P S_t=\pi S_t=\pi'\) when the original is stationary.
The telescoping identity proves equality for all controlled words.
If \(K=I\), use any nonzero matrix \(A\) with \(A\mathbf1=0\) and
\(S_t=I+tA\), again sufficiently near the identity.
To see that distinct similarities give distinct numerical realizations,
suppose a similarity \(V\) fixes both \(\pi\) and every \(T_y\).
Then \(\pi T_wV=\pi T_w\) for every word.  Reachability forces \(V=I\).
The ratio of two similarities yielding identical realizations would
be such a \(V\), hence the similarities would be equal.
A continuum of distinct realizations cannot be removed by the finite
set of state permutations.
\end{proof}

The theorem concerns a local open positive class, an unanchored
preparation, and silent stochastic controls.  Structural zeros, fixed
preparations, recorded control outcomes, or narrower parametrizations
can change the conclusion.  It is not a claim about all hidden-model
classes.  Adding invisible hidden labels is an easier ambiguity:
one may append arbitrary hidden transitions whose pushed visible rows
stay fixed, as in Theorem~\ref{op:lumpability}.  Such redundant lifts
are not counterexamples to minimal-realization theorems and do not
establish distinct physical ontologies.

\begin{theorem}[A sharp two-control criterion in two states]
\label{op:two-controls}
Consider stationary strictly positive binary Mealy models with two
states, passive linear dimension two, and an unanchored initial
distribution.  Suppose two labelled controls are known exactly and
have the same matrices in every admissible model:
\[
 K_j=\begin{pmatrix}1-a_j&a_j\\b_j&1-b_j\end{pmatrix},\quad j=1,2 .
\]
Complete controlled population laws identify every numerical output
matrix and the labelled preparation in this class if and only if
\[
 \Delta=a_1b_2-a_2b_1\ne0 .
\]
When \(\Delta\ne0\), experiments with at most one control followed by
at most two outputs suffice for the reconstruction.
\end{theorem}
\begin{proof}
Write \(\pi=(r,1-r)\), \(e=T_0\mathbf1=(e_1,e_2)^T\),
\(d=e_1-e_2\), \(k=rd\), \(q=P(0)\), and
\(q_j=P(K_j\text{ then }0)\).
If \(d=0\), both \(T_0\mathbf1=q\mathbf1\) and
\(T_1\mathbf1=(1-q)\mathbf1\).  Repeated multiplication makes every
word iid Bernoulli, contradicting passive dimension two.  Thus \(d\ne0\).
The first-state mass after \(K_j\) is
\(r+b_j-(a_j+b_j)r\), hence
\[
 \begin{pmatrix}q_1-q\\q_2-q\end{pmatrix}
 =
 B\begin{pmatrix}d\\k\end{pmatrix},\qquad
 B=\begin{pmatrix}
 b_1&-(a_1+b_1)\\ b_2&-(a_2+b_2)
 \end{pmatrix},\quad \det B=\Delta .
\]
If \(\Delta\ne0\), solve for \(d,k\), then recover
\[
 r=k/d,\qquad e_2=q-k,\qquad e_1=e_2+d .
\]
At least one \(q_j-q\) is nonzero, because otherwise \(B\) would map
the nonzero vector \((d,k)^T\) to zero.  For such \(j\), define
\[
 A=\begin{pmatrix}\pi\\\pi K_j\end{pmatrix},\quad
 F=\begin{pmatrix}\mathbf1&e\end{pmatrix},\quad
 C_y=\begin{pmatrix}
 P(y)&P(y0)\\
 P(K_j\text{ then }y)&P(K_j\text{ then }y0)
 \end{pmatrix}.
\]
Then \(C_y=AT_yF\),
\(\det A=-(q_j-q)/d\), and \(\det F=-d\).
Both are invertible, so \(T_y=A^{-1}C_yF^{-1}\).
All quantities are obtained from the stated short experiments.

Conversely, direct multiplication gives
\[
 [K_1,K_2]=\Delta\begin{pmatrix}1&-1\\1&-1\end{pmatrix}.
\]
If \(\Delta=0\) and a control \(K\ne I\) exists, the nonconstant
similarities \(I+t(K-I)\) commute with both controls.
If both controls are identity, use any nontrivial near-identity
row-normalized similarity.  The positivity, stationarity, and
minimality argument of Theorem~\ref{op:single-control} provides a
continuum of admissible distinct realizations with all controlled
laws equal.  Identification is impossible.
\end{proof}

For a positive non-reset example,
\[
 K_1=\begin{pmatrix}3/4&1/4\\1/2&1/2\end{pmatrix},\qquad
 K_2=\begin{pmatrix}1/2&1/2\\1/4&3/4\end{pmatrix},\qquad
 \Delta=-3/16.
\]
For Example~\ref{op:mealy-pair},
\[
 q=13/45,\quad q_1=29/90,\quad q_2=49/180,\quad
 d=1/5,\quad k=4/45,
\]
and the short-word matrices for \(j=1\) are
\[
 C_0=\begin{pmatrix}13/45&9/100\\29/90&21/200\end{pmatrix},
 \qquad
 C_1=\begin{pmatrix}32/45&179/900\\61/90&337/1800\end{pmatrix}.
\]
The formulas recover the displayed \(\pi,T_0,T_1\) exactly.
Small \(|\Delta|\), small \(|d|\), or a small preparation contrast can
make the inverse numerically unstable; this population theorem
contains no finite-sample stability bound.

\begin{proposition}[Pure preparations and a finite-state positive control]
\label{op:anchored-preparations}
For a two-state model with two exactly anchored labelled pure
preparations, put
\(e_i=P_i(0)\), \(p_i(y)=P_i(y)\), \(z_i(y)=P_i(y0)\).
If \(d=e_1-e_2\ne0\), then
\[
 (T_y)_{i1}=\frac{z_i(y)-e_2p_i(y)}d,\qquad
 (T_y)_{i2}=\frac{e_1p_i(y)-z_i(y)}d.
\]
More generally, an observable \(n\)-state realization is identified by
an exact labelled cyclic permutation control and an exact reset to
one labelled state, using experiments with at most \(n\) output symbols.
\end{proposition}
\begin{proof}
For two states, the equations
\((T_y)_{i1}+(T_y)_{i2}=p_i(y)\) and
\(e_1(T_y)_{i1}+e_2(T_y)_{i2}=z_i(y)\)
give the formulas.
For \(n\) states, a reset followed by powers of the cycle prepares
every pure labelled state.  Let
\(V_k=\operatorname{span}\{T_v\mathbf1:|v|\le k\}\).
The initial space \(V_0\) has dimension one.
If \(V_k=V_{k+1}\), each \(T_y\) maps \(V_k\) into itself, so all later
spaces are identical.  Observability therefore implies that full
dimension is reached by \(k=n-1\).
Choose \(n\) spanning suffix columns to form an invertible matrix \(F\).
Their probabilities under the pure preparations give every entry of
\(F\).  Probabilities of \(yv\) under the same preparations give
\(T_yF\), so recover \(T_y\).  Probabilities of those suffixes from the
original preparation give \(\pi F\), so recover \(\pi\).
\end{proof}

The cycle and reset contain structural zeros and permit repeated uses.
The bound concerns the number of output symbols per setting; it does
not count control uses, distinct settings, repeated trials, duration,
or calibration cost.  Together with the one-control obstruction, two
types are minimal only for the stated universal ideal control-design
question over positive minimal models.  Restricted model classes may
require fewer.

\subsection{Calibration intervals and statistical separation}
\label{op:calibration}

Exact controls and approximately known controls define different
inverse problems.  This can be shown quantitatively without assigning
a prior to the uncertain control.

\begin{theorem}[An explicit calibration-induced identification interval]
\label{op:calibration-interval}
Use the positive model of Example~\ref{op:mealy-pair} and the two
controls displayed after Theorem~\ref{op:two-controls}.
Keep \(K_1\) fixed exactly, but allow \(K_2\) to lie in its full
maximum-row total-variation ball of radius \(\varepsilon\).
Keep the preparation unanchored.  Among strictly positive stationary
models with the same complete controlled output laws, the first
preparation coordinate includes the interval
\[
 \left[\frac49,\frac49+\min\left\{\frac1{15},
                                      \frac{8\varepsilon}{9}\right\}\right].
\]
This is a contained interval, not the complete identified set.
\end{theorem}
\begin{proof}
For \(0\le t\le1/10\), use
\[
 S_t=\begin{pmatrix}1-t&t\\2t&1-2t\end{pmatrix},
 \qquad \det S_t=1-3t\ge7/10.
\]
Direct multiplication gives
\[
 S_t^{-1}K_1S_t=K_1,\qquad
 S_t^{-1}K_2S_t-K_2
 =\begin{pmatrix}3t/4&-3t/4\\3t/4&-3t/4\end{pmatrix}.
\]
Thus the maximum-row total variation error is \(3t/4\).
The transformed preparation is
\(\pi S_t=((4+6t)/9,(5-6t)/9)\).
For completeness, the eight numerators of the two transformed output
matrices, over the positive common denominator \(20(1-3t)\), are
\[
 \begin{array}{cc}
 6-15t-t^2&2-5t+t^2\\
 1-8t-t^2&3-12t+t^2\\[2pt]
 4-3t-35t^2&8-37t+35t^2\\
 7-4t-35t^2&9-36t+35t^2
 \end{array}.
\]
On the stated interval they are bounded below, in the same order, by
\[
 449/100,\ 3/2,\ 19/100,\ 9/5,\
 67/20,\ 43/10,\ 25/4,\ 27/5,
\]
all strictly positive.  Normalization and stationarity are preserved by
similarity.  Simultaneously transforming \(T_y,\pi,K_2\), while \(K_1\)
stays fixed, makes every controlled word telescope to its original
probability.  The admissible segment is
\(0\le t\le\min\{1/10,4\varepsilon/3\}\).
Its first preparation coordinate is \(4/9+2t/3\), giving the interval.
\end{proof}

For example, \(\varepsilon=1/100\) gives
\([4/9,34/75]\), of width \(2/225\).
This is uncertainty in a chosen latent coordinate, not uncertainty in
the common predictions.  An independently fixed preparation or a
model for extra calibration outcomes can break the equivalence.

\begin{proposition}[Finite positive menus with open calibration classes]
\label{op:open-calibration}
Let \(n\ge2\), with finitely many strictly positive output matrices,
a positive unanchored preparation, and finitely many strictly positive
silent controls.  Suppose their admissible classes contain local
neighborhoods of the output/preparation parameters and relative-open
neighborhoods of each control within row-stochastic matrices.
If the realization is linearly minimal, there is a local continuum of
distinct admissible simultaneous-similarity realizations with the same
complete controlled laws.
\end{proposition}
\begin{proof}
Choose a nonzero matrix \(A\) with \(A\mathbf1=0\), and
\(S_t=I+tA\) sufficiently near \(I\).
Transform every output matrix and every control by \(S_t^{-1}(\cdot)S_t\),
and the preparation by \(\pi S_t\).
The sum of the transformed output matrices and each transformed control
are row-stochastic, and the transformed preparation sums to one, because
\(S_t\mathbf1=\mathbf1\).
Strict positivity and the finitely many open-neighborhood conditions
persist for sufficiently small \(t\).
All probabilities telescope.  The trivial stabilizer argument in
Theorem~\ref{op:single-control} proves distinctness.
\end{proof}

The open-neighborhood premise excludes exact structural zeros and other
constraints that obstruct the path.  This qualitative result provides
no universal numerical interval width.

\subsubsection{The reset experiment with bounded nuisance}

For the fixed pair in Example~\ref{op:mealy-pair}, suppose the intended
reset to state two may place a candidate-specific mass \(u\in[0,\delta]\)
on state one.  The next-zero laws range over
\[
 p_A=(1+u)/5,\qquad p_B=(1+2u)/7.
\]
For \(0\le\delta<1/5\), the two probability intervals have strictly
positive separation
\[
 g_\delta=\frac15-\frac{1+2\delta}7
          =\frac{2-10\delta}{35}.
\]
At \(\delta=1/5\), take the exact \(R_2\) in the first model and
\[
 S^{-1}R_2S=
 \begin{pmatrix}1/5&4/5\\1/5&4/5\end{pmatrix}
\]
in the second.  The complete controlled laws, not only the next output,
then agree by simultaneous similarity.  The same witness is admissible
for every larger \(\delta\).  Thus this fixed-pair nuisance class has
a sharp \(20\%\) threshold for uniform separation by arbitrary records.
With an additional total-variation readout allowance \(\epsilon\) on
each binary law, \(g_\delta-2\epsilon>0\) is a sufficient one-output
separation condition.  The full-history sharpness claim above concerns
the explicitly specified reset nuisance, not every possible
history-dependent readout-error class.

With perfect resets, perfect readout, and independent repeated reset
trials, the simple hypotheses are Bernoulli\((1/7)\) versus
Bernoulli\((1/5)\) for the zero indicator.
Lemma~\ref{op:np} yields optimal size-\(1/20\) randomized powers
\[
 \begin{array}{c|cc}
 n&467&468\\ \hline
 \text{power}&0.9498952948087144\ldots&
               0.9505519541817742\ldots
 \end{array}
\]
At \(n=468\), reject for at least \(81\) zeros and randomize at \(80\)
with probability \(0.998625929855047\ldots\), defined exactly by
\[
 \gamma=\frac{1/20-
   \sum_{k=81}^{468}\binom{468}{k}(1/7)^k(6/7)^{468-k}}
   {\binom{468}{80}(1/7)^{80}(6/7)^{388}}.
\]
Discarding an extra iid trial cannot reduce attainable full-record
power, so the \(467/468\) boundary proves the minimum for this specified
randomized fixed-horizon test class.  It does not include calibration
trials or apply to the two-control reconstruction theorem.

For a different, robust experiment suppose the conditional zero
probability after every observed past is at most \(p_-\) under one
class and at least \(p_+\) under the other, with \(g=p_+-p_->0\).
Threshold the mean at \((p_-+p_+)/2\).
The conditional exponential bound used in
Lemma~\ref{op:bounded}, applied to the bounded martingale differences
of the zero indicators, gives each error at most
\(\exp(-ng^2/2)\).
It suffices that \(n\ge\lceil2\log(20)/g^2\rceil\).
For \(g=2/35\), this is \(1835\) trials.
Independence is unnecessary for this conditional-mean argument, but
the interval must hold under the actual filtration at every trial.
An unconditional mean gap does not supply it.

\subsection{Sampled generators and the information supplied by a clock}
\label{op:time}

A finite continuous-time Markov generator \(Q\) has nonnegative
off-diagonal entries and zero row sums, and produces transition
matrices \(e^{tQ}\) for \(t\ge0\).
Recovering \(Q\) from sampled transitions is the classical embedding
and matrix-exponential inverse problem
\cite{op:baake,op:yue}.  It presupposes a state basis and time coordinate.

\begin{theorem}[Two positive generators with all integer samples equal]
\label{op:generator-alias}
Let
\[
 C=\begin{pmatrix}0&1&0\\0&0&1\\1&0&0\end{pmatrix},\quad
 \Pi=\tfrac13\mathbf1\mathbf1^T,\quad R=I-\Pi,\quad A=C-C^2,
\]
and define
\[
 Q_0=-12R,\qquad Q_1=-12R+\frac{2\pi}{\sqrt3}A .
\]
Both are irreducible generators with uniform stationary distribution
and strictly positive off-diagonal rates.
They obey \(e^{nQ_0}=e^{nQ_1}\) for every integer \(n\ge0\), but are
neither similar nor related by a positive scalar time rescaling,
even when a similarity is also allowed.
One is reversible and the other has nonzero stationary circulation.
\end{theorem}
\begin{proof}
The off-diagonal rates are \(4\) for \(Q_0\) and
\(4\pm2\pi/\sqrt3>0\) for \(Q_1\).
The inequality follows, for example, from
\(\pi<22/7<2\sqrt3\).
Rows and columns sum to zero, so the uniform distribution is
stationary; strict positive rates imply irreducibility and uniqueness.
Direct algebra gives
\[
 A\Pi=\Pi A=0,\quad A^2=-3R,\quad AR=RA=A.
\]
Summing the exponential series on the invariant subspaces yields
\[
 e^{tQ_0}=\Pi+e^{-12t}R,\qquad
 e^{tQ_1}=\Pi+e^{-12t}
 \left[\cos(2\pi t)R+\frac{\sin(2\pi t)}{\sqrt3}A\right].
\]
At integer \(t\) the kernels agree.
The spectra are \(0,-12,-12\) and \(0,-12\pm2\pi i\), respectively.
Similarity preserves spectra, and a positive scalar cannot turn real
nonzero eigenvalues into nonreal ones.
The symmetric generator \(Q_0\) is reversible for the uniform law.
For \(Q_1\), the oriented stationary flux difference along a cycle edge
is \(\frac13[(4+2\pi/\sqrt3)-(4-2\pi/\sqrt3)]
=4\pi/(3\sqrt3)>0\).
\end{proof}

\begin{proposition}[Exactly the integer-grid intervention menu]
\label{op:grid-policy}
Use either generator in Theorem~\ref{op:generator-alias}, a common
initial law, and the same instantaneous labelled boundary controls
and instruments.  Suppose every control and every readout, including
silent state maps, occurs on the integer time grid, the between-boundary
generator stays fixed, and the policy depends only on the common
observed history and private randomness.
Then the complete finite controlled observation-history laws agree.
\end{proposition}
\begin{proof}
A fixed outcome-history probability is a row product of the initial
law, transition matrices for nonnegative integer durations, and the
corresponding boundary instrument matrices, followed by \(\mathbf1\).
Every factor is common to the two descriptions, so the probability
is common.  For an adaptive randomized policy, condition successively
on each history; its next action, scheduling choice, and stop decision
have the same conditional law.  Induction gives the result.
\end{proof}

Integer-grid readouts alone are insufficient if silent controls are
allowed at off-grid times.  Persistent drives that change the generator,
continuous monitoring, and observations of jump or waiting times also
lie outside this menu and may supply distinguishing information.

An anchored preparation of state one and a readout distinguishing
states two and three at time \(t=1/4\) give the contrast
\[
 (e^{Q_1/4})_{12}-(e^{Q_1/4})_{13}
 =\frac{2e^{-3}}{\sqrt3}
 =0.0574891546486971\ldots ,
\]
whereas the contrast is zero for \(Q_0\).
An unconditioned stationary uniform preparation has identical uniform
marginals under both models and does not provide this witness.
The contrast assumes a calibrated quarter-time setting; it is not
claimed optimal and carries no trial-count guarantee.

\begin{theorem}[Exact identification from two incommensurate times]
\label{op:two-times}
Let \(Q,B\) be finite real matrices represented in the same state basis.
Let \(t_1,t_2>0\) have irrational ratio.
If
\[
 e^{t_1Q}=e^{t_1B},\qquad e^{t_2Q}=e^{t_2B},
\]
then \(Q=B\).
\end{theorem}
\begin{proof}
Denote the two common exponential matrices by \(M_1,M_2\).
They are invertible and commute, since they are both functions of \(Q\).
For every pair of integers \(m,n\),
\[
 e^{(mt_1+nt_2)Q}=M_1^mM_2^n=e^{(mt_1+nt_2)B}.
\]
The subgroup \(\{mt_1+nt_2:m,n\in\mathbb Z\}\) is dense in
\(\mathbb R\).  One proof is that irrational rotation modulo \(t_2\)
has arbitrarily small positive return differences by the pigeonhole
principle; integer multiples of those differences approximate any
real number, and all such multiples remain in the subgroup.
Continuity gives equality of the exponentials for every real time.
Differentiating at zero yields \(Q=B\).
\end{proof}

The proof does not assume \(Q\) and \(B\) commute with one another.
Negative powers are algebraic inverses, not proposed physical operations
at negative time.  The result requires full exact transition matrices,
a common basis, homogeneous dynamics, and exact times.
Finite-precision measurements cannot establish exact irrationality of
a time ratio, and this uniqueness theorem is not a stability theorem.

\begin{theorem}[No uniform inverse stability over unbounded rates]
\label{op:rate-instability}
Fix a nonzero real \(c\) and any finite collection of strictly positive
observation times.  With the matrices \(R,A\) above, put
\[
 Q_a=-aR,\qquad B_a=-aR+cA,\qquad a>3|c|.
\]
These are irreducible Markov generators.  Their difference is the fixed
nonzero matrix \(cA\), while
\(\|e^{tQ_a}-e^{tB_a}\|\to0\) as \(a\to\infty\) at every specified
time, for any matrix norm.
Thus no uniform inverse error modulus holds over the unrestricted-rate
class.  In particular, even when the menu is injective, as for two
incommensurate times, its inverse is not uniformly continuous.
\end{theorem}
\begin{proof}
Off-diagonal entries are \(a/3\) or \(a/3\pm c\), all positive.
The same invariant-subspace calculation gives
\[
 e^{tB_a}-e^{tQ_a}
 =e^{-at}\left[
   (\cos(\sqrt3ct)-1)R+
   \frac{\sin(\sqrt3ct)}{\sqrt3}A\right].
\]
For fixed positive \(t\), the bracket is bounded independently of \(a\)
and the prefactor tends to zero.  A finite maximum over observation
times tends to zero as well, while
\(\|B_a-Q_a\|=|c|\|A\|>0\).
\end{proof}

On a specified compact generator class, the exact two-time map of
Theorem~\ref{op:two-times} is continuous and injective, so its inverse
on its image is uniformly continuous by compactness.
That qualitative exception provides no quantitative modulus, and no
joint finite-data result for uncertain timing and uncertain transition
matrices is obtained here.

\subsection{Scope of the mathematical conclusions}
\label{op:scope}

The positive results identify precisely what their declared experiments
determine: autonomous coarse dynamics under a fiber condition; a
classical predictive-state lower bound from a joint-law matrix; valid
finite-sample or sequential tests under explicit sampling assumptions;
finite-state reconstruction under exact anchored controls; and a
generator under full exact observations at two incommensurate times.
The counterexamples identify specific missing information: short-block
agreement, a final channel without fragment access, a single silent
control, open calibration uncertainty, or sampling on a coarse time
grid can leave distinct descriptions observationally equivalent.

All these conclusions are conditional on conventional mathematical
primitives and the stated observation maps.  They do not select a
physical tensor partition, demonstrate that a detector is independent,
derive a probability rule from a prephysical proposal, or provide an
emergent physical clock.  A physical correspondence must identify which
preparations, interventions, outcomes, reset procedures, and timings
realize the premises and must confront the appropriate competing
observation laws.

\clearpage
\begin{thebibliography}{99}
\bibitem{bodirsky}
M. Bodirsky. Cores of countably categorical structures.
\emph{Logical Methods in Computer Science} \textbf{3}(1:2), 1--16 (2007).
\url{https://arxiv.org/abs/cs/0612069v2}.

\bibitem{stacksinitial}
The Stacks Project Authors. Initial and final objects, Definition 4.12.1.
\emph{The Stacks Project}, Tag 002B.
\url{https://stacks.math.columbia.edu/tag/002B}.

\bibitem{abramsky}
S. Abramsky. No-cloning in categorical quantum mechanics.
In \emph{Semantic Techniques in Quantum Computation}, pp. 1--28,
Cambridge University Press (2010).
Author manuscript: \url{https://arxiv.org/abs/0910.2401}.

\bibitem{sorkin}
R. D. Sorkin. Quantum mechanics as quantum measure theory.
\emph{Modern Physics Letters A} \textbf{9}, 3119--3127 (1994).
\url{https://arxiv.org/abs/gr-qc/9401003}.

\bibitem{boes}
P. Boes and M. Navascu\'es. Composing decoherence functionals.
\emph{Physical Review A} \textbf{95}, 022114 (2017).
\url{https://doi.org/10.1103/PhysRevA.95.022114}.

\bibitem{shannon}
C. E. Shannon. A mathematical theory of communication.
\emph{Bell System Technical Journal} \textbf{27}, 379--423 and 623--656 (1948).
\url{https://doi.org/10.1002/j.1538-7305.1948.tb01338.x};
\url{https://doi.org/10.1002/j.1538-7305.1948.tb00917.x}.

\bibitem{pitman}
J. Pitman. \emph{Combinatorial Stochastic Processes}.
Lecture Notes in Mathematics \textbf{1875}, Springer (2006).
\url{https://doi.org/10.1007/b11601500}.

\bibitem{clemencon}
S. Cl\'emen\c{c}on, I. Colin and A. Bellet.
Scaling-up empirical risk minimization: optimization of incomplete $U$-statistics.
\emph{Journal of Machine Learning Research} \textbf{17}(76), 1--36 (2016).
\url{https://jmlr.org/papers/v17/15-012.html}.

\bibitem{hoeffding}
W. Hoeffding. Probability inequalities for sums of bounded random variables.
\emph{Journal of the American Statistical Association} \textbf{58}(301), 13--30 (1963).
\url{https://doi.org/10.1080/01621459.1963.10500830}.

\bibitem{horodecki}
R. Horodecki, P. Horodecki and M. Horodecki.
Violating Bell inequality by mixed spin-$1/2$ states: necessary and sufficient condition.
\emph{Physics Letters A} \textbf{200}, 340--344 (1995).
\url{https://doi.org/10.1016/0375-9601(95)00214-N}.

\bibitem{chenghall}
S. Cheng and M. J. W. Hall.
Anisotropic invariance and the distribution of quantum correlations.
\emph{Physical Review Letters} \textbf{118}, 010401 (2017).
\url{https://doi.org/10.1103/PhysRevLett.118.010401}.
Author manuscript: \url{https://arxiv.org/abs/1610.09302v3}.

\bibitem{caves}
C. M. Caves, C. A. Fuchs and R. Schack.
Unknown quantum states: the quantum de Finetti representation.
\emph{Journal of Mathematical Physics} \textbf{43}, 4537--4559 (2002).
\url{https://doi.org/10.1063/1.1494475}.

\bibitem{gaussianNotes} Y. Polyanskiy and Y. Wu, \emph{Information Theory}, MIT 6.441 course notes (2016), Chapters 17 and 19, on Gaussian channels, input constraints and continuous-time models. \url{https://ocw.mit.edu/courses/6-441-information-theory-spring-2016/}.
\bibitem{justusGaussian} K. Justus, \emph{Abstract Channels and Gaussian Capacity}, manuscript in the accompanying Mathematical Supplement (10 November 2026), especially the main Gaussian theorem and Proposition 6. AI assistance is disclosed in the study.
\bibitem{justusRecovery} K. Justus, \emph{Stochastic Recovery from Information Loss and Equality of Channel Capacity}, manuscript in the accompanying Mathematical Supplement (10 November 2026). AI assistance is disclosed in the study.
\bibitem{justusCalibration} K. Justus, \emph{Calibration Uncertainty and Finite-Sample Identification}, manuscript in the accompanying Mathematical Supplement (10 November 2026). AI assistance is disclosed in the study.

\bibitem{rec:operational-study} K. Justus, \emph{Operational Identification of Records, Controls and Continuous-Time Dynamics}, unpublished manuscript (10 November 2026), Proposition 11. Retained as a related source edition; AI assistance is disclosed in the study.

\bibitem{justusJoint} K. Justus, \emph{Calibration, Identified Sets and Uniform Limits}, mathematical study in the accompanying supplement (10 November 2026). The supplied companion edition is titled \emph{Mathematical Supplement: Calibration, identified sets and uniform limits}. Same-program source with AI assistance.
\bibitem{justusWeakGain} K. Justus, \emph{Statistical Precision under Weak Readout Gain}, mathematical study in the accompanying supplement (10 November 2026). Same-program source with AI assistance.
\bibitem{justusAssessment} K. Justus, \emph{Scientific Assessment and Further Results}, supplied Collections companion (10 November 2026), Sections 3.2--3.3. It examines the separately preserved 81-page and 420-page source edition; its page references and counts concern that edition.
\bibitem{beckerEntropy} S. Becker, N. Datta, M. G. Jabbour and M. E. Shirokov, \emph{Optimal continuity bound for the von Neumann entropy under energy constraints}, arXiv:2410.02686v2 (2025), Section 2. \url{https://arxiv.org/abs/2410.02686}. The all-positive-temperature Gibbs hypothesis is stronger than convergence of one partition sum.

\bibitem{op:wolfer}
G. Wolfer and S. Watanabe.
\emph{Geometric Aspects of Data-Processing of Markov Chains}.
Author manuscript, arXiv:2203.04575v3 (2023).
\url{https://arxiv.org/abs/2203.04575v3}.

\bibitem{op:ryabko}
B. Ryabko.
Low-Entropy Stochastic Processes for Generating \(k\)-Distributed and Normal
Sequences, and the Relationship of These Processes with Random Number Generators.
\emph{Mathematics} \textbf{7}(9), 838 (2019).
\url{https://doi.org/10.3390/math7090838}.

\bibitem{op:neyman}
J. Neyman and E. S. Pearson.
On the Problem of the Most Efficient Tests of Statistical Hypotheses.
\emph{Philosophical Transactions of the Royal Society of London A}
\textbf{231}, 289--337 (1933).
\url{https://doi.org/10.1098/rsta.1933.0009}.

\bibitem{op:ohta}
Y. Ohta.
\emph{On the Realization of Hidden Markov Models and Tensor Decomposition}.
Author manuscript, arXiv:2008.11487 (2020).
\url{https://arxiv.org/abs/2008.11487}.

\bibitem{op:vanluyten}
B. Vanluyten, J. C. Willems, and B. De Moor.
Equivalence of state representations for hidden Markov models.
\emph{Systems \& Control Letters} \textbf{57}, 410--419 (2008).
\url{https://doi.org/10.1016/j.sysconle.2007.10.004}.

\bibitem{op:eckart}
C. Eckart and G. Young.
The Approximation of One Matrix by Another of Lower Rank.
\emph{Psychometrika} \textbf{1}, 211--218 (1936).
\url{https://doi.org/10.1007/BF02288367}.

\bibitem{op:mcdiarmid}
C. McDiarmid.
On the method of bounded differences.
In \emph{Surveys in Combinatorics, 1989}, pp. 148--188.
Cambridge University Press (1989).
\url{https://doi.org/10.1017/CBO9781107359949.008}.

\bibitem{op:howard}
S. R. Howard, A. Ramdas, J. McAuliffe, and J. Sekhon.
Time-uniform Chernoff bounds via nonnegative supermartingales.
\emph{Probability Surveys} \textbf{17}, 257--317 (2020).
\url{https://doi.org/10.1214/18-PS321}.

\bibitem{op:larsson}
M. Larsson, A. Ramdas, and J. Ruf.
\emph{The numeraire e-variable and reverse information projection}.
Author manuscript, arXiv:2402.18810v2 (2024).
\url{https://arxiv.org/abs/2402.18810v2}.

\bibitem{op:crow}
D. Crow and R. Joynt.
Classical simulation of quantum dephasing and depolarizing noise.
\emph{Physical Review A} \textbf{89}, 042123 (2014).
\url{https://doi.org/10.1103/PhysRevA.89.042123}.

\bibitem{op:baake}
M. Baake and J. Sumner.
Notes on Markov embedding.
\emph{Linear Algebra and its Applications} \textbf{594}, 262--299 (2020).
\url{https://arxiv.org/abs/1903.08736v2}.

\bibitem{op:yue}
Z. Yue, J. Thunberg, and J. Goncalves.
\emph{Inverse Problems for Matrix Exponential in System Identification:
System Aliasing}.
Author manuscript, arXiv:1605.06973 (2016).
\url{https://arxiv.org/abs/1605.06973}.

\bibitem{fin:signless}
D. Cvetkovi\'c, P. Rowlinson, and S. K. Simi\'c.
Signless Laplacians of finite graphs.
\emph{Linear Algebra and its Applications} \textbf{423}, 155--171 (2007).
\url{https://doi.org/10.1016/j.laa.2007.01.009}.

\bibitem{fin:ising}
W. Vinci, K. Markstr\"om, S. Boixo, A. Roy, F. M. Spedalieri,
P. A. Warburton, and S. Severini.
Hearing the Shape of the Ising Model with a Programmable
Superconducting-Flux Annealer.
\emph{Scientific Reports} \textbf{4}, 5703 (2014).
\url{https://doi.org/10.1038/srep05703}.

\bibitem{fin:extendibility}
T. Konstantopoulos and L. Yuan.
On the extendibility of finitely exchangeable probability measures.
\emph{Transactions of the American Mathematical Society}
\textbf{371}, 7067--7092 (2019).
\url{https://doi.org/10.1090/tran/7701}.

\end{thebibliography}
\end{document}
